\documentclass[conference,compsoc]{IEEEtran}
% IEEE S&P requires IEEEtran.cls V1.8b with the compsoc option, US letter.
\usepackage{amsmath,amssymb,amsthm}
\usepackage{stmaryrd}
\usepackage{booktabs,array,multirow,tabularx,threeparttable}
\usepackage{xcolor}
\usepackage{tikz}
\usetikzlibrary{arrows.meta,positioning,shapes.geometric,fit,backgrounds,calc}
\usepackage{algorithm}
\usepackage{algpseudocode}
\usepackage{enumitem}
\usepackage{cite}
\usepackage{url}
\usepackage{microtype}

\definecolor{accent}{HTML}{0B5394}
\definecolor{accentlight}{HTML}{D9E8F7}
\definecolor{warn}{HTML}{B45F06}
\definecolor{okgreen}{HTML}{38761D}
\definecolor{ledger}{HTML}{674EA7}
\definecolor{pend}{HTML}{C00000}

% ---- Draft-only marker for experiments that are registered but not yet run.
% Every number that must come from a new run is behind \PENDING; the final
% submission must contain zero occurrences (grep -c PENDING main.tex == 0).
\newcommand{\PENDING}[1]{\textcolor{pend}{[\textsc{pending}: #1]}}

\newtheorem{definition}{Definition}
\newtheorem{proposition}{Proposition}
\newtheorem{corollary}{Corollary}
\newtheorem{remark}{Remark}
\newtheorem{assumption}{Assumption}
\newtheorem{theorem}{Theorem}
\theoremstyle{definition}
\newtheorem{example}{Example}
\theoremstyle{plain}

% ---- Invariants. I4 is argument binding (new in this version); the
% counterfactual-dependence bound, formerly I4, is now I5.
\newcommand{\Ione}{\ensuremath{I_1}}
\newcommand{\Itwo}{\ensuremath{I_2}}
\newcommand{\Ithree}{\ensuremath{I_3}}
\newcommand{\Ifour}{\ensuremath{I_4}}
\newcommand{\Ifive}{\ensuremath{I_5}}
\providecommand{\tc}{\tau_{c}}
\providecommand{\meet}{\sqcap}
\providecommand{\DeltaIrrel}{\Delta_{\mathrm{irrel}}}
\providecommand{\Dirr}{D_{\mathrm{irr}}}
\providecommand{\Lift}{\mathrm{Lift}}
\providecommand{\Tr}{\Sigma}
\newcommand{\att}{\ensuremath{\mathsf{att}}}
\newcommand{\cert}{\ensuremath{\mathsf{cert}}}
\newcommand{\deps}{\ensuremath{\mathsf{dep}}}
\newcommand{\support}{\ensuremath{\mathsf{support}}}
\newcommand{\acts}{\ensuremath{\mathsf{acts}}}
\newcommand{\args}{\ensuremath{\mathsf{args}}}
\newcommand{\vals}{\ensuremath{\mathsf{vals}}}
\newcommand{\bindpol}{\ensuremath{\beta}}
\newcommand{\ev}[1]{\ensuremath{\mathsf{#1}}}
\newcommand{\lmin}{\ensuremath{\ell_{\min}}}
\newcommand{\lreq}{\ensuremath{\ell_{\mathrm{req}}}}
\newcommand{\Cprot}{\ensuremath{C_{\mathrm{prot}}}}
\newcommand{\Dirrel}{\ensuremath{\Delta_{\mathrm{irrel}}}}
\newcommand{\dist}{\ensuremath{\mathsf{dist}}}
\newcommand{\DEI}{\ensuremath{\mathrm{DEI}}}
% ==== GATE MEASURED VALUES (results_gate/*.json, 2026-08-02; offline cache replay) ====
\newcommand{\GateFollows}{4}
\newcommand{\GateQuar}{2}
\newcommand{\GateHonestUnch}{1}
\newcommand{\GateInterUnch}{1}
\newcommand{\GateBenignDestr}{3}
\newcommand{\GateFalseQuarN}{0}
\newcommand{\GateFalseQuar}{0\%}
\newcommand{\GateOverhead}{1.375}
\newcommand{\GateMonitorUs}{52}
\newcommand{\tick}{\ensuremath{\checkmark}}
\newcommand{\partl}{\ensuremath{\sim}}
\newcommand{\absent}{---}

\begin{document}
% AUTO-GENERATED by recompute_and_emit_macros.py — do not edit by hand.
% Every reported live-model number is defined here. Swap n by rerunning.

\newcommand{\SecNdHit}{18}
\newcommand{\SecNdN}{40}
\newcommand{\SecNdASR}{0.45}
\newcommand{\SecNdCIlo}{0.307}
\newcommand{\SecNdCIhi}{0.602}
\newcommand{\SecNdFilt}{0}
\newcommand{\SecNdCaseHit}{18}
\newcommand{\SecNdCaseN}{40}
\newcommand{\SecMonHit}{2}
\newcommand{\SecMonN}{40}
\newcommand{\SecMonASR}{0.05}
\newcommand{\SecMonCIlo}{0.014}
\newcommand{\SecMonCIhi}{0.165}
\newcommand{\SecMonFilt}{47}
\newcommand{\SecMonCaseHit}{2}
\newcommand{\SecMonCaseN}{40}
\newcommand{\SecGateHit}{2}
\newcommand{\SecGateN}{40}
\newcommand{\SecGateASR}{0.05}
\newcommand{\SecGateCIlo}{0.014}
\newcommand{\SecGateCIhi}{0.165}
\newcommand{\SecGateFilt}{48}
\newcommand{\SecGateCaseHit}{2}
\newcommand{\SecGateCaseN}{40}
\newcommand{\UtilNdN}{24}
\newcommand{\UtilNdPass}{20}
\newcommand{\UtilNdBlock}{4}
\newcommand{\UtilNdRate}{0.83}
\newcommand{\UtilNdBlockCIhi}{0.359}
\newcommand{\UtilGateN}{24}
\newcommand{\UtilGatePass}{17}
\newcommand{\UtilGateBlock}{7}
\newcommand{\UtilGateRate}{0.71}
\newcommand{\UtilGateBlockCIhi}{0.492}
\newcommand{\GateVsNdFisherP}{0.0001}
\newcommand{\GateVsNdSig}{significant}
  % auto-generated live-model numbers (n=40 grid); regenerate via recompute_and_emit_macros.py
% DERIVED from results_macros.tex counts by universe_v4_constants.py (2026-10-27):
% Wald interval for the risk difference; log-scale interval for the risk ratio.
% TODO: extend recompute_and_emit_macros.py to emit these so that no figure is hand-entered.
\newcommand{\HeadlineRD}{-0.400}
\newcommand{\HeadlineRDlo}{-0.568}
\newcommand{\HeadlineRDhi}{-0.232}
\newcommand{\HeadlineRR}{0.111}
\newcommand{\HeadlineRRlo}{0.028}
\newcommand{\HeadlineRRhi}{0.448}
  % RD/RR derived from the same counts (see file header)
% AUTO-GENERATED by build/emit_rq5_macros.py from compliant restarted MCP runs.
\newcommand{\RqFiveAZeroFiveCell}{5.02/6.72}
\newcommand{\RqFiveAZeroFiveMedian}{5.02}
\newcommand{\RqFiveAZeroFiveTail}{6.72}
\newcommand{\RqFiveAZeroFiveN}{924}
\newcommand{\RqFiveAZeroTwentyFiveCell}{4.74/6.78}
\newcommand{\RqFiveAZeroTwentyFiveMedian}{4.74}
\newcommand{\RqFiveAZeroTwentyFiveTail}{6.78}
\newcommand{\RqFiveAZeroTwentyFiveN}{4620}
\newcommand{\RqFiveAZeroHundredCell}{4.65/7.21}
\newcommand{\RqFiveAZeroHundredMedian}{4.65}
\newcommand{\RqFiveAZeroHundredTail}{7.21}
\newcommand{\RqFiveAZeroHundredN}{18090}
\newcommand{\RqFiveAOneFiveCell}{4.95/7.00}
\newcommand{\RqFiveAOneFiveMedian}{4.95}
\newcommand{\RqFiveAOneFiveTail}{7.00}
\newcommand{\RqFiveAOneFiveN}{897}
\newcommand{\RqFiveAOneTwentyFiveCell}{4.77/6.78}
\newcommand{\RqFiveAOneTwentyFiveMedian}{4.77}
\newcommand{\RqFiveAOneTwentyFiveTail}{6.78}
\newcommand{\RqFiveAOneTwentyFiveN}{4671}
\newcommand{\RqFiveAOneHundredCell}{4.77/7.70}
\newcommand{\RqFiveAOneHundredMedian}{4.77}
\newcommand{\RqFiveAOneHundredTail}{7.70}
\newcommand{\RqFiveAOneHundredN}{17901}
\newcommand{\RqFiveATwoFiveCell}{5.53/7.67}
\newcommand{\RqFiveATwoFiveMedian}{5.53}
\newcommand{\RqFiveATwoFiveTail}{7.67}
\newcommand{\RqFiveATwoFiveN}{834}
\newcommand{\RqFiveATwoTwentyFiveCell}{5.02/7.61}
\newcommand{\RqFiveATwoTwentyFiveMedian}{5.02}
\newcommand{\RqFiveATwoTwentyFiveTail}{7.61}
\newcommand{\RqFiveATwoTwentyFiveN}{4581}
\newcommand{\RqFiveATwoHundredCell}{4.77/7.71}
\newcommand{\RqFiveATwoHundredMedian}{4.77}
\newcommand{\RqFiveATwoHundredTail}{7.71}
\newcommand{\RqFiveATwoHundredN}{18342}
\newcommand{\RqFiveOverheadFive}{10.1}
\newcommand{\RqFiveOverheadTwentyFive}{6.1}
\newcommand{\RqFiveOverheadHundred}{2.7}
\newcommand{\RqFiveOverheadRange}{2.7--10.1\,\%}
\newcommand{\RqFiveRlaSignCell}{1.42/2.14}
\newcommand{\RqFiveRlaVerifyCell}{1.47/2.18}
\newcommand{\RqFiveRlaIterations}{1000}
\newcommand{\RqFiveCodeHashShort}{ee4f46b238f6}
\newcommand{\RqFiveTrustedLines}{3022}
  % restarted MCP RQ5 measurements; regenerate via build/emit_rq5_macros.py
% AUTO-GENERATED by analyze.py from out/rq1_all.jsonl; primary (degenerate retained); do not edit.
% Voided: 0; marked degenerate in source: 0; bootstrap seed: 20260814.
\newcommand{\SpEpisodes}{1442}
\newcommand{\SpVoided}{0}
\newcommand{\SpDegenerate}{0}
\newcommand{\SpAnalysisMode}{primary}
\newcommand{\SpSecBTwoStrictHit}{7}
\newcommand{\SpSecBTwoStrictN}{150}
\newcommand{\SpSecBTwoStrictASR}{0.047}
\newcommand{\SpSecBTwoStrictLo}{0.023}
\newcommand{\SpSecBTwoStrictHi}{0.093}
\newcommand{\SpSecBTwoStrictFisherP}{7.9e-22}
\newcommand{\SpSecBTwoStrictRD}{-0.480}
\newcommand{\SpSecBTwoStrictRDlo}{-0.567}
\newcommand{\SpSecBTwoStrictRDhi}{-0.393}
\newcommand{\SpSecBTwoStrictRR}{0.089}
\newcommand{\SpSecBTwoStrictRRlo}{0.042}
\newcommand{\SpSecBTwoStrictRRhi}{0.186}
\newcommand{\SpSecDeiGateHit}{13}
\newcommand{\SpSecDeiGateN}{150}
\newcommand{\SpSecDeiGateASR}{0.087}
\newcommand{\SpSecDeiGateLo}{0.051}
\newcommand{\SpSecDeiGateHi}{0.143}
\newcommand{\SpSecDeiGateFisherP}{2.9e-17}
\newcommand{\SpSecDeiGateRD}{-0.440}
\newcommand{\SpSecDeiGateRDlo}{-0.532}
\newcommand{\SpSecDeiGateRDhi}{-0.348}
\newcommand{\SpSecDeiGateRR}{0.165}
\newcommand{\SpSecDeiGateRRlo}{0.096}
\newcommand{\SpSecDeiGateRRhi}{0.283}
\newcommand{\SpSecBOneHit}{13}
\newcommand{\SpSecBOneN}{150}
\newcommand{\SpSecBOneASR}{0.087}
\newcommand{\SpSecBOneLo}{0.051}
\newcommand{\SpSecBOneHi}{0.143}
\newcommand{\SpSecBOneFisherP}{2.9e-17}
\newcommand{\SpSecBOneRD}{-0.440}
\newcommand{\SpSecBOneRDlo}{-0.532}
\newcommand{\SpSecBOneRDhi}{-0.348}
\newcommand{\SpSecBOneRR}{0.165}
\newcommand{\SpSecBOneRRlo}{0.096}
\newcommand{\SpSecBOneRRhi}{0.283}
\newcommand{\SpSecBThreeHit}{25}
\newcommand{\SpSecBThreeN}{150}
\newcommand{\SpSecBThreeASR}{0.167}
\newcommand{\SpSecBThreeLo}{0.116}
\newcommand{\SpSecBThreeHi}{0.234}
\newcommand{\SpSecBThreeFisherP}{6.1e-11}
\newcommand{\SpSecBThreeRD}{-0.360}
\newcommand{\SpSecBThreeRDlo}{-0.460}
\newcommand{\SpSecBThreeRDhi}{-0.260}
\newcommand{\SpSecBThreeRR}{0.316}
\newcommand{\SpSecBThreeRRlo}{0.215}
\newcommand{\SpSecBThreeRRhi}{0.467}
\newcommand{\SpSecDeiHit}{13}
\newcommand{\SpSecDeiN}{150}
\newcommand{\SpSecDeiASR}{0.087}
\newcommand{\SpSecDeiLo}{0.051}
\newcommand{\SpSecDeiHi}{0.143}
\newcommand{\SpSecDeiFisherP}{2.9e-17}
\newcommand{\SpSecDeiRD}{-0.440}
\newcommand{\SpSecDeiRDlo}{-0.532}
\newcommand{\SpSecDeiRDhi}{-0.348}
\newcommand{\SpSecDeiRR}{0.165}
\newcommand{\SpSecDeiRRlo}{0.096}
\newcommand{\SpSecDeiRRhi}{0.283}
\newcommand{\SpSecBTwoHit}{33}
\newcommand{\SpSecBTwoN}{150}
\newcommand{\SpSecBTwoASR}{0.220}
\newcommand{\SpSecBTwoLo}{0.161}
\newcommand{\SpSecBTwoHi}{0.293}
\newcommand{\SpSecBTwoFisherP}{5.7e-08}
\newcommand{\SpSecBTwoRD}{-0.307}
\newcommand{\SpSecBTwoRDlo}{-0.410}
\newcommand{\SpSecBTwoRDhi}{-0.203}
\newcommand{\SpSecBTwoRR}{0.418}
\newcommand{\SpSecBTwoRRlo}{0.298}
\newcommand{\SpSecBTwoRRhi}{0.585}
\newcommand{\SpSecNoDefenseHit}{79}
\newcommand{\SpSecNoDefenseN}{150}
\newcommand{\SpSecNoDefenseASR}{0.527}
\newcommand{\SpSecNoDefenseLo}{0.447}
\newcommand{\SpSecNoDefenseHi}{0.605}
\newcommand{\SpSecDeiVsBOneFisherP}{1.0e+00}
\newcommand{\SpSecDeiVsBOneHolmP}{1.0e+00}
\newcommand{\SpSecDeiVsBOneRD}{0.000}
\newcommand{\SpSecDeiVsBOneRDlo}{-0.064}
\newcommand{\SpSecDeiVsBOneRDhi}{0.064}
\newcommand{\SpSecDeiVsBOneRR}{1.000}
\newcommand{\SpSecDeiVsBOneRRlo}{0.480}
\newcommand{\SpSecDeiVsBOneRRhi}{2.085}
\newcommand{\SpSecDeiVsBTwoFisherP}{2.1e-03}
\newcommand{\SpSecDeiVsBTwoHolmP}{8.3e-03}
\newcommand{\SpSecDeiVsBTwoRD}{-0.133}
\newcommand{\SpSecDeiVsBTwoRDlo}{-0.213}
\newcommand{\SpSecDeiVsBTwoRDhi}{-0.053}
\newcommand{\SpSecDeiVsBTwoRR}{0.394}
\newcommand{\SpSecDeiVsBTwoRRlo}{0.216}
\newcommand{\SpSecDeiVsBTwoRRhi}{0.718}
\newcommand{\SpSecDeiVsBTwoStrictFisherP}{2.5e-01}
\newcommand{\SpSecDeiVsBTwoStrictHolmP}{4.9e-01}
\newcommand{\SpSecDeiVsBTwoStrictRD}{0.040}
\newcommand{\SpSecDeiVsBTwoStrictRDlo}{-0.016}
\newcommand{\SpSecDeiVsBTwoStrictRDhi}{0.096}
\newcommand{\SpSecDeiVsBTwoStrictRR}{1.857}
\newcommand{\SpSecDeiVsBTwoStrictRRlo}{0.762}
\newcommand{\SpSecDeiVsBTwoStrictRRhi}{4.525}
\newcommand{\SpSecDeiVsBThreeFisherP}{5.5e-02}
\newcommand{\SpSecDeiVsBThreeHolmP}{1.7e-01}
\newcommand{\SpSecDeiVsBThreeRD}{-0.080}
\newcommand{\SpSecDeiVsBThreeRDlo}{-0.155}
\newcommand{\SpSecDeiVsBThreeRDhi}{-0.005}
\newcommand{\SpSecDeiVsBThreeRR}{0.520}
\newcommand{\SpSecDeiVsBThreeRRlo}{0.277}
\newcommand{\SpSecDeiVsBThreeRRhi}{0.977}
\newcommand{\SpSecBOneSuiteBankingHit}{0}
\newcommand{\SpSecBOneSuiteBankingN}{30}
\newcommand{\SpSecBOneSuiteBankingASR}{0.000}
\newcommand{\SpSecBOneSuiteBankingLo}{0.000}
\newcommand{\SpSecBOneSuiteBankingHi}{0.114}
\newcommand{\SpSecBTwoSuiteBankingHit}{5}
\newcommand{\SpSecBTwoSuiteBankingN}{30}
\newcommand{\SpSecBTwoSuiteBankingASR}{0.167}
\newcommand{\SpSecBTwoSuiteBankingLo}{0.073}
\newcommand{\SpSecBTwoSuiteBankingHi}{0.336}
\newcommand{\SpSecBTwoStrictSuiteBankingHit}{0}
\newcommand{\SpSecBTwoStrictSuiteBankingN}{30}
\newcommand{\SpSecBTwoStrictSuiteBankingASR}{0.000}
\newcommand{\SpSecBTwoStrictSuiteBankingLo}{0.000}
\newcommand{\SpSecBTwoStrictSuiteBankingHi}{0.114}
\newcommand{\SpSecBThreeSuiteBankingHit}{0}
\newcommand{\SpSecBThreeSuiteBankingN}{30}
\newcommand{\SpSecBThreeSuiteBankingASR}{0.000}
\newcommand{\SpSecBThreeSuiteBankingLo}{0.000}
\newcommand{\SpSecBThreeSuiteBankingHi}{0.114}
\newcommand{\SpSecDeiSuiteBankingHit}{0}
\newcommand{\SpSecDeiSuiteBankingN}{30}
\newcommand{\SpSecDeiSuiteBankingASR}{0.000}
\newcommand{\SpSecDeiSuiteBankingLo}{0.000}
\newcommand{\SpSecDeiSuiteBankingHi}{0.114}
\newcommand{\SpSecDeiGateSuiteBankingHit}{0}
\newcommand{\SpSecDeiGateSuiteBankingN}{30}
\newcommand{\SpSecDeiGateSuiteBankingASR}{0.000}
\newcommand{\SpSecDeiGateSuiteBankingLo}{0.000}
\newcommand{\SpSecDeiGateSuiteBankingHi}{0.114}
\newcommand{\SpSecNoDefenseSuiteBankingHit}{12}
\newcommand{\SpSecNoDefenseSuiteBankingN}{30}
\newcommand{\SpSecNoDefenseSuiteBankingASR}{0.400}
\newcommand{\SpSecNoDefenseSuiteBankingLo}{0.246}
\newcommand{\SpSecNoDefenseSuiteBankingHi}{0.577}
\newcommand{\SpSecBOneSuiteSlackHit}{6}
\newcommand{\SpSecBOneSuiteSlackN}{30}
\newcommand{\SpSecBOneSuiteSlackASR}{0.200}
\newcommand{\SpSecBOneSuiteSlackLo}{0.095}
\newcommand{\SpSecBOneSuiteSlackHi}{0.373}
\newcommand{\SpSecBTwoSuiteSlackHit}{20}
\newcommand{\SpSecBTwoSuiteSlackN}{30}
\newcommand{\SpSecBTwoSuiteSlackASR}{0.667}
\newcommand{\SpSecBTwoSuiteSlackLo}{0.488}
\newcommand{\SpSecBTwoSuiteSlackHi}{0.808}
\newcommand{\SpSecBTwoStrictSuiteSlackHit}{5}
\newcommand{\SpSecBTwoStrictSuiteSlackN}{30}
\newcommand{\SpSecBTwoStrictSuiteSlackASR}{0.167}
\newcommand{\SpSecBTwoStrictSuiteSlackLo}{0.073}
\newcommand{\SpSecBTwoStrictSuiteSlackHi}{0.336}
\newcommand{\SpSecBThreeSuiteSlackHit}{6}
\newcommand{\SpSecBThreeSuiteSlackN}{30}
\newcommand{\SpSecBThreeSuiteSlackASR}{0.200}
\newcommand{\SpSecBThreeSuiteSlackLo}{0.095}
\newcommand{\SpSecBThreeSuiteSlackHi}{0.373}
\newcommand{\SpSecDeiSuiteSlackHit}{6}
\newcommand{\SpSecDeiSuiteSlackN}{30}
\newcommand{\SpSecDeiSuiteSlackASR}{0.200}
\newcommand{\SpSecDeiSuiteSlackLo}{0.095}
\newcommand{\SpSecDeiSuiteSlackHi}{0.373}
\newcommand{\SpSecDeiGateSuiteSlackHit}{5}
\newcommand{\SpSecDeiGateSuiteSlackN}{30}
\newcommand{\SpSecDeiGateSuiteSlackASR}{0.167}
\newcommand{\SpSecDeiGateSuiteSlackLo}{0.073}
\newcommand{\SpSecDeiGateSuiteSlackHi}{0.336}
\newcommand{\SpSecNoDefenseSuiteSlackHit}{28}
\newcommand{\SpSecNoDefenseSuiteSlackN}{30}
\newcommand{\SpSecNoDefenseSuiteSlackASR}{0.933}
\newcommand{\SpSecNoDefenseSuiteSlackLo}{0.787}
\newcommand{\SpSecNoDefenseSuiteSlackHi}{0.982}
\newcommand{\SpSecBOneSuiteTravelHit}{7}
\newcommand{\SpSecBOneSuiteTravelN}{30}
\newcommand{\SpSecBOneSuiteTravelASR}{0.233}
\newcommand{\SpSecBOneSuiteTravelLo}{0.118}
\newcommand{\SpSecBOneSuiteTravelHi}{0.409}
\newcommand{\SpSecBTwoSuiteTravelHit}{8}
\newcommand{\SpSecBTwoSuiteTravelN}{30}
\newcommand{\SpSecBTwoSuiteTravelASR}{0.267}
\newcommand{\SpSecBTwoSuiteTravelLo}{0.142}
\newcommand{\SpSecBTwoSuiteTravelHi}{0.444}
\newcommand{\SpSecBTwoStrictSuiteTravelHit}{2}
\newcommand{\SpSecBTwoStrictSuiteTravelN}{30}
\newcommand{\SpSecBTwoStrictSuiteTravelASR}{0.067}
\newcommand{\SpSecBTwoStrictSuiteTravelLo}{0.018}
\newcommand{\SpSecBTwoStrictSuiteTravelHi}{0.213}
\newcommand{\SpSecBThreeSuiteTravelHit}{17}
\newcommand{\SpSecBThreeSuiteTravelN}{30}
\newcommand{\SpSecBThreeSuiteTravelASR}{0.567}
\newcommand{\SpSecBThreeSuiteTravelLo}{0.392}
\newcommand{\SpSecBThreeSuiteTravelHi}{0.726}
\newcommand{\SpSecDeiSuiteTravelHit}{7}
\newcommand{\SpSecDeiSuiteTravelN}{30}
\newcommand{\SpSecDeiSuiteTravelASR}{0.233}
\newcommand{\SpSecDeiSuiteTravelLo}{0.118}
\newcommand{\SpSecDeiSuiteTravelHi}{0.409}
\newcommand{\SpSecDeiGateSuiteTravelHit}{8}
\newcommand{\SpSecDeiGateSuiteTravelN}{30}
\newcommand{\SpSecDeiGateSuiteTravelASR}{0.267}
\newcommand{\SpSecDeiGateSuiteTravelLo}{0.142}
\newcommand{\SpSecDeiGateSuiteTravelHi}{0.444}
\newcommand{\SpSecNoDefenseSuiteTravelHit}{24}
\newcommand{\SpSecNoDefenseSuiteTravelN}{30}
\newcommand{\SpSecNoDefenseSuiteTravelASR}{0.800}
\newcommand{\SpSecNoDefenseSuiteTravelLo}{0.627}
\newcommand{\SpSecNoDefenseSuiteTravelHi}{0.905}
\newcommand{\SpSecBOneSuiteWorkspaceHit}{0}
\newcommand{\SpSecBOneSuiteWorkspaceN}{60}
\newcommand{\SpSecBOneSuiteWorkspaceASR}{0.000}
\newcommand{\SpSecBOneSuiteWorkspaceLo}{0.000}
\newcommand{\SpSecBOneSuiteWorkspaceHi}{0.060}
\newcommand{\SpSecBTwoSuiteWorkspaceHit}{0}
\newcommand{\SpSecBTwoSuiteWorkspaceN}{60}
\newcommand{\SpSecBTwoSuiteWorkspaceASR}{0.000}
\newcommand{\SpSecBTwoSuiteWorkspaceLo}{0.000}
\newcommand{\SpSecBTwoSuiteWorkspaceHi}{0.060}
\newcommand{\SpSecBTwoStrictSuiteWorkspaceHit}{0}
\newcommand{\SpSecBTwoStrictSuiteWorkspaceN}{60}
\newcommand{\SpSecBTwoStrictSuiteWorkspaceASR}{0.000}
\newcommand{\SpSecBTwoStrictSuiteWorkspaceLo}{0.000}
\newcommand{\SpSecBTwoStrictSuiteWorkspaceHi}{0.060}
\newcommand{\SpSecBThreeSuiteWorkspaceHit}{2}
\newcommand{\SpSecBThreeSuiteWorkspaceN}{60}
\newcommand{\SpSecBThreeSuiteWorkspaceASR}{0.033}
\newcommand{\SpSecBThreeSuiteWorkspaceLo}{0.009}
\newcommand{\SpSecBThreeSuiteWorkspaceHi}{0.114}
\newcommand{\SpSecDeiSuiteWorkspaceHit}{0}
\newcommand{\SpSecDeiSuiteWorkspaceN}{60}
\newcommand{\SpSecDeiSuiteWorkspaceASR}{0.000}
\newcommand{\SpSecDeiSuiteWorkspaceLo}{0.000}
\newcommand{\SpSecDeiSuiteWorkspaceHi}{0.060}
\newcommand{\SpSecDeiGateSuiteWorkspaceHit}{0}
\newcommand{\SpSecDeiGateSuiteWorkspaceN}{60}
\newcommand{\SpSecDeiGateSuiteWorkspaceASR}{0.000}
\newcommand{\SpSecDeiGateSuiteWorkspaceLo}{0.000}
\newcommand{\SpSecDeiGateSuiteWorkspaceHi}{0.060}
\newcommand{\SpSecNoDefenseSuiteWorkspaceHit}{15}
\newcommand{\SpSecNoDefenseSuiteWorkspaceN}{60}
\newcommand{\SpSecNoDefenseSuiteWorkspaceASR}{0.250}
\newcommand{\SpSecNoDefenseSuiteWorkspaceLo}{0.158}
\newcommand{\SpSecNoDefenseSuiteWorkspaceHi}{0.372}
\newcommand{\SpSecBOneTauTuningHit}{0}
\newcommand{\SpSecBOneTauTuningN}{32}
\newcommand{\SpSecBOneTauTuningASR}{0.000}
\newcommand{\SpSecBOneTauTuningLo}{0.000}
\newcommand{\SpSecBOneTauTuningHi}{0.107}
\newcommand{\SpSecBTwoTauTuningHit}{0}
\newcommand{\SpSecBTwoTauTuningN}{32}
\newcommand{\SpSecBTwoTauTuningASR}{0.000}
\newcommand{\SpSecBTwoTauTuningLo}{0.000}
\newcommand{\SpSecBTwoTauTuningHi}{0.107}
\newcommand{\SpSecBTwoStrictTauTuningHit}{0}
\newcommand{\SpSecBTwoStrictTauTuningN}{32}
\newcommand{\SpSecBTwoStrictTauTuningASR}{0.000}
\newcommand{\SpSecBTwoStrictTauTuningLo}{0.000}
\newcommand{\SpSecBTwoStrictTauTuningHi}{0.107}
\newcommand{\SpSecBThreeTauTuningHit}{1}
\newcommand{\SpSecBThreeTauTuningN}{32}
\newcommand{\SpSecBThreeTauTuningASR}{0.031}
\newcommand{\SpSecBThreeTauTuningLo}{0.006}
\newcommand{\SpSecBThreeTauTuningHi}{0.157}
\newcommand{\SpSecDeiTauTuningHit}{0}
\newcommand{\SpSecDeiTauTuningN}{32}
\newcommand{\SpSecDeiTauTuningASR}{0.000}
\newcommand{\SpSecDeiTauTuningLo}{0.000}
\newcommand{\SpSecDeiTauTuningHi}{0.107}
\newcommand{\SpSecDeiGateTauTuningHit}{0}
\newcommand{\SpSecDeiGateTauTuningN}{32}
\newcommand{\SpSecDeiGateTauTuningASR}{0.000}
\newcommand{\SpSecDeiGateTauTuningLo}{0.000}
\newcommand{\SpSecDeiGateTauTuningHi}{0.107}
\newcommand{\SpSecNoDefenseTauTuningHit}{7}
\newcommand{\SpSecNoDefenseTauTuningN}{32}
\newcommand{\SpSecNoDefenseTauTuningASR}{0.219}
\newcommand{\SpSecNoDefenseTauTuningLo}{0.110}
\newcommand{\SpSecNoDefenseTauTuningHi}{0.388}
\newcommand{\SpSecBOneTauUntouchedHit}{0}
\newcommand{\SpSecBOneTauUntouchedN}{28}
\newcommand{\SpSecBOneTauUntouchedASR}{0.000}
\newcommand{\SpSecBOneTauUntouchedLo}{0.000}
\newcommand{\SpSecBOneTauUntouchedHi}{0.121}
\newcommand{\SpSecBTwoTauUntouchedHit}{0}
\newcommand{\SpSecBTwoTauUntouchedN}{28}
\newcommand{\SpSecBTwoTauUntouchedASR}{0.000}
\newcommand{\SpSecBTwoTauUntouchedLo}{0.000}
\newcommand{\SpSecBTwoTauUntouchedHi}{0.121}
\newcommand{\SpSecBTwoStrictTauUntouchedHit}{0}
\newcommand{\SpSecBTwoStrictTauUntouchedN}{28}
\newcommand{\SpSecBTwoStrictTauUntouchedASR}{0.000}
\newcommand{\SpSecBTwoStrictTauUntouchedLo}{0.000}
\newcommand{\SpSecBTwoStrictTauUntouchedHi}{0.121}
\newcommand{\SpSecBThreeTauUntouchedHit}{1}
\newcommand{\SpSecBThreeTauUntouchedN}{28}
\newcommand{\SpSecBThreeTauUntouchedASR}{0.036}
\newcommand{\SpSecBThreeTauUntouchedLo}{0.006}
\newcommand{\SpSecBThreeTauUntouchedHi}{0.177}
\newcommand{\SpSecDeiTauUntouchedHit}{0}
\newcommand{\SpSecDeiTauUntouchedN}{28}
\newcommand{\SpSecDeiTauUntouchedASR}{0.000}
\newcommand{\SpSecDeiTauUntouchedLo}{0.000}
\newcommand{\SpSecDeiTauUntouchedHi}{0.121}
\newcommand{\SpSecDeiGateTauUntouchedHit}{0}
\newcommand{\SpSecDeiGateTauUntouchedN}{28}
\newcommand{\SpSecDeiGateTauUntouchedASR}{0.000}
\newcommand{\SpSecDeiGateTauUntouchedLo}{0.000}
\newcommand{\SpSecDeiGateTauUntouchedHi}{0.121}
\newcommand{\SpSecNoDefenseTauUntouchedHit}{8}
\newcommand{\SpSecNoDefenseTauUntouchedN}{28}
\newcommand{\SpSecNoDefenseTauUntouchedASR}{0.286}
\newcommand{\SpSecNoDefenseTauUntouchedLo}{0.153}
\newcommand{\SpSecNoDefenseTauUntouchedHi}{0.471}
\newcommand{\SpUtilNoDefenseFail}{21}
\newcommand{\SpUtilNoDefenseN}{56}
\newcommand{\SpUtilNoDefenseRate}{62.5\%}
\newcommand{\SpUtilNoDefenseBlockRate}{0.375}
\newcommand{\SpUtilNoDefenseBlockLo}{0.260}
\newcommand{\SpUtilNoDefenseBlockHi}{0.506}
\newcommand{\SpUtilNoDefenseQuarEpisodes}{0}
\newcommand{\SpUtilNoDefenseQuarRate}{0.000}
\newcommand{\SpUtilNoDefenseQuarLo}{0.000}
\newcommand{\SpUtilNoDefenseQuarHi}{0.064}
\newcommand{\SpUtilNoDefenseFailDefense}{0}
\newcommand{\SpUtilNoDefenseFailModel}{21}
\newcommand{\SpUtilNoDefenseFailDegenerate}{0}
\newcommand{\SpUtilDeiGateFail}{40}
\newcommand{\SpUtilDeiGateN}{56}
\newcommand{\SpUtilDeiGateRate}{28.6\%}
\newcommand{\SpUtilDeiGateBlockRate}{0.714}
\newcommand{\SpUtilDeiGateBlockLo}{0.585}
\newcommand{\SpUtilDeiGateBlockHi}{0.816}
\newcommand{\SpUtilDeiGateQuarEpisodes}{4}
\newcommand{\SpUtilDeiGateQuarRate}{0.071}
\newcommand{\SpUtilDeiGateQuarLo}{0.028}
\newcommand{\SpUtilDeiGateQuarHi}{0.170}
\newcommand{\SpUtilDeiGateFailDefense}{19}
\newcommand{\SpUtilDeiGateFailModel}{21}
\newcommand{\SpUtilDeiGateFailDegenerate}{0}
\newcommand{\SpUtilDeiGateFisherP}{0.00}
\newcommand{\SpUtilDeiGateMcNemarP}{0.000}
\newcommand{\SpUtilDeiGateDiscordantB}{20}
\newcommand{\SpUtilDeiGateDiscordantC}{1}
\newcommand{\SpUtilDeiGatePairedN}{56}
\newcommand{\SpUtilDeiGatePairedRD}{-0.339}
\newcommand{\SpUtilDeiGatePairedRDlo}{-0.458}
\newcommand{\SpUtilDeiGatePairedRDhi}{-0.198}
\newcommand{\SpUtilDeiGateClusterN}{56}
\newcommand{\SpUtilDeiGateClusterRDlo}{-0.464}
\newcommand{\SpUtilDeiGateClusterRDhi}{-0.214}
\newcommand{\SpUtilBThreeFail}{39}
\newcommand{\SpUtilBThreeN}{56}
\newcommand{\SpUtilBThreeRate}{30.4\%}
\newcommand{\SpUtilBThreeBlockRate}{0.696}
\newcommand{\SpUtilBThreeBlockLo}{0.567}
\newcommand{\SpUtilBThreeBlockHi}{0.801}
\newcommand{\SpUtilBThreeQuarEpisodes}{0}
\newcommand{\SpUtilBThreeQuarRate}{0.000}
\newcommand{\SpUtilBThreeQuarLo}{0.000}
\newcommand{\SpUtilBThreeQuarHi}{0.064}
\newcommand{\SpUtilBThreeFailDefense}{15}
\newcommand{\SpUtilBThreeFailModel}{24}
\newcommand{\SpUtilBThreeFailDegenerate}{0}
\newcommand{\SpUtilBThreeFisherP}{0.00}
\newcommand{\SpUtilBThreeMcNemarP}{0.000}
\newcommand{\SpUtilBThreeDiscordantB}{19}
\newcommand{\SpUtilBThreeDiscordantC}{1}
\newcommand{\SpUtilBThreePairedN}{56}
\newcommand{\SpUtilBThreePairedRD}{-0.321}
\newcommand{\SpUtilBThreePairedRDlo}{-0.439}
\newcommand{\SpUtilBThreePairedRDhi}{-0.182}
\newcommand{\SpUtilBThreeClusterN}{56}
\newcommand{\SpUtilBThreeClusterRDlo}{-0.446}
\newcommand{\SpUtilBThreeClusterRDhi}{-0.196}
\newcommand{\SpUtilBTwoStrictFail}{43}
\newcommand{\SpUtilBTwoStrictN}{56}
\newcommand{\SpUtilBTwoStrictRate}{23.2\%}
\newcommand{\SpUtilBTwoStrictBlockRate}{0.768}
\newcommand{\SpUtilBTwoStrictBlockLo}{0.642}
\newcommand{\SpUtilBTwoStrictBlockHi}{0.859}
\newcommand{\SpUtilBTwoStrictQuarEpisodes}{0}
\newcommand{\SpUtilBTwoStrictQuarRate}{0.000}
\newcommand{\SpUtilBTwoStrictQuarLo}{0.000}
\newcommand{\SpUtilBTwoStrictQuarHi}{0.064}
\newcommand{\SpUtilBTwoStrictFailDefense}{26}
\newcommand{\SpUtilBTwoStrictFailModel}{17}
\newcommand{\SpUtilBTwoStrictFailDegenerate}{0}
\newcommand{\SpUtilBTwoStrictFisherP}{0.00}
\newcommand{\SpUtilBTwoStrictMcNemarP}{0.000}
\newcommand{\SpUtilBTwoStrictDiscordantB}{22}
\newcommand{\SpUtilBTwoStrictDiscordantC}{0}
\newcommand{\SpUtilBTwoStrictPairedN}{56}
\newcommand{\SpUtilBTwoStrictPairedRD}{-0.393}
\newcommand{\SpUtilBTwoStrictPairedRDlo}{-0.505}
\newcommand{\SpUtilBTwoStrictPairedRDhi}{-0.255}
\newcommand{\SpUtilBTwoStrictClusterN}{56}
\newcommand{\SpUtilBTwoStrictClusterRDlo}{-0.518}
\newcommand{\SpUtilBTwoStrictClusterRDhi}{-0.268}
\newcommand{\SpUtilBTwoFail}{34}
\newcommand{\SpUtilBTwoN}{56}
\newcommand{\SpUtilBTwoRate}{39.3\%}
\newcommand{\SpUtilBTwoBlockRate}{0.607}
\newcommand{\SpUtilBTwoBlockLo}{0.476}
\newcommand{\SpUtilBTwoBlockHi}{0.724}
\newcommand{\SpUtilBTwoQuarEpisodes}{0}
\newcommand{\SpUtilBTwoQuarRate}{0.000}
\newcommand{\SpUtilBTwoQuarLo}{0.000}
\newcommand{\SpUtilBTwoQuarHi}{0.064}
\newcommand{\SpUtilBTwoFailDefense}{16}
\newcommand{\SpUtilBTwoFailModel}{18}
\newcommand{\SpUtilBTwoFailDegenerate}{0}
\newcommand{\SpUtilBTwoFisherP}{0.02}
\newcommand{\SpUtilBTwoMcNemarP}{0.001}
\newcommand{\SpUtilBTwoDiscordantB}{14}
\newcommand{\SpUtilBTwoDiscordantC}{1}
\newcommand{\SpUtilBTwoPairedN}{56}
\newcommand{\SpUtilBTwoPairedRD}{-0.232}
\newcommand{\SpUtilBTwoPairedRDlo}{-0.342}
\newcommand{\SpUtilBTwoPairedRDhi}{-0.108}
\newcommand{\SpUtilBTwoClusterN}{56}
\newcommand{\SpUtilBTwoClusterRDlo}{-0.357}
\newcommand{\SpUtilBTwoClusterRDhi}{-0.107}
\newcommand{\SpUtilDeiFail}{35}
\newcommand{\SpUtilDeiN}{56}
\newcommand{\SpUtilDeiRate}{37.5\%}
\newcommand{\SpUtilDeiBlockRate}{0.625}
\newcommand{\SpUtilDeiBlockLo}{0.494}
\newcommand{\SpUtilDeiBlockHi}{0.740}
\newcommand{\SpUtilDeiQuarEpisodes}{0}
\newcommand{\SpUtilDeiQuarRate}{0.000}
\newcommand{\SpUtilDeiQuarLo}{0.000}
\newcommand{\SpUtilDeiQuarHi}{0.064}
\newcommand{\SpUtilDeiFailDefense}{16}
\newcommand{\SpUtilDeiFailModel}{19}
\newcommand{\SpUtilDeiFailDegenerate}{0}
\newcommand{\SpUtilDeiFisherP}{0.01}
\newcommand{\SpUtilDeiMcNemarP}{0.004}
\newcommand{\SpUtilDeiDiscordantB}{18}
\newcommand{\SpUtilDeiDiscordantC}{4}
\newcommand{\SpUtilDeiPairedN}{56}
\newcommand{\SpUtilDeiPairedRD}{-0.250}
\newcommand{\SpUtilDeiPairedRDlo}{-0.386}
\newcommand{\SpUtilDeiPairedRDhi}{-0.095}
\newcommand{\SpUtilDeiClusterN}{56}
\newcommand{\SpUtilDeiClusterRDlo}{-0.393}
\newcommand{\SpUtilDeiClusterRDhi}{-0.107}
\newcommand{\SpUtilBOneFail}{39}
\newcommand{\SpUtilBOneN}{56}
\newcommand{\SpUtilBOneRate}{30.4\%}
\newcommand{\SpUtilBOneBlockRate}{0.696}
\newcommand{\SpUtilBOneBlockLo}{0.567}
\newcommand{\SpUtilBOneBlockHi}{0.801}
\newcommand{\SpUtilBOneQuarEpisodes}{0}
\newcommand{\SpUtilBOneQuarRate}{0.000}
\newcommand{\SpUtilBOneQuarLo}{0.000}
\newcommand{\SpUtilBOneQuarHi}{0.064}
\newcommand{\SpUtilBOneFailDefense}{20}
\newcommand{\SpUtilBOneFailModel}{19}
\newcommand{\SpUtilBOneFailDegenerate}{0}
\newcommand{\SpUtilBOneFisherP}{0.00}
\newcommand{\SpUtilBOneMcNemarP}{0.000}
\newcommand{\SpUtilBOneDiscordantB}{19}
\newcommand{\SpUtilBOneDiscordantC}{1}
\newcommand{\SpUtilBOnePairedN}{56}
\newcommand{\SpUtilBOnePairedRD}{-0.321}
\newcommand{\SpUtilBOnePairedRDlo}{-0.439}
\newcommand{\SpUtilBOnePairedRDhi}{-0.182}
\newcommand{\SpUtilBOneClusterN}{56}
\newcommand{\SpUtilBOneClusterRDlo}{-0.446}
\newcommand{\SpUtilBOneClusterRDhi}{-0.196}
\newcommand{\SpUtilDeiGateVsBTwoStrictPairedN}{56}
\newcommand{\SpUtilDeiGateVsBTwoStrictPairedRD}{0.054}
\newcommand{\SpUtilDeiGateVsBTwoStrictPairedRDlo}{-0.026}
\newcommand{\SpUtilDeiGateVsBTwoStrictPairedRDhi}{0.135}
\newcommand{\SpUtilDeiGateVsBTwoStrictMcNemarP}{0.375}
\newcommand{\SpUtilDeiGateVsBTwoStrictHolmP}{0.750}
\newcommand{\SpUtilDeiGateVsBTwoStrictClusterN}{56}
\newcommand{\SpUtilDeiGateVsBTwoStrictClusterRDlo}{-0.018}
\newcommand{\SpUtilDeiGateVsBTwoStrictClusterRDhi}{0.125}
\newcommand{\SpUtilDeiGateVsBOnePairedN}{56}
\newcommand{\SpUtilDeiGateVsBOnePairedRD}{-0.018}
\newcommand{\SpUtilDeiGateVsBOnePairedRDlo}{-0.111}
\newcommand{\SpUtilDeiGateVsBOnePairedRDhi}{0.075}
\newcommand{\SpUtilDeiGateVsBOneMcNemarP}{1.000}
\newcommand{\SpUtilDeiGateVsBOneHolmP}{1.000}
\newcommand{\SpUtilDeiGateVsBOneClusterN}{56}
\newcommand{\SpUtilDeiGateVsBOneClusterRDlo}{-0.107}
\newcommand{\SpUtilDeiGateVsBOneClusterRDhi}{0.071}
\newcommand{\SpUtilBOneProtectedUseFail}{22}
\newcommand{\SpUtilBOneProtectedUseN}{27}
\newcommand{\SpUtilBOneProtectedUseBlockRate}{0.815}
\newcommand{\SpUtilBOneProtectedUseBlockLo}{0.633}
\newcommand{\SpUtilBOneProtectedUseBlockHi}{0.918}
\newcommand{\SpUtilBTwoProtectedUseFail}{17}
\newcommand{\SpUtilBTwoProtectedUseN}{27}
\newcommand{\SpUtilBTwoProtectedUseBlockRate}{0.630}
\newcommand{\SpUtilBTwoProtectedUseBlockLo}{0.442}
\newcommand{\SpUtilBTwoProtectedUseBlockHi}{0.785}
\newcommand{\SpUtilBTwoStrictProtectedUseFail}{26}
\newcommand{\SpUtilBTwoStrictProtectedUseN}{26}
\newcommand{\SpUtilBTwoStrictProtectedUseBlockRate}{1.000}
\newcommand{\SpUtilBTwoStrictProtectedUseBlockLo}{0.871}
\newcommand{\SpUtilBTwoStrictProtectedUseBlockHi}{1.000}
\newcommand{\SpUtilBThreeProtectedUseFail}{17}
\newcommand{\SpUtilBThreeProtectedUseN}{24}
\newcommand{\SpUtilBThreeProtectedUseBlockRate}{0.708}
\newcommand{\SpUtilBThreeProtectedUseBlockLo}{0.508}
\newcommand{\SpUtilBThreeProtectedUseBlockHi}{0.851}
\newcommand{\SpUtilDeiProtectedUseFail}{19}
\newcommand{\SpUtilDeiProtectedUseN}{26}
\newcommand{\SpUtilDeiProtectedUseBlockRate}{0.731}
\newcommand{\SpUtilDeiProtectedUseBlockLo}{0.539}
\newcommand{\SpUtilDeiProtectedUseBlockHi}{0.863}
\newcommand{\SpUtilDeiGateProtectedUseFail}{22}
\newcommand{\SpUtilDeiGateProtectedUseN}{25}
\newcommand{\SpUtilDeiGateProtectedUseBlockRate}{0.880}
\newcommand{\SpUtilDeiGateProtectedUseBlockLo}{0.700}
\newcommand{\SpUtilDeiGateProtectedUseBlockHi}{0.958}
\newcommand{\SpUtilNoDefenseProtectedUseFail}{4}
\newcommand{\SpUtilNoDefenseProtectedUseN}{24}
\newcommand{\SpUtilNoDefenseProtectedUseBlockRate}{0.167}
\newcommand{\SpUtilNoDefenseProtectedUseBlockLo}{0.067}
\newcommand{\SpUtilNoDefenseProtectedUseBlockHi}{0.359}
\newcommand{\SpUtilBOneNoProtectedUseFail}{17}
\newcommand{\SpUtilBOneNoProtectedUseN}{29}
\newcommand{\SpUtilBOneNoProtectedUseBlockRate}{0.586}
\newcommand{\SpUtilBOneNoProtectedUseBlockLo}{0.407}
\newcommand{\SpUtilBOneNoProtectedUseBlockHi}{0.745}
\newcommand{\SpUtilBTwoNoProtectedUseFail}{17}
\newcommand{\SpUtilBTwoNoProtectedUseN}{29}
\newcommand{\SpUtilBTwoNoProtectedUseBlockRate}{0.586}
\newcommand{\SpUtilBTwoNoProtectedUseBlockLo}{0.407}
\newcommand{\SpUtilBTwoNoProtectedUseBlockHi}{0.745}
\newcommand{\SpUtilBTwoStrictNoProtectedUseFail}{17}
\newcommand{\SpUtilBTwoStrictNoProtectedUseN}{30}
\newcommand{\SpUtilBTwoStrictNoProtectedUseBlockRate}{0.567}
\newcommand{\SpUtilBTwoStrictNoProtectedUseBlockLo}{0.392}
\newcommand{\SpUtilBTwoStrictNoProtectedUseBlockHi}{0.726}
\newcommand{\SpUtilBThreeNoProtectedUseFail}{22}
\newcommand{\SpUtilBThreeNoProtectedUseN}{32}
\newcommand{\SpUtilBThreeNoProtectedUseBlockRate}{0.688}
\newcommand{\SpUtilBThreeNoProtectedUseBlockLo}{0.514}
\newcommand{\SpUtilBThreeNoProtectedUseBlockHi}{0.820}
\newcommand{\SpUtilDeiNoProtectedUseFail}{16}
\newcommand{\SpUtilDeiNoProtectedUseN}{30}
\newcommand{\SpUtilDeiNoProtectedUseBlockRate}{0.533}
\newcommand{\SpUtilDeiNoProtectedUseBlockLo}{0.361}
\newcommand{\SpUtilDeiNoProtectedUseBlockHi}{0.698}
\newcommand{\SpUtilDeiGateNoProtectedUseFail}{18}
\newcommand{\SpUtilDeiGateNoProtectedUseN}{31}
\newcommand{\SpUtilDeiGateNoProtectedUseBlockRate}{0.581}
\newcommand{\SpUtilDeiGateNoProtectedUseBlockLo}{0.408}
\newcommand{\SpUtilDeiGateNoProtectedUseBlockHi}{0.736}
\newcommand{\SpUtilNoDefenseNoProtectedUseFail}{17}
\newcommand{\SpUtilNoDefenseNoProtectedUseN}{32}
\newcommand{\SpUtilNoDefenseNoProtectedUseBlockRate}{0.531}
\newcommand{\SpUtilNoDefenseNoProtectedUseBlockLo}{0.364}
\newcommand{\SpUtilNoDefenseNoProtectedUseBlockHi}{0.691}
\newcommand{\SpUtilBOneDevelopmentUntouchedFail}{39}
\newcommand{\SpUtilBOneDevelopmentUntouchedN}{56}
\newcommand{\SpUtilBOneDevelopmentUntouchedBlockRate}{0.696}
\newcommand{\SpUtilBOneDevelopmentUntouchedBlockLo}{0.567}
\newcommand{\SpUtilBOneDevelopmentUntouchedBlockHi}{0.801}
\newcommand{\SpUtilBTwoDevelopmentUntouchedFail}{34}
\newcommand{\SpUtilBTwoDevelopmentUntouchedN}{56}
\newcommand{\SpUtilBTwoDevelopmentUntouchedBlockRate}{0.607}
\newcommand{\SpUtilBTwoDevelopmentUntouchedBlockLo}{0.476}
\newcommand{\SpUtilBTwoDevelopmentUntouchedBlockHi}{0.724}
\newcommand{\SpUtilBTwoStrictDevelopmentUntouchedFail}{43}
\newcommand{\SpUtilBTwoStrictDevelopmentUntouchedN}{56}
\newcommand{\SpUtilBTwoStrictDevelopmentUntouchedBlockRate}{0.768}
\newcommand{\SpUtilBTwoStrictDevelopmentUntouchedBlockLo}{0.642}
\newcommand{\SpUtilBTwoStrictDevelopmentUntouchedBlockHi}{0.859}
\newcommand{\SpUtilBThreeDevelopmentUntouchedFail}{39}
\newcommand{\SpUtilBThreeDevelopmentUntouchedN}{56}
\newcommand{\SpUtilBThreeDevelopmentUntouchedBlockRate}{0.696}
\newcommand{\SpUtilBThreeDevelopmentUntouchedBlockLo}{0.567}
\newcommand{\SpUtilBThreeDevelopmentUntouchedBlockHi}{0.801}
\newcommand{\SpUtilDeiDevelopmentUntouchedFail}{35}
\newcommand{\SpUtilDeiDevelopmentUntouchedN}{56}
\newcommand{\SpUtilDeiDevelopmentUntouchedBlockRate}{0.625}
\newcommand{\SpUtilDeiDevelopmentUntouchedBlockLo}{0.494}
\newcommand{\SpUtilDeiDevelopmentUntouchedBlockHi}{0.740}
\newcommand{\SpUtilDeiGateDevelopmentUntouchedFail}{40}
\newcommand{\SpUtilDeiGateDevelopmentUntouchedN}{56}
\newcommand{\SpUtilDeiGateDevelopmentUntouchedBlockRate}{0.714}
\newcommand{\SpUtilDeiGateDevelopmentUntouchedBlockLo}{0.585}
\newcommand{\SpUtilDeiGateDevelopmentUntouchedBlockHi}{0.816}
\newcommand{\SpUtilNoDefenseDevelopmentUntouchedFail}{21}
\newcommand{\SpUtilNoDefenseDevelopmentUntouchedN}{56}
\newcommand{\SpUtilNoDefenseDevelopmentUntouchedBlockRate}{0.375}
\newcommand{\SpUtilNoDefenseDevelopmentUntouchedBlockLo}{0.260}
\newcommand{\SpUtilNoDefenseDevelopmentUntouchedBlockHi}{0.506}
\newcommand{\SpSweepTauThirtySecAdmit}{39}
\newcommand{\SpSweepTauThirtySecN}{184}
\newcommand{\SpSweepTauThirtySecRate}{0.212}
\newcommand{\SpSweepTauThirtyBenignAdmit}{22}
\newcommand{\SpSweepTauThirtyBenignN}{51}
\newcommand{\SpSweepTauThirtyBenignRate}{0.431}
\newcommand{\SpSweepTauFortyFiveSecAdmit}{39}
\newcommand{\SpSweepTauFortyFiveSecN}{184}
\newcommand{\SpSweepTauFortyFiveSecRate}{0.212}
\newcommand{\SpSweepTauFortyFiveBenignAdmit}{22}
\newcommand{\SpSweepTauFortyFiveBenignN}{51}
\newcommand{\SpSweepTauFortyFiveBenignRate}{0.431}
\newcommand{\SpSweepTauFiftySecAdmit}{39}
\newcommand{\SpSweepTauFiftySecN}{184}
\newcommand{\SpSweepTauFiftySecRate}{0.212}
\newcommand{\SpSweepTauFiftyBenignAdmit}{22}
\newcommand{\SpSweepTauFiftyBenignN}{51}
\newcommand{\SpSweepTauFiftyBenignRate}{0.431}
\newcommand{\SpSweepTauSixtySecAdmit}{0}
\newcommand{\SpSweepTauSixtySecN}{184}
\newcommand{\SpSweepTauSixtySecRate}{0.000}
\newcommand{\SpSweepTauSixtyBenignAdmit}{0}
\newcommand{\SpSweepTauSixtyBenignN}{51}
\newcommand{\SpSweepTauSixtyBenignRate}{0.000}
\newcommand{\SpSweepTauSixtySevenSecAdmit}{0}
\newcommand{\SpSweepTauSixtySevenSecN}{184}
\newcommand{\SpSweepTauSixtySevenSecRate}{0.000}
\newcommand{\SpSweepTauSixtySevenBenignAdmit}{0}
\newcommand{\SpSweepTauSixtySevenBenignN}{51}
\newcommand{\SpSweepTauSixtySevenBenignRate}{0.000}
\newcommand{\SpSweepEllAllTrustedSecAdmit}{19}
\newcommand{\SpSweepEllAllTrustedSecN}{184}
\newcommand{\SpSweepEllAllTrustedSecRate}{0.103}
\newcommand{\SpSweepEllAllTrustedBenignAdmit}{22}
\newcommand{\SpSweepEllAllTrustedBenignN}{51}
\newcommand{\SpSweepEllAllTrustedBenignRate}{0.431}
  % registered seven-arm campaign; regenerate via analyze.py
% AUTO-GENERATED by analyze.py from out/rq2_probes.jsonl; primary (degenerate retained); do not edit.
% Voided: 0; marked degenerate in source: 0; bootstrap seed: 20260814.
\newcommand{\RqTwoEpisodes}{640}
\newcommand{\RqTwoVoided}{0}
\newcommand{\RqTwoDegenerate}{0}
\newcommand{\RqTwoAnalysisMode}{primary}
\newcommand{\RqTwoSecNoDefenseHit}{28}
\newcommand{\RqTwoSecNoDefenseN}{80}
\newcommand{\RqTwoSecNoDefenseASR}{0.350}
\newcommand{\RqTwoSecNoDefenseLo}{0.255}
\newcommand{\RqTwoSecNoDefenseHi}{0.459}
\newcommand{\RqTwoSecDeiGateHit}{17}
\newcommand{\RqTwoSecDeiGateN}{80}
\newcommand{\RqTwoSecDeiGateASR}{0.212}
\newcommand{\RqTwoSecDeiGateLo}{0.137}
\newcommand{\RqTwoSecDeiGateHi}{0.314}
\newcommand{\RqTwoSecDeiGateFisherP}{7.8e-02}
\newcommand{\RqTwoSecDeiGateRD}{-0.137}
\newcommand{\RqTwoSecDeiGateRDlo}{-0.275}
\newcommand{\RqTwoSecDeiGateRDhi}{0.000}
\newcommand{\RqTwoSecDeiGateRR}{0.607}
\newcommand{\RqTwoSecDeiGateRRlo}{0.362}
\newcommand{\RqTwoSecDeiGateRRhi}{1.018}
\newcommand{\RqTwoSecBOneHit}{32}
\newcommand{\RqTwoSecBOneN}{80}
\newcommand{\RqTwoSecBOneASR}{0.400}
\newcommand{\RqTwoSecBOneLo}{0.300}
\newcommand{\RqTwoSecBOneHi}{0.510}
\newcommand{\RqTwoSecBOneFisherP}{6.2e-01}
\newcommand{\RqTwoSecBOneRD}{0.050}
\newcommand{\RqTwoSecBOneRDlo}{-0.100}
\newcommand{\RqTwoSecBOneRDhi}{0.200}
\newcommand{\RqTwoSecBOneRR}{1.143}
\newcommand{\RqTwoSecBOneRRlo}{0.765}
\newcommand{\RqTwoSecBOneRRhi}{1.708}
\newcommand{\RqTwoSecDeiHit}{32}
\newcommand{\RqTwoSecDeiN}{80}
\newcommand{\RqTwoSecDeiASR}{0.400}
\newcommand{\RqTwoSecDeiLo}{0.300}
\newcommand{\RqTwoSecDeiHi}{0.510}
\newcommand{\RqTwoSecDeiFisherP}{6.2e-01}
\newcommand{\RqTwoSecDeiRD}{0.050}
\newcommand{\RqTwoSecDeiRDlo}{-0.100}
\newcommand{\RqTwoSecDeiRDhi}{0.200}
\newcommand{\RqTwoSecDeiRR}{1.143}
\newcommand{\RqTwoSecDeiRRlo}{0.765}
\newcommand{\RqTwoSecDeiRRhi}{1.708}
\newcommand{\RqTwoSecBOneSuiteWorkspaceHit}{32}
\newcommand{\RqTwoSecBOneSuiteWorkspaceN}{80}
\newcommand{\RqTwoSecBOneSuiteWorkspaceASR}{0.400}
\newcommand{\RqTwoSecBOneSuiteWorkspaceLo}{0.300}
\newcommand{\RqTwoSecBOneSuiteWorkspaceHi}{0.510}
\newcommand{\RqTwoSecDeiSuiteWorkspaceHit}{32}
\newcommand{\RqTwoSecDeiSuiteWorkspaceN}{80}
\newcommand{\RqTwoSecDeiSuiteWorkspaceASR}{0.400}
\newcommand{\RqTwoSecDeiSuiteWorkspaceLo}{0.300}
\newcommand{\RqTwoSecDeiSuiteWorkspaceHi}{0.510}
\newcommand{\RqTwoSecDeiGateSuiteWorkspaceHit}{17}
\newcommand{\RqTwoSecDeiGateSuiteWorkspaceN}{80}
\newcommand{\RqTwoSecDeiGateSuiteWorkspaceASR}{0.212}
\newcommand{\RqTwoSecDeiGateSuiteWorkspaceLo}{0.137}
\newcommand{\RqTwoSecDeiGateSuiteWorkspaceHi}{0.314}
\newcommand{\RqTwoSecNoDefenseSuiteWorkspaceHit}{28}
\newcommand{\RqTwoSecNoDefenseSuiteWorkspaceN}{80}
\newcommand{\RqTwoSecNoDefenseSuiteWorkspaceASR}{0.350}
\newcommand{\RqTwoSecNoDefenseSuiteWorkspaceLo}{0.255}
\newcommand{\RqTwoSecNoDefenseSuiteWorkspaceHi}{0.459}
\newcommand{\RqTwoUtilNoDefenseFail}{36}
\newcommand{\RqTwoUtilNoDefenseN}{80}
\newcommand{\RqTwoUtilNoDefenseRate}{55.0\%}
\newcommand{\RqTwoUtilNoDefenseBlockRate}{0.450}
\newcommand{\RqTwoUtilNoDefenseBlockLo}{0.346}
\newcommand{\RqTwoUtilNoDefenseBlockHi}{0.559}
\newcommand{\RqTwoUtilNoDefenseQuarEpisodes}{0}
\newcommand{\RqTwoUtilNoDefenseQuarRate}{0.000}
\newcommand{\RqTwoUtilNoDefenseQuarLo}{0.000}
\newcommand{\RqTwoUtilNoDefenseQuarHi}{0.046}
\newcommand{\RqTwoUtilNoDefenseFailDefense}{0}
\newcommand{\RqTwoUtilNoDefenseFailModel}{36}
\newcommand{\RqTwoUtilNoDefenseFailDegenerate}{0}
\newcommand{\RqTwoUtilDeiFail}{37}
\newcommand{\RqTwoUtilDeiN}{80}
\newcommand{\RqTwoUtilDeiRate}{53.8\%}
\newcommand{\RqTwoUtilDeiBlockRate}{0.463}
\newcommand{\RqTwoUtilDeiBlockLo}{0.357}
\newcommand{\RqTwoUtilDeiBlockHi}{0.571}
\newcommand{\RqTwoUtilDeiQuarEpisodes}{0}
\newcommand{\RqTwoUtilDeiQuarRate}{0.000}
\newcommand{\RqTwoUtilDeiQuarLo}{0.000}
\newcommand{\RqTwoUtilDeiQuarHi}{0.046}
\newcommand{\RqTwoUtilDeiFailDefense}{0}
\newcommand{\RqTwoUtilDeiFailModel}{37}
\newcommand{\RqTwoUtilDeiFailDegenerate}{0}
\newcommand{\RqTwoUtilDeiFisherP}{1.00}
\newcommand{\RqTwoUtilDeiMcNemarP}{1.000}
\newcommand{\RqTwoUtilDeiDiscordantB}{10}
\newcommand{\RqTwoUtilDeiDiscordantC}{9}
\newcommand{\RqTwoUtilDeiPairedN}{80}
\newcommand{\RqTwoUtilDeiPairedRD}{-0.013}
\newcommand{\RqTwoUtilDeiPairedRDlo}{-0.117}
\newcommand{\RqTwoUtilDeiPairedRDhi}{0.092}
\newcommand{\RqTwoUtilDeiClusterN}{20}
\newcommand{\RqTwoUtilDeiClusterRDlo}{-0.100}
\newcommand{\RqTwoUtilDeiClusterRDhi}{0.075}
\newcommand{\RqTwoUtilDeiGateFail}{39}
\newcommand{\RqTwoUtilDeiGateN}{80}
\newcommand{\RqTwoUtilDeiGateRate}{51.2\%}
\newcommand{\RqTwoUtilDeiGateBlockRate}{0.487}
\newcommand{\RqTwoUtilDeiGateBlockLo}{0.381}
\newcommand{\RqTwoUtilDeiGateBlockHi}{0.595}
\newcommand{\RqTwoUtilDeiGateQuarEpisodes}{4}
\newcommand{\RqTwoUtilDeiGateQuarRate}{0.050}
\newcommand{\RqTwoUtilDeiGateQuarLo}{0.020}
\newcommand{\RqTwoUtilDeiGateQuarHi}{0.122}
\newcommand{\RqTwoUtilDeiGateFailDefense}{4}
\newcommand{\RqTwoUtilDeiGateFailModel}{35}
\newcommand{\RqTwoUtilDeiGateFailDegenerate}{0}
\newcommand{\RqTwoUtilDeiGateFisherP}{0.75}
\newcommand{\RqTwoUtilDeiGateMcNemarP}{0.629}
\newcommand{\RqTwoUtilDeiGateDiscordantB}{10}
\newcommand{\RqTwoUtilDeiGateDiscordantC}{7}
\newcommand{\RqTwoUtilDeiGatePairedN}{80}
\newcommand{\RqTwoUtilDeiGatePairedRD}{-0.038}
\newcommand{\RqTwoUtilDeiGatePairedRDlo}{-0.135}
\newcommand{\RqTwoUtilDeiGatePairedRDhi}{0.062}
\newcommand{\RqTwoUtilDeiGateClusterN}{20}
\newcommand{\RqTwoUtilDeiGateClusterRDlo}{-0.150}
\newcommand{\RqTwoUtilDeiGateClusterRDhi}{0.075}
\newcommand{\RqTwoUtilBOneFail}{28}
\newcommand{\RqTwoUtilBOneN}{80}
\newcommand{\RqTwoUtilBOneRate}{65.0\%}
\newcommand{\RqTwoUtilBOneBlockRate}{0.350}
\newcommand{\RqTwoUtilBOneBlockLo}{0.255}
\newcommand{\RqTwoUtilBOneBlockHi}{0.459}
\newcommand{\RqTwoUtilBOneQuarEpisodes}{0}
\newcommand{\RqTwoUtilBOneQuarRate}{0.000}
\newcommand{\RqTwoUtilBOneQuarLo}{0.000}
\newcommand{\RqTwoUtilBOneQuarHi}{0.046}
\newcommand{\RqTwoUtilBOneFailDefense}{0}
\newcommand{\RqTwoUtilBOneFailModel}{28}
\newcommand{\RqTwoUtilBOneFailDegenerate}{0}
\newcommand{\RqTwoUtilBOneFisherP}{0.26}
\newcommand{\RqTwoUtilBOneMcNemarP}{0.057}
\newcommand{\RqTwoUtilBOneDiscordantB}{3}
\newcommand{\RqTwoUtilBOneDiscordantC}{11}
\newcommand{\RqTwoUtilBOnePairedN}{80}
\newcommand{\RqTwoUtilBOnePairedRD}{0.100}
\newcommand{\RqTwoUtilBOnePairedRDlo}{0.011}
\newcommand{\RqTwoUtilBOnePairedRDhi}{0.186}
\newcommand{\RqTwoUtilBOneClusterN}{20}
\newcommand{\RqTwoUtilBOneClusterRDlo}{0.013}
\newcommand{\RqTwoUtilBOneClusterRDhi}{0.188}
\newcommand{\RqTwoUtilBOneProtectedUseFail}{0}
\newcommand{\RqTwoUtilBOneProtectedUseN}{35}
\newcommand{\RqTwoUtilBOneProtectedUseBlockRate}{0.000}
\newcommand{\RqTwoUtilBOneProtectedUseBlockLo}{0.000}
\newcommand{\RqTwoUtilBOneProtectedUseBlockHi}{0.099}
\newcommand{\RqTwoUtilDeiProtectedUseFail}{0}
\newcommand{\RqTwoUtilDeiProtectedUseN}{29}
\newcommand{\RqTwoUtilDeiProtectedUseBlockRate}{0.000}
\newcommand{\RqTwoUtilDeiProtectedUseBlockLo}{0.000}
\newcommand{\RqTwoUtilDeiProtectedUseBlockHi}{0.117}
\newcommand{\RqTwoUtilDeiGateProtectedUseFail}{4}
\newcommand{\RqTwoUtilDeiGateProtectedUseN}{29}
\newcommand{\RqTwoUtilDeiGateProtectedUseBlockRate}{0.138}
\newcommand{\RqTwoUtilDeiGateProtectedUseBlockLo}{0.055}
\newcommand{\RqTwoUtilDeiGateProtectedUseBlockHi}{0.306}
\newcommand{\RqTwoUtilNoDefenseProtectedUseFail}{0}
\newcommand{\RqTwoUtilNoDefenseProtectedUseN}{28}
\newcommand{\RqTwoUtilNoDefenseProtectedUseBlockRate}{0.000}
\newcommand{\RqTwoUtilNoDefenseProtectedUseBlockLo}{0.000}
\newcommand{\RqTwoUtilNoDefenseProtectedUseBlockHi}{0.121}
\newcommand{\RqTwoUtilBOneNoProtectedUseFail}{28}
\newcommand{\RqTwoUtilBOneNoProtectedUseN}{45}
\newcommand{\RqTwoUtilBOneNoProtectedUseBlockRate}{0.622}
\newcommand{\RqTwoUtilBOneNoProtectedUseBlockLo}{0.476}
\newcommand{\RqTwoUtilBOneNoProtectedUseBlockHi}{0.749}
\newcommand{\RqTwoUtilDeiNoProtectedUseFail}{37}
\newcommand{\RqTwoUtilDeiNoProtectedUseN}{51}
\newcommand{\RqTwoUtilDeiNoProtectedUseBlockRate}{0.725}
\newcommand{\RqTwoUtilDeiNoProtectedUseBlockLo}{0.591}
\newcommand{\RqTwoUtilDeiNoProtectedUseBlockHi}{0.829}
\newcommand{\RqTwoUtilDeiGateNoProtectedUseFail}{35}
\newcommand{\RqTwoUtilDeiGateNoProtectedUseN}{51}
\newcommand{\RqTwoUtilDeiGateNoProtectedUseBlockRate}{0.686}
\newcommand{\RqTwoUtilDeiGateNoProtectedUseBlockLo}{0.550}
\newcommand{\RqTwoUtilDeiGateNoProtectedUseBlockHi}{0.797}
\newcommand{\RqTwoUtilNoDefenseNoProtectedUseFail}{36}
\newcommand{\RqTwoUtilNoDefenseNoProtectedUseN}{52}
\newcommand{\RqTwoUtilNoDefenseNoProtectedUseBlockRate}{0.692}
\newcommand{\RqTwoUtilNoDefenseNoProtectedUseBlockLo}{0.557}
\newcommand{\RqTwoUtilNoDefenseNoProtectedUseBlockHi}{0.801}
\newcommand{\RqTwoSweepTauThirtySecAdmit}{90}
\newcommand{\RqTwoSweepTauThirtySecN}{90}
\newcommand{\RqTwoSweepTauThirtySecRate}{1.000}
\newcommand{\RqTwoSweepTauThirtyBenignAdmit}{58}
\newcommand{\RqTwoSweepTauThirtyBenignN}{58}
\newcommand{\RqTwoSweepTauThirtyBenignRate}{1.000}
\newcommand{\RqTwoSweepTauFortyFiveSecAdmit}{90}
\newcommand{\RqTwoSweepTauFortyFiveSecN}{90}
\newcommand{\RqTwoSweepTauFortyFiveSecRate}{1.000}
\newcommand{\RqTwoSweepTauFortyFiveBenignAdmit}{58}
\newcommand{\RqTwoSweepTauFortyFiveBenignN}{58}
\newcommand{\RqTwoSweepTauFortyFiveBenignRate}{1.000}
\newcommand{\RqTwoSweepTauFiftySecAdmit}{90}
\newcommand{\RqTwoSweepTauFiftySecN}{90}
\newcommand{\RqTwoSweepTauFiftySecRate}{1.000}
\newcommand{\RqTwoSweepTauFiftyBenignAdmit}{58}
\newcommand{\RqTwoSweepTauFiftyBenignN}{58}
\newcommand{\RqTwoSweepTauFiftyBenignRate}{1.000}
\newcommand{\RqTwoSweepTauSixtySecAdmit}{0}
\newcommand{\RqTwoSweepTauSixtySecN}{90}
\newcommand{\RqTwoSweepTauSixtySecRate}{0.000}
\newcommand{\RqTwoSweepTauSixtyBenignAdmit}{0}
\newcommand{\RqTwoSweepTauSixtyBenignN}{58}
\newcommand{\RqTwoSweepTauSixtyBenignRate}{0.000}
\newcommand{\RqTwoSweepTauSixtySevenSecAdmit}{0}
\newcommand{\RqTwoSweepTauSixtySevenSecN}{90}
\newcommand{\RqTwoSweepTauSixtySevenSecRate}{0.000}
\newcommand{\RqTwoSweepTauSixtySevenBenignAdmit}{0}
\newcommand{\RqTwoSweepTauSixtySevenBenignN}{58}
\newcommand{\RqTwoSweepTauSixtySevenBenignRate}{0.000}
\newcommand{\RqTwoSweepEllAllTrustedSecAdmit}{90}
\newcommand{\RqTwoSweepEllAllTrustedSecN}{90}
\newcommand{\RqTwoSweepEllAllTrustedSecRate}{1.000}
\newcommand{\RqTwoSweepEllAllTrustedBenignAdmit}{58}
\newcommand{\RqTwoSweepEllAllTrustedBenignN}{58}
\newcommand{\RqTwoSweepEllAllTrustedBenignRate}{1.000}
\newcommand{\RqTwoProbeFOneWrongFileColocatedNdCount}{0}
\newcommand{\RqTwoProbeFOneWrongFileColocatedNdN}{20}
\newcommand{\RqTwoProbeFOneWrongFileColocatedNdRate}{0.000}
\newcommand{\RqTwoProbeFOneWrongFileColocatedDeiCount}{1}
\newcommand{\RqTwoProbeFOneWrongFileColocatedDeiN}{20}
\newcommand{\RqTwoProbeFOneWrongFileColocatedDeiRate}{0.050}
\newcommand{\RqTwoProbeFOneWrongFileColocatedBOneCount}{1}
\newcommand{\RqTwoProbeFOneWrongFileColocatedBOneN}{20}
\newcommand{\RqTwoProbeFOneWrongFileColocatedBOneRate}{0.050}
\newcommand{\RqTwoProbeFOneWrongFileColocatedAnyQuar}{5}
\newcommand{\RqTwoProbeFOneWrongFileColocatedAnyQuarN}{5}
\newcommand{\RqTwoProbeFOneWrongFileColocatedAnyQuarRate}{1.000}
\newcommand{\RqTwoProbeFOneWrongFileColocatedAnyQuarLo}{0.566}
\newcommand{\RqTwoProbeFOneWrongFileColocatedAnyQuarHi}{1.000}
\newcommand{\RqTwoProbeFOneWrongFileColocatedAnyFq}{1}
\newcommand{\RqTwoProbeFOneWrongFileColocatedAnyFqN}{1}
\newcommand{\RqTwoProbeFOneWrongFileColocatedAnyFqRate}{1.000}
\newcommand{\RqTwoProbeFOneWrongFileColocatedAnyFqLo}{0.207}
\newcommand{\RqTwoProbeFOneWrongFileColocatedAnyFqHi}{1.000}
\newcommand{\RqTwoProbeFOneWrongFileColocatedDeltaTwoQuar}{5}
\newcommand{\RqTwoProbeFOneWrongFileColocatedDeltaTwoQuarN}{5}
\newcommand{\RqTwoProbeFOneWrongFileColocatedDeltaTwoQuarRate}{1.000}
\newcommand{\RqTwoProbeFOneWrongFileColocatedDeltaTwoQuarLo}{0.566}
\newcommand{\RqTwoProbeFOneWrongFileColocatedDeltaTwoQuarHi}{1.000}
\newcommand{\RqTwoProbeFOneWrongFileColocatedDeltaTwoFq}{1}
\newcommand{\RqTwoProbeFOneWrongFileColocatedDeltaTwoFqN}{1}
\newcommand{\RqTwoProbeFOneWrongFileColocatedDeltaTwoFqRate}{1.000}
\newcommand{\RqTwoProbeFOneWrongFileColocatedDeltaTwoFqLo}{0.207}
\newcommand{\RqTwoProbeFOneWrongFileColocatedDeltaTwoFqHi}{1.000}
\newcommand{\RqTwoProbeFOneWrongFileColocatedDeltaThreeQuar}{4}
\newcommand{\RqTwoProbeFOneWrongFileColocatedDeltaThreeQuarN}{5}
\newcommand{\RqTwoProbeFOneWrongFileColocatedDeltaThreeQuarRate}{0.800}
\newcommand{\RqTwoProbeFOneWrongFileColocatedDeltaThreeQuarLo}{0.376}
\newcommand{\RqTwoProbeFOneWrongFileColocatedDeltaThreeQuarHi}{0.964}
\newcommand{\RqTwoProbeFOneWrongFileColocatedDeltaThreeFq}{1}
\newcommand{\RqTwoProbeFOneWrongFileColocatedDeltaThreeFqN}{1}
\newcommand{\RqTwoProbeFOneWrongFileColocatedDeltaThreeFqRate}{1.000}
\newcommand{\RqTwoProbeFOneWrongFileColocatedDeltaThreeFqLo}{0.207}
\newcommand{\RqTwoProbeFOneWrongFileColocatedDeltaThreeFqHi}{1.000}
\newcommand{\RqTwoProbeFOneWrongFileColocatedDeltaOneQuar}{0}
\newcommand{\RqTwoProbeFOneWrongFileColocatedDeltaOneQuarN}{0}
\newcommand{\RqTwoProbeFOneWrongFileColocatedDeltaOneQuarRate}{--}
\newcommand{\RqTwoProbeFOneWrongFileColocatedDeltaOneQuarLo}{--}
\newcommand{\RqTwoProbeFOneWrongFileColocatedDeltaOneQuarHi}{--}
\newcommand{\RqTwoProbeFOneWrongFileColocatedDeltaOneFq}{0}
\newcommand{\RqTwoProbeFOneWrongFileColocatedDeltaOneFqN}{0}
\newcommand{\RqTwoProbeFOneWrongFileColocatedDeltaOneFqRate}{--}
\newcommand{\RqTwoProbeFOneWrongFileColocatedDeltaOneFqLo}{--}
\newcommand{\RqTwoProbeFOneWrongFileColocatedDeltaOneFqHi}{--}
\newcommand{\RqTwoProbeFOneWrongFileSeparateNdCount}{2}
\newcommand{\RqTwoProbeFOneWrongFileSeparateNdN}{20}
\newcommand{\RqTwoProbeFOneWrongFileSeparateNdRate}{0.100}
\newcommand{\RqTwoProbeFOneWrongFileSeparateDeiCount}{1}
\newcommand{\RqTwoProbeFOneWrongFileSeparateDeiN}{20}
\newcommand{\RqTwoProbeFOneWrongFileSeparateDeiRate}{0.050}
\newcommand{\RqTwoProbeFOneWrongFileSeparateBOneCount}{2}
\newcommand{\RqTwoProbeFOneWrongFileSeparateBOneN}{20}
\newcommand{\RqTwoProbeFOneWrongFileSeparateBOneRate}{0.100}
\newcommand{\RqTwoProbeFOneWrongFileSeparateAnyQuar}{12}
\newcommand{\RqTwoProbeFOneWrongFileSeparateAnyQuarN}{12}
\newcommand{\RqTwoProbeFOneWrongFileSeparateAnyQuarRate}{1.000}
\newcommand{\RqTwoProbeFOneWrongFileSeparateAnyQuarLo}{0.757}
\newcommand{\RqTwoProbeFOneWrongFileSeparateAnyQuarHi}{1.000}
\newcommand{\RqTwoProbeFOneWrongFileSeparateAnyFq}{7}
\newcommand{\RqTwoProbeFOneWrongFileSeparateAnyFqN}{8}
\newcommand{\RqTwoProbeFOneWrongFileSeparateAnyFqRate}{0.875}
\newcommand{\RqTwoProbeFOneWrongFileSeparateAnyFqLo}{0.529}
\newcommand{\RqTwoProbeFOneWrongFileSeparateAnyFqHi}{0.978}
\newcommand{\RqTwoProbeFOneWrongFileSeparateDeltaTwoQuar}{11}
\newcommand{\RqTwoProbeFOneWrongFileSeparateDeltaTwoQuarN}{12}
\newcommand{\RqTwoProbeFOneWrongFileSeparateDeltaTwoQuarRate}{0.917}
\newcommand{\RqTwoProbeFOneWrongFileSeparateDeltaTwoQuarLo}{0.646}
\newcommand{\RqTwoProbeFOneWrongFileSeparateDeltaTwoQuarHi}{0.985}
\newcommand{\RqTwoProbeFOneWrongFileSeparateDeltaTwoFq}{7}
\newcommand{\RqTwoProbeFOneWrongFileSeparateDeltaTwoFqN}{8}
\newcommand{\RqTwoProbeFOneWrongFileSeparateDeltaTwoFqRate}{0.875}
\newcommand{\RqTwoProbeFOneWrongFileSeparateDeltaTwoFqLo}{0.529}
\newcommand{\RqTwoProbeFOneWrongFileSeparateDeltaTwoFqHi}{0.978}
\newcommand{\RqTwoProbeFOneWrongFileSeparateDeltaThreeQuar}{9}
\newcommand{\RqTwoProbeFOneWrongFileSeparateDeltaThreeQuarN}{12}
\newcommand{\RqTwoProbeFOneWrongFileSeparateDeltaThreeQuarRate}{0.750}
\newcommand{\RqTwoProbeFOneWrongFileSeparateDeltaThreeQuarLo}{0.468}
\newcommand{\RqTwoProbeFOneWrongFileSeparateDeltaThreeQuarHi}{0.911}
\newcommand{\RqTwoProbeFOneWrongFileSeparateDeltaThreeFq}{3}
\newcommand{\RqTwoProbeFOneWrongFileSeparateDeltaThreeFqN}{8}
\newcommand{\RqTwoProbeFOneWrongFileSeparateDeltaThreeFqRate}{0.375}
\newcommand{\RqTwoProbeFOneWrongFileSeparateDeltaThreeFqLo}{0.137}
\newcommand{\RqTwoProbeFOneWrongFileSeparateDeltaThreeFqHi}{0.694}
\newcommand{\RqTwoProbeFOneWrongFileSeparateDeltaOneQuar}{0}
\newcommand{\RqTwoProbeFOneWrongFileSeparateDeltaOneQuarN}{0}
\newcommand{\RqTwoProbeFOneWrongFileSeparateDeltaOneQuarRate}{--}
\newcommand{\RqTwoProbeFOneWrongFileSeparateDeltaOneQuarLo}{--}
\newcommand{\RqTwoProbeFOneWrongFileSeparateDeltaOneQuarHi}{--}
\newcommand{\RqTwoProbeFOneWrongFileSeparateDeltaOneFq}{0}
\newcommand{\RqTwoProbeFOneWrongFileSeparateDeltaOneFqN}{0}
\newcommand{\RqTwoProbeFOneWrongFileSeparateDeltaOneFqRate}{--}
\newcommand{\RqTwoProbeFOneWrongFileSeparateDeltaOneFqLo}{--}
\newcommand{\RqTwoProbeFOneWrongFileSeparateDeltaOneFqHi}{--}
\newcommand{\RqTwoProbeFTwoPrematureSendColocatedNdCount}{12}
\newcommand{\RqTwoProbeFTwoPrematureSendColocatedNdN}{20}
\newcommand{\RqTwoProbeFTwoPrematureSendColocatedNdRate}{0.600}
\newcommand{\RqTwoProbeFTwoPrematureSendColocatedDeiCount}{14}
\newcommand{\RqTwoProbeFTwoPrematureSendColocatedDeiN}{20}
\newcommand{\RqTwoProbeFTwoPrematureSendColocatedDeiRate}{0.700}
\newcommand{\RqTwoProbeFTwoPrematureSendColocatedBOneCount}{15}
\newcommand{\RqTwoProbeFTwoPrematureSendColocatedBOneN}{20}
\newcommand{\RqTwoProbeFTwoPrematureSendColocatedBOneRate}{0.750}
\newcommand{\RqTwoProbeFTwoPrematureSendColocatedAnyQuar}{14}
\newcommand{\RqTwoProbeFTwoPrematureSendColocatedAnyQuarN}{14}
\newcommand{\RqTwoProbeFTwoPrematureSendColocatedAnyQuarRate}{1.000}
\newcommand{\RqTwoProbeFTwoPrematureSendColocatedAnyQuarLo}{0.785}
\newcommand{\RqTwoProbeFTwoPrematureSendColocatedAnyQuarHi}{1.000}
\newcommand{\RqTwoProbeFTwoPrematureSendColocatedAnyFq}{0}
\newcommand{\RqTwoProbeFTwoPrematureSendColocatedAnyFqN}{0}
\newcommand{\RqTwoProbeFTwoPrematureSendColocatedAnyFqRate}{--}
\newcommand{\RqTwoProbeFTwoPrematureSendColocatedAnyFqLo}{--}
\newcommand{\RqTwoProbeFTwoPrematureSendColocatedAnyFqHi}{--}
\newcommand{\RqTwoProbeFTwoPrematureSendColocatedDeltaTwoQuar}{0}
\newcommand{\RqTwoProbeFTwoPrematureSendColocatedDeltaTwoQuarN}{14}
\newcommand{\RqTwoProbeFTwoPrematureSendColocatedDeltaTwoQuarRate}{0.000}
\newcommand{\RqTwoProbeFTwoPrematureSendColocatedDeltaTwoQuarLo}{0.000}
\newcommand{\RqTwoProbeFTwoPrematureSendColocatedDeltaTwoQuarHi}{0.215}
\newcommand{\RqTwoProbeFTwoPrematureSendColocatedDeltaTwoFq}{0}
\newcommand{\RqTwoProbeFTwoPrematureSendColocatedDeltaTwoFqN}{0}
\newcommand{\RqTwoProbeFTwoPrematureSendColocatedDeltaTwoFqRate}{--}
\newcommand{\RqTwoProbeFTwoPrematureSendColocatedDeltaTwoFqLo}{--}
\newcommand{\RqTwoProbeFTwoPrematureSendColocatedDeltaTwoFqHi}{--}
\newcommand{\RqTwoProbeFTwoPrematureSendColocatedDeltaThreeQuar}{14}
\newcommand{\RqTwoProbeFTwoPrematureSendColocatedDeltaThreeQuarN}{14}
\newcommand{\RqTwoProbeFTwoPrematureSendColocatedDeltaThreeQuarRate}{1.000}
\newcommand{\RqTwoProbeFTwoPrematureSendColocatedDeltaThreeQuarLo}{0.785}
\newcommand{\RqTwoProbeFTwoPrematureSendColocatedDeltaThreeQuarHi}{1.000}
\newcommand{\RqTwoProbeFTwoPrematureSendColocatedDeltaThreeFq}{0}
\newcommand{\RqTwoProbeFTwoPrematureSendColocatedDeltaThreeFqN}{0}
\newcommand{\RqTwoProbeFTwoPrematureSendColocatedDeltaThreeFqRate}{--}
\newcommand{\RqTwoProbeFTwoPrematureSendColocatedDeltaThreeFqLo}{--}
\newcommand{\RqTwoProbeFTwoPrematureSendColocatedDeltaThreeFqHi}{--}
\newcommand{\RqTwoProbeFTwoPrematureSendColocatedDeltaOneQuar}{0}
\newcommand{\RqTwoProbeFTwoPrematureSendColocatedDeltaOneQuarN}{0}
\newcommand{\RqTwoProbeFTwoPrematureSendColocatedDeltaOneQuarRate}{--}
\newcommand{\RqTwoProbeFTwoPrematureSendColocatedDeltaOneQuarLo}{--}
\newcommand{\RqTwoProbeFTwoPrematureSendColocatedDeltaOneQuarHi}{--}
\newcommand{\RqTwoProbeFTwoPrematureSendColocatedDeltaOneFq}{0}
\newcommand{\RqTwoProbeFTwoPrematureSendColocatedDeltaOneFqN}{0}
\newcommand{\RqTwoProbeFTwoPrematureSendColocatedDeltaOneFqRate}{--}
\newcommand{\RqTwoProbeFTwoPrematureSendColocatedDeltaOneFqLo}{--}
\newcommand{\RqTwoProbeFTwoPrematureSendColocatedDeltaOneFqHi}{--}
\newcommand{\RqTwoProbeFThreeWrongOptionColocatedNdCount}{14}
\newcommand{\RqTwoProbeFThreeWrongOptionColocatedNdN}{20}
\newcommand{\RqTwoProbeFThreeWrongOptionColocatedNdRate}{0.700}
\newcommand{\RqTwoProbeFThreeWrongOptionColocatedDeiCount}{16}
\newcommand{\RqTwoProbeFThreeWrongOptionColocatedDeiN}{20}
\newcommand{\RqTwoProbeFThreeWrongOptionColocatedDeiRate}{0.800}
\newcommand{\RqTwoProbeFThreeWrongOptionColocatedBOneCount}{14}
\newcommand{\RqTwoProbeFThreeWrongOptionColocatedBOneN}{20}
\newcommand{\RqTwoProbeFThreeWrongOptionColocatedBOneRate}{0.700}
\newcommand{\RqTwoProbeFThreeWrongOptionColocatedAnyQuar}{16}
\newcommand{\RqTwoProbeFThreeWrongOptionColocatedAnyQuarN}{16}
\newcommand{\RqTwoProbeFThreeWrongOptionColocatedAnyQuarRate}{1.000}
\newcommand{\RqTwoProbeFThreeWrongOptionColocatedAnyQuarLo}{0.806}
\newcommand{\RqTwoProbeFThreeWrongOptionColocatedAnyQuarHi}{1.000}
\newcommand{\RqTwoProbeFThreeWrongOptionColocatedAnyFq}{20}
\newcommand{\RqTwoProbeFThreeWrongOptionColocatedAnyFqN}{20}
\newcommand{\RqTwoProbeFThreeWrongOptionColocatedAnyFqRate}{1.000}
\newcommand{\RqTwoProbeFThreeWrongOptionColocatedAnyFqLo}{0.839}
\newcommand{\RqTwoProbeFThreeWrongOptionColocatedAnyFqHi}{1.000}
\newcommand{\RqTwoProbeFThreeWrongOptionColocatedDeltaTwoQuar}{16}
\newcommand{\RqTwoProbeFThreeWrongOptionColocatedDeltaTwoQuarN}{16}
\newcommand{\RqTwoProbeFThreeWrongOptionColocatedDeltaTwoQuarRate}{1.000}
\newcommand{\RqTwoProbeFThreeWrongOptionColocatedDeltaTwoQuarLo}{0.806}
\newcommand{\RqTwoProbeFThreeWrongOptionColocatedDeltaTwoQuarHi}{1.000}
\newcommand{\RqTwoProbeFThreeWrongOptionColocatedDeltaTwoFq}{20}
\newcommand{\RqTwoProbeFThreeWrongOptionColocatedDeltaTwoFqN}{20}
\newcommand{\RqTwoProbeFThreeWrongOptionColocatedDeltaTwoFqRate}{1.000}
\newcommand{\RqTwoProbeFThreeWrongOptionColocatedDeltaTwoFqLo}{0.839}
\newcommand{\RqTwoProbeFThreeWrongOptionColocatedDeltaTwoFqHi}{1.000}
\newcommand{\RqTwoProbeFThreeWrongOptionColocatedDeltaThreeQuar}{0}
\newcommand{\RqTwoProbeFThreeWrongOptionColocatedDeltaThreeQuarN}{16}
\newcommand{\RqTwoProbeFThreeWrongOptionColocatedDeltaThreeQuarRate}{0.000}
\newcommand{\RqTwoProbeFThreeWrongOptionColocatedDeltaThreeQuarLo}{0.000}
\newcommand{\RqTwoProbeFThreeWrongOptionColocatedDeltaThreeQuarHi}{0.194}
\newcommand{\RqTwoProbeFThreeWrongOptionColocatedDeltaThreeFq}{0}
\newcommand{\RqTwoProbeFThreeWrongOptionColocatedDeltaThreeFqN}{20}
\newcommand{\RqTwoProbeFThreeWrongOptionColocatedDeltaThreeFqRate}{0.000}
\newcommand{\RqTwoProbeFThreeWrongOptionColocatedDeltaThreeFqLo}{0.000}
\newcommand{\RqTwoProbeFThreeWrongOptionColocatedDeltaThreeFqHi}{0.161}
\newcommand{\RqTwoProbeFThreeWrongOptionColocatedDeltaOneQuar}{0}
\newcommand{\RqTwoProbeFThreeWrongOptionColocatedDeltaOneQuarN}{0}
\newcommand{\RqTwoProbeFThreeWrongOptionColocatedDeltaOneQuarRate}{--}
\newcommand{\RqTwoProbeFThreeWrongOptionColocatedDeltaOneQuarLo}{--}
\newcommand{\RqTwoProbeFThreeWrongOptionColocatedDeltaOneQuarHi}{--}
\newcommand{\RqTwoProbeFThreeWrongOptionColocatedDeltaOneFq}{0}
\newcommand{\RqTwoProbeFThreeWrongOptionColocatedDeltaOneFqN}{0}
\newcommand{\RqTwoProbeFThreeWrongOptionColocatedDeltaOneFqRate}{--}
\newcommand{\RqTwoProbeFThreeWrongOptionColocatedDeltaOneFqLo}{--}
\newcommand{\RqTwoProbeFThreeWrongOptionColocatedDeltaOneFqHi}{--}
\newcommand{\RqTwoSeparableDeltaThreeK}{23}
\newcommand{\RqTwoSeparableDeltaThreeN}{26}
\newcommand{\RqTwoSeparableDeltaThreeProbes}{9}
\newcommand{\RqTwoSeparableDeltaThreeRate}{0.885}
\newcommand{\RqTwoSeparableDeltaThreeLo}{0.710}
\newcommand{\RqTwoSeparableDeltaThreeHi}{0.960}
\newcommand{\RqTwoSeparableDeltaThreeClusterLo}{0.704}
\newcommand{\RqTwoSeparableDeltaThreeClusterHi}{1.000}
\newcommand{\RqTwoColocatedDeltaThreeK}{4}
\newcommand{\RqTwoColocatedDeltaThreeN}{21}
\newcommand{\RqTwoColocatedDeltaThreeProbes}{9}
\newcommand{\RqTwoColocatedDeltaThreeRate}{0.190}
\newcommand{\RqTwoColocatedDeltaThreeLo}{0.077}
\newcommand{\RqTwoColocatedDeltaThreeHi}{0.400}
\newcommand{\RqTwoColocatedDeltaThreeClusterLo}{0.034}
\newcommand{\RqTwoColocatedDeltaThreeClusterHi}{0.500}
\newcommand{\RqTwoHonestDeltaThreeK}{4}
\newcommand{\RqTwoHonestDeltaThreeN}{29}
\newcommand{\RqTwoHonestDeltaThreeProbes}{10}
\newcommand{\RqTwoHonestDeltaThreeRate}{0.138}
\newcommand{\RqTwoHonestDeltaThreeLo}{0.055}
\newcommand{\RqTwoHonestDeltaThreeHi}{0.306}
\newcommand{\RqTwoHonestDeltaThreeClusterLo}{0.028}
\newcommand{\RqTwoHonestDeltaThreeClusterHi}{0.333}
\newcommand{\RqTwoColocatedDeltaTwoK}{21}
\newcommand{\RqTwoColocatedDeltaTwoN}{21}
\newcommand{\RqTwoColocatedDeltaTwoProbes}{9}
\newcommand{\RqTwoColocatedDeltaTwoRate}{1.000}
\newcommand{\RqTwoColocatedDeltaTwoLo}{0.845}
\newcommand{\RqTwoColocatedDeltaTwoHi}{1.000}
\newcommand{\RqTwoColocatedDeltaTwoClusterLo}{1.000}
\newcommand{\RqTwoColocatedDeltaTwoClusterHi}{1.000}
\newcommand{\RqTwoHonestDeltaTwoK}{28}
\newcommand{\RqTwoHonestDeltaTwoN}{29}
\newcommand{\RqTwoHonestDeltaTwoProbes}{10}
\newcommand{\RqTwoHonestDeltaTwoRate}{0.966}
\newcommand{\RqTwoHonestDeltaTwoLo}{0.828}
\newcommand{\RqTwoHonestDeltaTwoHi}{0.994}
\newcommand{\RqTwoHonestDeltaTwoClusterLo}{0.880}
\newcommand{\RqTwoHonestDeltaTwoClusterHi}{1.000}
\newcommand{\RqTwoDeltaTwoTp}{15}
\newcommand{\RqTwoDeltaTwoFn}{12}
\newcommand{\RqTwoDeltaTwoFp}{6}
\newcommand{\RqTwoDeltaTwoTn}{1}
\newcommand{\RqTwoDeltaTwoMissRate}{0.444}
\newcommand{\RqTwoDeltaTwoMissLo}{0.276}
\newcommand{\RqTwoDeltaTwoMissHi}{0.627}
\newcommand{\RqTwoDeltaTwoFpRate}{0.857}
\newcommand{\RqTwoDeltaTwoFpLo}{0.487}
\newcommand{\RqTwoDeltaTwoFpHi}{0.974}
\newcommand{\RqTwoDeltaThreeTp}{12}
\newcommand{\RqTwoDeltaThreeFn}{15}
\newcommand{\RqTwoDeltaThreeFp}{5}
\newcommand{\RqTwoDeltaThreeTn}{2}
\newcommand{\RqTwoDeltaThreeMissRate}{0.556}
\newcommand{\RqTwoDeltaThreeMissLo}{0.373}
\newcommand{\RqTwoDeltaThreeMissHi}{0.724}
\newcommand{\RqTwoDeltaThreeFpRate}{0.714}
\newcommand{\RqTwoDeltaThreeFpLo}{0.359}
\newcommand{\RqTwoDeltaThreeFpHi}{0.918}
  % registered probe and oracle campaign; regenerate via analyze.py
% AUTO-GENERATED by build/emit_rq1_table_macros.py from RQ1 JSONL.
\newcommand{\RqOneNodefenseCell}{0.527 (79/150)}
\newcommand{\RqOneNodefenseCI}{[0.447, 0.605]}
\newcommand{\RqOneNodefenseFiltered}{0}
\newcommand{\RqOneDeiCell}{0.087 (13/150)}
\newcommand{\RqOneDeiCI}{[0.051, 0.143]}
\newcommand{\RqOneDeiFiltered}{114}
\newcommand{\RqOneDeigateCell}{0.087 (13/150)}
\newcommand{\RqOneDeigateCI}{[0.051, 0.143]}
\newcommand{\RqOneDeigateFiltered}{110}
\newcommand{\RqOneBOneCell}{0.087 (13/150)}
\newcommand{\RqOneBOneCI}{[0.051, 0.143]}
\newcommand{\RqOneBOneFiltered}{134}
\newcommand{\RqOneBTwoCell}{0.220 (33/150)}
\newcommand{\RqOneBTwoCI}{[0.161, 0.293]}
\newcommand{\RqOneBTwoFiltered}{124}
\newcommand{\RqOneBTwoStrictCell}{0.047 (7/150)}
\newcommand{\RqOneBTwoStrictCI}{[0.023, 0.093]}
\newcommand{\RqOneBTwoStrictFiltered}{146}
\newcommand{\RqOneBThreeCell}{0.167 (25/150)}
\newcommand{\RqOneBThreeCI}{[0.116, 0.234]}
\newcommand{\RqOneBThreeFiltered}{105}
  % Table 4 cells; regenerate via build/emit_rq1_table_macros.py
% AUTO-GENERATED by analyze.py from out/rq3_all.jsonl; primary (degenerate retained); do not edit.
% Voided: 50; marked degenerate in source: 1; bootstrap seed: 20260814.
\newcommand{\RqThreeEpisodes}{1164}
\newcommand{\RqThreeVoided}{50}
\newcommand{\RqThreeDegenerate}{1}
\newcommand{\RqThreeAnalysisMode}{primary}
\newcommand{\RqThreeUtilBTwoStrictFail}{220}
\newcommand{\RqThreeUtilBTwoStrictN}{291}
\newcommand{\RqThreeUtilBTwoStrictRate}{24.4\%}
\newcommand{\RqThreeUtilBTwoStrictBlockRate}{0.756}
\newcommand{\RqThreeUtilBTwoStrictBlockLo}{0.704}
\newcommand{\RqThreeUtilBTwoStrictBlockHi}{0.802}
\newcommand{\RqThreeUtilBTwoStrictQuarEpisodes}{0}
\newcommand{\RqThreeUtilBTwoStrictQuarRate}{0.000}
\newcommand{\RqThreeUtilBTwoStrictQuarLo}{0.000}
\newcommand{\RqThreeUtilBTwoStrictQuarHi}{0.013}
\newcommand{\RqThreeUtilBTwoStrictFailDefense}{132}
\newcommand{\RqThreeUtilBTwoStrictFailModel}{88}
\newcommand{\RqThreeUtilBTwoStrictFailDegenerate}{0}
\newcommand{\RqThreeUtilBTwoStrictFisherP}{0.00}
\newcommand{\RqThreeUtilBTwoStrictMcNemarP}{0.000}
\newcommand{\RqThreeUtilBTwoStrictDiscordantB}{120}
\newcommand{\RqThreeUtilBTwoStrictDiscordantC}{11}
\newcommand{\RqThreeUtilBTwoStrictPairedN}{291}
\newcommand{\RqThreeUtilBTwoStrictPairedRD}{-0.375}
\newcommand{\RqThreeUtilBTwoStrictPairedRDlo}{-0.435}
\newcommand{\RqThreeUtilBTwoStrictPairedRDhi}{-0.308}
\newcommand{\RqThreeUtilBTwoStrictClusterN}{97}
\newcommand{\RqThreeUtilBTwoStrictClusterRDlo}{-0.467}
\newcommand{\RqThreeUtilBTwoStrictClusterRDhi}{-0.282}
\newcommand{\RqThreeUtilNoDefenseFail}{111}
\newcommand{\RqThreeUtilNoDefenseN}{291}
\newcommand{\RqThreeUtilNoDefenseRate}{61.9\%}
\newcommand{\RqThreeUtilNoDefenseBlockRate}{0.381}
\newcommand{\RqThreeUtilNoDefenseBlockLo}{0.328}
\newcommand{\RqThreeUtilNoDefenseBlockHi}{0.438}
\newcommand{\RqThreeUtilNoDefenseQuarEpisodes}{0}
\newcommand{\RqThreeUtilNoDefenseQuarRate}{0.000}
\newcommand{\RqThreeUtilNoDefenseQuarLo}{0.000}
\newcommand{\RqThreeUtilNoDefenseQuarHi}{0.013}
\newcommand{\RqThreeUtilNoDefenseFailDefense}{0}
\newcommand{\RqThreeUtilNoDefenseFailModel}{110}
\newcommand{\RqThreeUtilNoDefenseFailDegenerate}{1}
\newcommand{\RqThreeUtilDeiGateFail}{189}
\newcommand{\RqThreeUtilDeiGateN}{291}
\newcommand{\RqThreeUtilDeiGateRate}{35.1\%}
\newcommand{\RqThreeUtilDeiGateBlockRate}{0.649}
\newcommand{\RqThreeUtilDeiGateBlockLo}{0.593}
\newcommand{\RqThreeUtilDeiGateBlockHi}{0.702}
\newcommand{\RqThreeUtilDeiGateQuarEpisodes}{15}
\newcommand{\RqThreeUtilDeiGateQuarRate}{0.052}
\newcommand{\RqThreeUtilDeiGateQuarLo}{0.031}
\newcommand{\RqThreeUtilDeiGateQuarHi}{0.083}
\newcommand{\RqThreeUtilDeiGateFailDefense}{74}
\newcommand{\RqThreeUtilDeiGateFailModel}{115}
\newcommand{\RqThreeUtilDeiGateFailDegenerate}{0}
\newcommand{\RqThreeUtilDeiGateFisherP}{0.00}
\newcommand{\RqThreeUtilDeiGateMcNemarP}{0.000}
\newcommand{\RqThreeUtilDeiGateDiscordantB}{99}
\newcommand{\RqThreeUtilDeiGateDiscordantC}{21}
\newcommand{\RqThreeUtilDeiGatePairedN}{291}
\newcommand{\RqThreeUtilDeiGatePairedRD}{-0.268}
\newcommand{\RqThreeUtilDeiGatePairedRDlo}{-0.333}
\newcommand{\RqThreeUtilDeiGatePairedRDhi}{-0.199}
\newcommand{\RqThreeUtilDeiGateClusterN}{97}
\newcommand{\RqThreeUtilDeiGateClusterRDlo}{-0.357}
\newcommand{\RqThreeUtilDeiGateClusterRDhi}{-0.182}
\newcommand{\RqThreeUtilBOneFail}{205}
\newcommand{\RqThreeUtilBOneN}{291}
\newcommand{\RqThreeUtilBOneRate}{29.6\%}
\newcommand{\RqThreeUtilBOneBlockRate}{0.704}
\newcommand{\RqThreeUtilBOneBlockLo}{0.650}
\newcommand{\RqThreeUtilBOneBlockHi}{0.754}
\newcommand{\RqThreeUtilBOneQuarEpisodes}{0}
\newcommand{\RqThreeUtilBOneQuarRate}{0.000}
\newcommand{\RqThreeUtilBOneQuarLo}{0.000}
\newcommand{\RqThreeUtilBOneQuarHi}{0.013}
\newcommand{\RqThreeUtilBOneFailDefense}{89}
\newcommand{\RqThreeUtilBOneFailModel}{116}
\newcommand{\RqThreeUtilBOneFailDegenerate}{0}
\newcommand{\RqThreeUtilBOneFisherP}{0.00}
\newcommand{\RqThreeUtilBOneMcNemarP}{0.000}
\newcommand{\RqThreeUtilBOneDiscordantB}{104}
\newcommand{\RqThreeUtilBOneDiscordantC}{10}
\newcommand{\RqThreeUtilBOnePairedN}{291}
\newcommand{\RqThreeUtilBOnePairedRD}{-0.323}
\newcommand{\RqThreeUtilBOnePairedRDlo}{-0.382}
\newcommand{\RqThreeUtilBOnePairedRDhi}{-0.259}
\newcommand{\RqThreeUtilBOneClusterN}{97}
\newcommand{\RqThreeUtilBOneClusterRDlo}{-0.409}
\newcommand{\RqThreeUtilBOneClusterRDhi}{-0.237}
\newcommand{\RqThreeUtilDeiGateVsBTwoStrictPairedN}{291}
\newcommand{\RqThreeUtilDeiGateVsBTwoStrictPairedRD}{0.107}
\newcommand{\RqThreeUtilDeiGateVsBTwoStrictPairedRDlo}{0.057}
\newcommand{\RqThreeUtilDeiGateVsBTwoStrictPairedRDhi}{0.156}
\newcommand{\RqThreeUtilDeiGateVsBTwoStrictMcNemarP}{0.000}
\newcommand{\RqThreeUtilDeiGateVsBTwoStrictHolmP}{0.000}
\newcommand{\RqThreeUtilDeiGateVsBTwoStrictClusterN}{97}
\newcommand{\RqThreeUtilDeiGateVsBTwoStrictClusterRDlo}{0.045}
\newcommand{\RqThreeUtilDeiGateVsBTwoStrictClusterRDhi}{0.172}
\newcommand{\RqThreeUtilDeiGateVsBOnePairedN}{291}
\newcommand{\RqThreeUtilDeiGateVsBOnePairedRD}{0.055}
\newcommand{\RqThreeUtilDeiGateVsBOnePairedRDlo}{0.003}
\newcommand{\RqThreeUtilDeiGateVsBOnePairedRDhi}{0.106}
\newcommand{\RqThreeUtilDeiGateVsBOneMcNemarP}{0.052}
\newcommand{\RqThreeUtilDeiGateVsBOneHolmP}{0.052}
\newcommand{\RqThreeUtilDeiGateVsBOneClusterN}{97}
\newcommand{\RqThreeUtilDeiGateVsBOneClusterRDlo}{-0.007}
\newcommand{\RqThreeUtilDeiGateVsBOneClusterRDhi}{0.117}
\newcommand{\RqThreeUtilBOneProtectedUseFail}{97}
\newcommand{\RqThreeUtilBOneProtectedUseN}{119}
\newcommand{\RqThreeUtilBOneProtectedUseBlockRate}{0.815}
\newcommand{\RqThreeUtilBOneProtectedUseBlockLo}{0.736}
\newcommand{\RqThreeUtilBOneProtectedUseBlockHi}{0.875}
\newcommand{\RqThreeUtilBTwoStrictProtectedUseFail}{132}
\newcommand{\RqThreeUtilBTwoStrictProtectedUseN}{133}
\newcommand{\RqThreeUtilBTwoStrictProtectedUseBlockRate}{0.992}
\newcommand{\RqThreeUtilBTwoStrictProtectedUseBlockLo}{0.959}
\newcommand{\RqThreeUtilBTwoStrictProtectedUseBlockHi}{0.999}
\newcommand{\RqThreeUtilDeiGateProtectedUseFail}{86}
\newcommand{\RqThreeUtilDeiGateProtectedUseN}{120}
\newcommand{\RqThreeUtilDeiGateProtectedUseBlockRate}{0.717}
\newcommand{\RqThreeUtilDeiGateProtectedUseBlockLo}{0.630}
\newcommand{\RqThreeUtilDeiGateProtectedUseBlockHi}{0.790}
\newcommand{\RqThreeUtilNoDefenseProtectedUseFail}{19}
\newcommand{\RqThreeUtilNoDefenseProtectedUseN}{128}
\newcommand{\RqThreeUtilNoDefenseProtectedUseBlockRate}{0.148}
\newcommand{\RqThreeUtilNoDefenseProtectedUseBlockLo}{0.097}
\newcommand{\RqThreeUtilNoDefenseProtectedUseBlockHi}{0.220}
\newcommand{\RqThreeUtilBOneNoProtectedUseFail}{108}
\newcommand{\RqThreeUtilBOneNoProtectedUseN}{172}
\newcommand{\RqThreeUtilBOneNoProtectedUseBlockRate}{0.628}
\newcommand{\RqThreeUtilBOneNoProtectedUseBlockLo}{0.554}
\newcommand{\RqThreeUtilBOneNoProtectedUseBlockHi}{0.697}
\newcommand{\RqThreeUtilBTwoStrictNoProtectedUseFail}{88}
\newcommand{\RqThreeUtilBTwoStrictNoProtectedUseN}{158}
\newcommand{\RqThreeUtilBTwoStrictNoProtectedUseBlockRate}{0.557}
\newcommand{\RqThreeUtilBTwoStrictNoProtectedUseBlockLo}{0.479}
\newcommand{\RqThreeUtilBTwoStrictNoProtectedUseBlockHi}{0.632}
\newcommand{\RqThreeUtilDeiGateNoProtectedUseFail}{103}
\newcommand{\RqThreeUtilDeiGateNoProtectedUseN}{171}
\newcommand{\RqThreeUtilDeiGateNoProtectedUseBlockRate}{0.602}
\newcommand{\RqThreeUtilDeiGateNoProtectedUseBlockLo}{0.528}
\newcommand{\RqThreeUtilDeiGateNoProtectedUseBlockHi}{0.673}
\newcommand{\RqThreeUtilNoDefenseNoProtectedUseFail}{92}
\newcommand{\RqThreeUtilNoDefenseNoProtectedUseN}{163}
\newcommand{\RqThreeUtilNoDefenseNoProtectedUseBlockRate}{0.564}
\newcommand{\RqThreeUtilNoDefenseNoProtectedUseBlockLo}{0.488}
\newcommand{\RqThreeUtilNoDefenseNoProtectedUseBlockHi}{0.638}
\newcommand{\RqThreeUtilBOneDevelopmentTouchedFail}{83}
\newcommand{\RqThreeUtilBOneDevelopmentTouchedN}{123}
\newcommand{\RqThreeUtilBOneDevelopmentTouchedBlockRate}{0.675}
\newcommand{\RqThreeUtilBOneDevelopmentTouchedBlockLo}{0.588}
\newcommand{\RqThreeUtilBOneDevelopmentTouchedBlockHi}{0.751}
\newcommand{\RqThreeUtilBTwoStrictDevelopmentTouchedFail}{91}
\newcommand{\RqThreeUtilBTwoStrictDevelopmentTouchedN}{123}
\newcommand{\RqThreeUtilBTwoStrictDevelopmentTouchedBlockRate}{0.740}
\newcommand{\RqThreeUtilBTwoStrictDevelopmentTouchedBlockLo}{0.656}
\newcommand{\RqThreeUtilBTwoStrictDevelopmentTouchedBlockHi}{0.809}
\newcommand{\RqThreeUtilDeiGateDevelopmentTouchedFail}{69}
\newcommand{\RqThreeUtilDeiGateDevelopmentTouchedN}{123}
\newcommand{\RqThreeUtilDeiGateDevelopmentTouchedBlockRate}{0.561}
\newcommand{\RqThreeUtilDeiGateDevelopmentTouchedBlockLo}{0.473}
\newcommand{\RqThreeUtilDeiGateDevelopmentTouchedBlockHi}{0.646}
\newcommand{\RqThreeUtilNoDefenseDevelopmentTouchedFail}{46}
\newcommand{\RqThreeUtilNoDefenseDevelopmentTouchedN}{123}
\newcommand{\RqThreeUtilNoDefenseDevelopmentTouchedBlockRate}{0.374}
\newcommand{\RqThreeUtilNoDefenseDevelopmentTouchedBlockLo}{0.294}
\newcommand{\RqThreeUtilNoDefenseDevelopmentTouchedBlockHi}{0.462}
\newcommand{\RqThreeUtilBOneDevelopmentUntouchedFail}{122}
\newcommand{\RqThreeUtilBOneDevelopmentUntouchedN}{168}
\newcommand{\RqThreeUtilBOneDevelopmentUntouchedBlockRate}{0.726}
\newcommand{\RqThreeUtilBOneDevelopmentUntouchedBlockLo}{0.654}
\newcommand{\RqThreeUtilBOneDevelopmentUntouchedBlockHi}{0.788}
\newcommand{\RqThreeUtilBTwoStrictDevelopmentUntouchedFail}{129}
\newcommand{\RqThreeUtilBTwoStrictDevelopmentUntouchedN}{168}
\newcommand{\RqThreeUtilBTwoStrictDevelopmentUntouchedBlockRate}{0.768}
\newcommand{\RqThreeUtilBTwoStrictDevelopmentUntouchedBlockLo}{0.698}
\newcommand{\RqThreeUtilBTwoStrictDevelopmentUntouchedBlockHi}{0.825}
\newcommand{\RqThreeUtilDeiGateDevelopmentUntouchedFail}{120}
\newcommand{\RqThreeUtilDeiGateDevelopmentUntouchedN}{168}
\newcommand{\RqThreeUtilDeiGateDevelopmentUntouchedBlockRate}{0.714}
\newcommand{\RqThreeUtilDeiGateDevelopmentUntouchedBlockLo}{0.642}
\newcommand{\RqThreeUtilDeiGateDevelopmentUntouchedBlockHi}{0.777}
\newcommand{\RqThreeUtilNoDefenseDevelopmentUntouchedFail}{65}
\newcommand{\RqThreeUtilNoDefenseDevelopmentUntouchedN}{168}
\newcommand{\RqThreeUtilNoDefenseDevelopmentUntouchedBlockRate}{0.387}
\newcommand{\RqThreeUtilNoDefenseDevelopmentUntouchedBlockLo}{0.317}
\newcommand{\RqThreeUtilNoDefenseDevelopmentUntouchedBlockHi}{0.462}
\newcommand{\RqThreeSweepTauThirtySecAdmit}{0}
\newcommand{\RqThreeSweepTauThirtySecN}{0}
\newcommand{\RqThreeSweepTauThirtySecRate}{--}
\newcommand{\RqThreeSweepTauThirtyBenignAdmit}{59}
\newcommand{\RqThreeSweepTauThirtyBenignN}{120}
\newcommand{\RqThreeSweepTauThirtyBenignRate}{0.492}
\newcommand{\RqThreeSweepTauFortyFiveSecAdmit}{0}
\newcommand{\RqThreeSweepTauFortyFiveSecN}{0}
\newcommand{\RqThreeSweepTauFortyFiveSecRate}{--}
\newcommand{\RqThreeSweepTauFortyFiveBenignAdmit}{59}
\newcommand{\RqThreeSweepTauFortyFiveBenignN}{120}
\newcommand{\RqThreeSweepTauFortyFiveBenignRate}{0.492}
\newcommand{\RqThreeSweepTauFiftySecAdmit}{0}
\newcommand{\RqThreeSweepTauFiftySecN}{0}
\newcommand{\RqThreeSweepTauFiftySecRate}{--}
\newcommand{\RqThreeSweepTauFiftyBenignAdmit}{59}
\newcommand{\RqThreeSweepTauFiftyBenignN}{120}
\newcommand{\RqThreeSweepTauFiftyBenignRate}{0.492}
\newcommand{\RqThreeSweepTauSixtySecAdmit}{0}
\newcommand{\RqThreeSweepTauSixtySecN}{0}
\newcommand{\RqThreeSweepTauSixtySecRate}{--}
\newcommand{\RqThreeSweepTauSixtyBenignAdmit}{0}
\newcommand{\RqThreeSweepTauSixtyBenignN}{120}
\newcommand{\RqThreeSweepTauSixtyBenignRate}{0.000}
\newcommand{\RqThreeSweepTauSixtySevenSecAdmit}{0}
\newcommand{\RqThreeSweepTauSixtySevenSecN}{0}
\newcommand{\RqThreeSweepTauSixtySevenSecRate}{--}
\newcommand{\RqThreeSweepTauSixtySevenBenignAdmit}{0}
\newcommand{\RqThreeSweepTauSixtySevenBenignN}{120}
\newcommand{\RqThreeSweepTauSixtySevenBenignRate}{0.000}
\newcommand{\RqThreeSweepEllAllTrustedSecAdmit}{0}
\newcommand{\RqThreeSweepEllAllTrustedSecN}{0}
\newcommand{\RqThreeSweepEllAllTrustedSecRate}{--}
\newcommand{\RqThreeSweepEllAllTrustedBenignAdmit}{46}
\newcommand{\RqThreeSweepEllAllTrustedBenignN}{120}
\newcommand{\RqThreeSweepEllAllTrustedBenignRate}{0.383}
  % paired RQ3 utility analysis
% AUTO-GENERATED by analyze.py from out/rq4a_all.jsonl; primary (degenerate retained); do not edit.
% Voided: 0; marked degenerate in source: 0; bootstrap seed: 20260814.
\newcommand{\RqFourAEpisodes}{824}
\newcommand{\RqFourAVoided}{0}
\newcommand{\RqFourADegenerate}{0}
\newcommand{\RqFourAAnalysisMode}{primary}
\newcommand{\RqFourASecDeiHit}{16}
\newcommand{\RqFourASecDeiN}{150}
\newcommand{\RqFourASecDeiASR}{0.107}
\newcommand{\RqFourASecDeiLo}{0.067}
\newcommand{\RqFourASecDeiHi}{0.166}
\newcommand{\RqFourASecDeiFisherP}{7.4e-13}
\newcommand{\RqFourASecDeiRD}{-0.373}
\newcommand{\RqFourASecDeiRDlo}{-0.467}
\newcommand{\RqFourASecDeiRDhi}{-0.279}
\newcommand{\RqFourASecDeiRR}{0.222}
\newcommand{\RqFourASecDeiRRlo}{0.136}
\newcommand{\RqFourASecDeiRRhi}{0.364}
\newcommand{\RqFourASecNoDefenseHit}{72}
\newcommand{\RqFourASecNoDefenseN}{150}
\newcommand{\RqFourASecNoDefenseASR}{0.480}
\newcommand{\RqFourASecNoDefenseLo}{0.402}
\newcommand{\RqFourASecNoDefenseHi}{0.559}
\newcommand{\RqFourASecBOneHit}{13}
\newcommand{\RqFourASecBOneN}{150}
\newcommand{\RqFourASecBOneASR}{0.087}
\newcommand{\RqFourASecBOneLo}{0.051}
\newcommand{\RqFourASecBOneHi}{0.143}
\newcommand{\RqFourASecBOneFisherP}{1.4e-14}
\newcommand{\RqFourASecBOneRD}{-0.393}
\newcommand{\RqFourASecBOneRDlo}{-0.485}
\newcommand{\RqFourASecBOneRDhi}{-0.302}
\newcommand{\RqFourASecBOneRR}{0.181}
\newcommand{\RqFourASecBOneRRlo}{0.105}
\newcommand{\RqFourASecBOneRRhi}{0.312}
\newcommand{\RqFourASecDeiGateHit}{9}
\newcommand{\RqFourASecDeiGateN}{150}
\newcommand{\RqFourASecDeiGateASR}{0.060}
\newcommand{\RqFourASecDeiGateLo}{0.032}
\newcommand{\RqFourASecDeiGateHi}{0.110}
\newcommand{\RqFourASecDeiGateFisherP}{2.9e-17}
\newcommand{\RqFourASecDeiGateRD}{-0.420}
\newcommand{\RqFourASecDeiGateRDlo}{-0.509}
\newcommand{\RqFourASecDeiGateRDhi}{-0.331}
\newcommand{\RqFourASecDeiGateRR}{0.125}
\newcommand{\RqFourASecDeiGateRRlo}{0.065}
\newcommand{\RqFourASecDeiGateRRhi}{0.241}
\newcommand{\RqFourASecBOneSuiteBankingHit}{0}
\newcommand{\RqFourASecBOneSuiteBankingN}{30}
\newcommand{\RqFourASecBOneSuiteBankingASR}{0.000}
\newcommand{\RqFourASecBOneSuiteBankingLo}{0.000}
\newcommand{\RqFourASecBOneSuiteBankingHi}{0.114}
\newcommand{\RqFourASecDeiSuiteBankingHit}{0}
\newcommand{\RqFourASecDeiSuiteBankingN}{30}
\newcommand{\RqFourASecDeiSuiteBankingASR}{0.000}
\newcommand{\RqFourASecDeiSuiteBankingLo}{0.000}
\newcommand{\RqFourASecDeiSuiteBankingHi}{0.114}
\newcommand{\RqFourASecDeiGateSuiteBankingHit}{0}
\newcommand{\RqFourASecDeiGateSuiteBankingN}{30}
\newcommand{\RqFourASecDeiGateSuiteBankingASR}{0.000}
\newcommand{\RqFourASecDeiGateSuiteBankingLo}{0.000}
\newcommand{\RqFourASecDeiGateSuiteBankingHi}{0.114}
\newcommand{\RqFourASecNoDefenseSuiteBankingHit}{19}
\newcommand{\RqFourASecNoDefenseSuiteBankingN}{30}
\newcommand{\RqFourASecNoDefenseSuiteBankingASR}{0.633}
\newcommand{\RqFourASecNoDefenseSuiteBankingLo}{0.455}
\newcommand{\RqFourASecNoDefenseSuiteBankingHi}{0.781}
\newcommand{\RqFourASecBOneSuiteSlackHit}{4}
\newcommand{\RqFourASecBOneSuiteSlackN}{30}
\newcommand{\RqFourASecBOneSuiteSlackASR}{0.133}
\newcommand{\RqFourASecBOneSuiteSlackLo}{0.053}
\newcommand{\RqFourASecBOneSuiteSlackHi}{0.297}
\newcommand{\RqFourASecDeiSuiteSlackHit}{5}
\newcommand{\RqFourASecDeiSuiteSlackN}{30}
\newcommand{\RqFourASecDeiSuiteSlackASR}{0.167}
\newcommand{\RqFourASecDeiSuiteSlackLo}{0.073}
\newcommand{\RqFourASecDeiSuiteSlackHi}{0.336}
\newcommand{\RqFourASecDeiGateSuiteSlackHit}{5}
\newcommand{\RqFourASecDeiGateSuiteSlackN}{30}
\newcommand{\RqFourASecDeiGateSuiteSlackASR}{0.167}
\newcommand{\RqFourASecDeiGateSuiteSlackLo}{0.073}
\newcommand{\RqFourASecDeiGateSuiteSlackHi}{0.336}
\newcommand{\RqFourASecNoDefenseSuiteSlackHit}{19}
\newcommand{\RqFourASecNoDefenseSuiteSlackN}{30}
\newcommand{\RqFourASecNoDefenseSuiteSlackASR}{0.633}
\newcommand{\RqFourASecNoDefenseSuiteSlackLo}{0.455}
\newcommand{\RqFourASecNoDefenseSuiteSlackHi}{0.781}
\newcommand{\RqFourASecBOneSuiteTravelHit}{9}
\newcommand{\RqFourASecBOneSuiteTravelN}{30}
\newcommand{\RqFourASecBOneSuiteTravelASR}{0.300}
\newcommand{\RqFourASecBOneSuiteTravelLo}{0.167}
\newcommand{\RqFourASecBOneSuiteTravelHi}{0.479}
\newcommand{\RqFourASecDeiSuiteTravelHit}{11}
\newcommand{\RqFourASecDeiSuiteTravelN}{30}
\newcommand{\RqFourASecDeiSuiteTravelASR}{0.367}
\newcommand{\RqFourASecDeiSuiteTravelLo}{0.219}
\newcommand{\RqFourASecDeiSuiteTravelHi}{0.545}
\newcommand{\RqFourASecDeiGateSuiteTravelHit}{4}
\newcommand{\RqFourASecDeiGateSuiteTravelN}{30}
\newcommand{\RqFourASecDeiGateSuiteTravelASR}{0.133}
\newcommand{\RqFourASecDeiGateSuiteTravelLo}{0.053}
\newcommand{\RqFourASecDeiGateSuiteTravelHi}{0.297}
\newcommand{\RqFourASecNoDefenseSuiteTravelHit}{24}
\newcommand{\RqFourASecNoDefenseSuiteTravelN}{30}
\newcommand{\RqFourASecNoDefenseSuiteTravelASR}{0.800}
\newcommand{\RqFourASecNoDefenseSuiteTravelLo}{0.627}
\newcommand{\RqFourASecNoDefenseSuiteTravelHi}{0.905}
\newcommand{\RqFourASecBOneSuiteWorkspaceHit}{0}
\newcommand{\RqFourASecBOneSuiteWorkspaceN}{60}
\newcommand{\RqFourASecBOneSuiteWorkspaceASR}{0.000}
\newcommand{\RqFourASecBOneSuiteWorkspaceLo}{0.000}
\newcommand{\RqFourASecBOneSuiteWorkspaceHi}{0.060}
\newcommand{\RqFourASecDeiSuiteWorkspaceHit}{0}
\newcommand{\RqFourASecDeiSuiteWorkspaceN}{60}
\newcommand{\RqFourASecDeiSuiteWorkspaceASR}{0.000}
\newcommand{\RqFourASecDeiSuiteWorkspaceLo}{0.000}
\newcommand{\RqFourASecDeiSuiteWorkspaceHi}{0.060}
\newcommand{\RqFourASecDeiGateSuiteWorkspaceHit}{0}
\newcommand{\RqFourASecDeiGateSuiteWorkspaceN}{60}
\newcommand{\RqFourASecDeiGateSuiteWorkspaceASR}{0.000}
\newcommand{\RqFourASecDeiGateSuiteWorkspaceLo}{0.000}
\newcommand{\RqFourASecDeiGateSuiteWorkspaceHi}{0.060}
\newcommand{\RqFourASecNoDefenseSuiteWorkspaceHit}{10}
\newcommand{\RqFourASecNoDefenseSuiteWorkspaceN}{60}
\newcommand{\RqFourASecNoDefenseSuiteWorkspaceASR}{0.167}
\newcommand{\RqFourASecNoDefenseSuiteWorkspaceLo}{0.093}
\newcommand{\RqFourASecNoDefenseSuiteWorkspaceHi}{0.280}
\newcommand{\RqFourASecBOneTauTuningHit}{0}
\newcommand{\RqFourASecBOneTauTuningN}{32}
\newcommand{\RqFourASecBOneTauTuningASR}{0.000}
\newcommand{\RqFourASecBOneTauTuningLo}{0.000}
\newcommand{\RqFourASecBOneTauTuningHi}{0.107}
\newcommand{\RqFourASecDeiTauTuningHit}{0}
\newcommand{\RqFourASecDeiTauTuningN}{32}
\newcommand{\RqFourASecDeiTauTuningASR}{0.000}
\newcommand{\RqFourASecDeiTauTuningLo}{0.000}
\newcommand{\RqFourASecDeiTauTuningHi}{0.107}
\newcommand{\RqFourASecDeiGateTauTuningHit}{0}
\newcommand{\RqFourASecDeiGateTauTuningN}{32}
\newcommand{\RqFourASecDeiGateTauTuningASR}{0.000}
\newcommand{\RqFourASecDeiGateTauTuningLo}{0.000}
\newcommand{\RqFourASecDeiGateTauTuningHi}{0.107}
\newcommand{\RqFourASecNoDefenseTauTuningHit}{9}
\newcommand{\RqFourASecNoDefenseTauTuningN}{32}
\newcommand{\RqFourASecNoDefenseTauTuningASR}{0.281}
\newcommand{\RqFourASecNoDefenseTauTuningLo}{0.156}
\newcommand{\RqFourASecNoDefenseTauTuningHi}{0.454}
\newcommand{\RqFourASecBOneTauUntouchedHit}{0}
\newcommand{\RqFourASecBOneTauUntouchedN}{28}
\newcommand{\RqFourASecBOneTauUntouchedASR}{0.000}
\newcommand{\RqFourASecBOneTauUntouchedLo}{0.000}
\newcommand{\RqFourASecBOneTauUntouchedHi}{0.121}
\newcommand{\RqFourASecDeiTauUntouchedHit}{0}
\newcommand{\RqFourASecDeiTauUntouchedN}{28}
\newcommand{\RqFourASecDeiTauUntouchedASR}{0.000}
\newcommand{\RqFourASecDeiTauUntouchedLo}{0.000}
\newcommand{\RqFourASecDeiTauUntouchedHi}{0.121}
\newcommand{\RqFourASecDeiGateTauUntouchedHit}{0}
\newcommand{\RqFourASecDeiGateTauUntouchedN}{28}
\newcommand{\RqFourASecDeiGateTauUntouchedASR}{0.000}
\newcommand{\RqFourASecDeiGateTauUntouchedLo}{0.000}
\newcommand{\RqFourASecDeiGateTauUntouchedHi}{0.121}
\newcommand{\RqFourASecNoDefenseTauUntouchedHit}{1}
\newcommand{\RqFourASecNoDefenseTauUntouchedN}{28}
\newcommand{\RqFourASecNoDefenseTauUntouchedASR}{0.036}
\newcommand{\RqFourASecNoDefenseTauUntouchedLo}{0.006}
\newcommand{\RqFourASecNoDefenseTauUntouchedHi}{0.177}
\newcommand{\RqFourAUtilDeiFail}{29}
\newcommand{\RqFourAUtilDeiN}{56}
\newcommand{\RqFourAUtilDeiRate}{48.2\%}
\newcommand{\RqFourAUtilDeiBlockRate}{0.518}
\newcommand{\RqFourAUtilDeiBlockLo}{0.390}
\newcommand{\RqFourAUtilDeiBlockHi}{0.643}
\newcommand{\RqFourAUtilDeiQuarEpisodes}{0}
\newcommand{\RqFourAUtilDeiQuarRate}{0.000}
\newcommand{\RqFourAUtilDeiQuarLo}{0.000}
\newcommand{\RqFourAUtilDeiQuarHi}{0.064}
\newcommand{\RqFourAUtilDeiFailDefense}{19}
\newcommand{\RqFourAUtilDeiFailModel}{10}
\newcommand{\RqFourAUtilDeiFailDegenerate}{0}
\newcommand{\RqFourAUtilDeiFisherP}{0.00}
\newcommand{\RqFourAUtilDeiMcNemarP}{0.000}
\newcommand{\RqFourAUtilDeiDiscordantB}{17}
\newcommand{\RqFourAUtilDeiDiscordantC}{1}
\newcommand{\RqFourAUtilDeiPairedN}{56}
\newcommand{\RqFourAUtilDeiPairedRD}{-0.286}
\newcommand{\RqFourAUtilDeiPairedRDlo}{-0.403}
\newcommand{\RqFourAUtilDeiPairedRDhi}{-0.152}
\newcommand{\RqFourAUtilDeiClusterN}{56}
\newcommand{\RqFourAUtilDeiClusterRDlo}{-0.411}
\newcommand{\RqFourAUtilDeiClusterRDhi}{-0.161}
\newcommand{\RqFourAUtilNoDefenseFail}{13}
\newcommand{\RqFourAUtilNoDefenseN}{56}
\newcommand{\RqFourAUtilNoDefenseRate}{76.8\%}
\newcommand{\RqFourAUtilNoDefenseBlockRate}{0.232}
\newcommand{\RqFourAUtilNoDefenseBlockLo}{0.141}
\newcommand{\RqFourAUtilNoDefenseBlockHi}{0.358}
\newcommand{\RqFourAUtilNoDefenseQuarEpisodes}{0}
\newcommand{\RqFourAUtilNoDefenseQuarRate}{0.000}
\newcommand{\RqFourAUtilNoDefenseQuarLo}{0.000}
\newcommand{\RqFourAUtilNoDefenseQuarHi}{0.064}
\newcommand{\RqFourAUtilNoDefenseFailDefense}{0}
\newcommand{\RqFourAUtilNoDefenseFailModel}{13}
\newcommand{\RqFourAUtilNoDefenseFailDegenerate}{0}
\newcommand{\RqFourAUtilBOneFail}{34}
\newcommand{\RqFourAUtilBOneN}{56}
\newcommand{\RqFourAUtilBOneRate}{39.3\%}
\newcommand{\RqFourAUtilBOneBlockRate}{0.607}
\newcommand{\RqFourAUtilBOneBlockLo}{0.476}
\newcommand{\RqFourAUtilBOneBlockHi}{0.724}
\newcommand{\RqFourAUtilBOneQuarEpisodes}{0}
\newcommand{\RqFourAUtilBOneQuarRate}{0.000}
\newcommand{\RqFourAUtilBOneQuarLo}{0.000}
\newcommand{\RqFourAUtilBOneQuarHi}{0.064}
\newcommand{\RqFourAUtilBOneFailDefense}{25}
\newcommand{\RqFourAUtilBOneFailModel}{9}
\newcommand{\RqFourAUtilBOneFailDegenerate}{0}
\newcommand{\RqFourAUtilBOneFisherP}{0.00}
\newcommand{\RqFourAUtilBOneMcNemarP}{0.000}
\newcommand{\RqFourAUtilBOneDiscordantB}{22}
\newcommand{\RqFourAUtilBOneDiscordantC}{1}
\newcommand{\RqFourAUtilBOnePairedN}{56}
\newcommand{\RqFourAUtilBOnePairedRD}{-0.375}
\newcommand{\RqFourAUtilBOnePairedRDlo}{-0.495}
\newcommand{\RqFourAUtilBOnePairedRDhi}{-0.229}
\newcommand{\RqFourAUtilBOneClusterN}{56}
\newcommand{\RqFourAUtilBOneClusterRDlo}{-0.518}
\newcommand{\RqFourAUtilBOneClusterRDhi}{-0.250}
\newcommand{\RqFourAUtilDeiGateFail}{32}
\newcommand{\RqFourAUtilDeiGateN}{56}
\newcommand{\RqFourAUtilDeiGateRate}{42.9\%}
\newcommand{\RqFourAUtilDeiGateBlockRate}{0.571}
\newcommand{\RqFourAUtilDeiGateBlockLo}{0.441}
\newcommand{\RqFourAUtilDeiGateBlockHi}{0.692}
\newcommand{\RqFourAUtilDeiGateQuarEpisodes}{1}
\newcommand{\RqFourAUtilDeiGateQuarRate}{0.018}
\newcommand{\RqFourAUtilDeiGateQuarLo}{0.003}
\newcommand{\RqFourAUtilDeiGateQuarHi}{0.094}
\newcommand{\RqFourAUtilDeiGateFailDefense}{21}
\newcommand{\RqFourAUtilDeiGateFailModel}{11}
\newcommand{\RqFourAUtilDeiGateFailDegenerate}{0}
\newcommand{\RqFourAUtilDeiGateFisherP}{0.00}
\newcommand{\RqFourAUtilDeiGateMcNemarP}{0.000}
\newcommand{\RqFourAUtilDeiGateDiscordantB}{19}
\newcommand{\RqFourAUtilDeiGateDiscordantC}{0}
\newcommand{\RqFourAUtilDeiGatePairedN}{56}
\newcommand{\RqFourAUtilDeiGatePairedRD}{-0.339}
\newcommand{\RqFourAUtilDeiGatePairedRDlo}{-0.451}
\newcommand{\RqFourAUtilDeiGatePairedRDhi}{-0.208}
\newcommand{\RqFourAUtilDeiGateClusterN}{56}
\newcommand{\RqFourAUtilDeiGateClusterRDlo}{-0.464}
\newcommand{\RqFourAUtilDeiGateClusterRDhi}{-0.214}
\newcommand{\RqFourAUtilBOneProtectedUseFail}{25}
\newcommand{\RqFourAUtilBOneProtectedUseN}{31}
\newcommand{\RqFourAUtilBOneProtectedUseBlockRate}{0.806}
\newcommand{\RqFourAUtilBOneProtectedUseBlockLo}{0.637}
\newcommand{\RqFourAUtilBOneProtectedUseBlockHi}{0.908}
\newcommand{\RqFourAUtilDeiProtectedUseFail}{19}
\newcommand{\RqFourAUtilDeiProtectedUseN}{30}
\newcommand{\RqFourAUtilDeiProtectedUseBlockRate}{0.633}
\newcommand{\RqFourAUtilDeiProtectedUseBlockLo}{0.455}
\newcommand{\RqFourAUtilDeiProtectedUseBlockHi}{0.781}
\newcommand{\RqFourAUtilDeiGateProtectedUseFail}{21}
\newcommand{\RqFourAUtilDeiGateProtectedUseN}{30}
\newcommand{\RqFourAUtilDeiGateProtectedUseBlockRate}{0.700}
\newcommand{\RqFourAUtilDeiGateProtectedUseBlockLo}{0.521}
\newcommand{\RqFourAUtilDeiGateProtectedUseBlockHi}{0.833}
\newcommand{\RqFourAUtilNoDefenseProtectedUseFail}{2}
\newcommand{\RqFourAUtilNoDefenseProtectedUseN}{28}
\newcommand{\RqFourAUtilNoDefenseProtectedUseBlockRate}{0.071}
\newcommand{\RqFourAUtilNoDefenseProtectedUseBlockLo}{0.020}
\newcommand{\RqFourAUtilNoDefenseProtectedUseBlockHi}{0.226}
\newcommand{\RqFourAUtilBOneNoProtectedUseFail}{9}
\newcommand{\RqFourAUtilBOneNoProtectedUseN}{25}
\newcommand{\RqFourAUtilBOneNoProtectedUseBlockRate}{0.360}
\newcommand{\RqFourAUtilBOneNoProtectedUseBlockLo}{0.202}
\newcommand{\RqFourAUtilBOneNoProtectedUseBlockHi}{0.555}
\newcommand{\RqFourAUtilDeiNoProtectedUseFail}{10}
\newcommand{\RqFourAUtilDeiNoProtectedUseN}{26}
\newcommand{\RqFourAUtilDeiNoProtectedUseBlockRate}{0.385}
\newcommand{\RqFourAUtilDeiNoProtectedUseBlockLo}{0.224}
\newcommand{\RqFourAUtilDeiNoProtectedUseBlockHi}{0.575}
\newcommand{\RqFourAUtilDeiGateNoProtectedUseFail}{11}
\newcommand{\RqFourAUtilDeiGateNoProtectedUseN}{26}
\newcommand{\RqFourAUtilDeiGateNoProtectedUseBlockRate}{0.423}
\newcommand{\RqFourAUtilDeiGateNoProtectedUseBlockLo}{0.255}
\newcommand{\RqFourAUtilDeiGateNoProtectedUseBlockHi}{0.611}
\newcommand{\RqFourAUtilNoDefenseNoProtectedUseFail}{11}
\newcommand{\RqFourAUtilNoDefenseNoProtectedUseN}{28}
\newcommand{\RqFourAUtilNoDefenseNoProtectedUseBlockRate}{0.393}
\newcommand{\RqFourAUtilNoDefenseNoProtectedUseBlockLo}{0.236}
\newcommand{\RqFourAUtilNoDefenseNoProtectedUseBlockHi}{0.576}
\newcommand{\RqFourAUtilBOneDevelopmentUntouchedFail}{34}
\newcommand{\RqFourAUtilBOneDevelopmentUntouchedN}{56}
\newcommand{\RqFourAUtilBOneDevelopmentUntouchedBlockRate}{0.607}
\newcommand{\RqFourAUtilBOneDevelopmentUntouchedBlockLo}{0.476}
\newcommand{\RqFourAUtilBOneDevelopmentUntouchedBlockHi}{0.724}
\newcommand{\RqFourAUtilDeiDevelopmentUntouchedFail}{29}
\newcommand{\RqFourAUtilDeiDevelopmentUntouchedN}{56}
\newcommand{\RqFourAUtilDeiDevelopmentUntouchedBlockRate}{0.518}
\newcommand{\RqFourAUtilDeiDevelopmentUntouchedBlockLo}{0.390}
\newcommand{\RqFourAUtilDeiDevelopmentUntouchedBlockHi}{0.643}
\newcommand{\RqFourAUtilDeiGateDevelopmentUntouchedFail}{32}
\newcommand{\RqFourAUtilDeiGateDevelopmentUntouchedN}{56}
\newcommand{\RqFourAUtilDeiGateDevelopmentUntouchedBlockRate}{0.571}
\newcommand{\RqFourAUtilDeiGateDevelopmentUntouchedBlockLo}{0.441}
\newcommand{\RqFourAUtilDeiGateDevelopmentUntouchedBlockHi}{0.692}
\newcommand{\RqFourAUtilNoDefenseDevelopmentUntouchedFail}{13}
\newcommand{\RqFourAUtilNoDefenseDevelopmentUntouchedN}{56}
\newcommand{\RqFourAUtilNoDefenseDevelopmentUntouchedBlockRate}{0.232}
\newcommand{\RqFourAUtilNoDefenseDevelopmentUntouchedBlockLo}{0.141}
\newcommand{\RqFourAUtilNoDefenseDevelopmentUntouchedBlockHi}{0.358}
\newcommand{\RqFourASweepTauThirtySecAdmit}{48}
\newcommand{\RqFourASweepTauThirtySecN}{194}
\newcommand{\RqFourASweepTauThirtySecRate}{0.247}
\newcommand{\RqFourASweepTauThirtyBenignAdmit}{23}
\newcommand{\RqFourASweepTauThirtyBenignN}{60}
\newcommand{\RqFourASweepTauThirtyBenignRate}{0.383}
\newcommand{\RqFourASweepTauFortyFiveSecAdmit}{48}
\newcommand{\RqFourASweepTauFortyFiveSecN}{194}
\newcommand{\RqFourASweepTauFortyFiveSecRate}{0.247}
\newcommand{\RqFourASweepTauFortyFiveBenignAdmit}{23}
\newcommand{\RqFourASweepTauFortyFiveBenignN}{60}
\newcommand{\RqFourASweepTauFortyFiveBenignRate}{0.383}
\newcommand{\RqFourASweepTauFiftySecAdmit}{48}
\newcommand{\RqFourASweepTauFiftySecN}{194}
\newcommand{\RqFourASweepTauFiftySecRate}{0.247}
\newcommand{\RqFourASweepTauFiftyBenignAdmit}{23}
\newcommand{\RqFourASweepTauFiftyBenignN}{60}
\newcommand{\RqFourASweepTauFiftyBenignRate}{0.383}
\newcommand{\RqFourASweepTauSixtySecAdmit}{0}
\newcommand{\RqFourASweepTauSixtySecN}{194}
\newcommand{\RqFourASweepTauSixtySecRate}{0.000}
\newcommand{\RqFourASweepTauSixtyBenignAdmit}{0}
\newcommand{\RqFourASweepTauSixtyBenignN}{60}
\newcommand{\RqFourASweepTauSixtyBenignRate}{0.000}
\newcommand{\RqFourASweepTauSixtySevenSecAdmit}{0}
\newcommand{\RqFourASweepTauSixtySevenSecN}{194}
\newcommand{\RqFourASweepTauSixtySevenSecRate}{0.000}
\newcommand{\RqFourASweepTauSixtySevenBenignAdmit}{0}
\newcommand{\RqFourASweepTauSixtySevenBenignN}{60}
\newcommand{\RqFourASweepTauSixtySevenBenignRate}{0.000}
\newcommand{\RqFourASweepEllAllTrustedSecAdmit}{25}
\newcommand{\RqFourASweepEllAllTrustedSecN}{194}
\newcommand{\RqFourASweepEllAllTrustedSecRate}{0.129}
\newcommand{\RqFourASweepEllAllTrustedBenignAdmit}{23}
\newcommand{\RqFourASweepEllAllTrustedBenignN}{60}
\newcommand{\RqFourASweepEllAllTrustedBenignRate}{0.383}
  % gpt-4.1-mini replication
% AUTO-GENERATED by analyze.py from out/rq4c_all.jsonl; primary (degenerate retained); do not edit.
% Voided: 20; marked degenerate in source: 0; bootstrap seed: 20260814.
\newcommand{\RqFourCEpisodes}{618}
\newcommand{\RqFourCVoided}{20}
\newcommand{\RqFourCDegenerate}{0}
\newcommand{\RqFourCAnalysisMode}{primary}
\newcommand{\RqFourCSecBOneHit}{12}
\newcommand{\RqFourCSecBOneN}{150}
\newcommand{\RqFourCSecBOneASR}{0.080}
\newcommand{\RqFourCSecBOneLo}{0.046}
\newcommand{\RqFourCSecBOneHi}{0.135}
\newcommand{\RqFourCSecBOneFisherP}{9.4e-19}
\newcommand{\RqFourCSecBOneRD}{-0.460}
\newcommand{\RqFourCSecBOneRDlo}{-0.551}
\newcommand{\RqFourCSecBOneRDhi}{-0.369}
\newcommand{\RqFourCSecBOneRR}{0.148}
\newcommand{\RqFourCSecBOneRRlo}{0.084}
\newcommand{\RqFourCSecBOneRRhi}{0.260}
\newcommand{\RqFourCSecDeiHit}{15}
\newcommand{\RqFourCSecDeiN}{150}
\newcommand{\RqFourCSecDeiASR}{0.100}
\newcommand{\RqFourCSecDeiLo}{0.062}
\newcommand{\RqFourCSecDeiHi}{0.158}
\newcommand{\RqFourCSecDeiFisherP}{8.9e-17}
\newcommand{\RqFourCSecDeiRD}{-0.440}
\newcommand{\RqFourCSecDeiRDlo}{-0.533}
\newcommand{\RqFourCSecDeiRDhi}{-0.347}
\newcommand{\RqFourCSecDeiRR}{0.185}
\newcommand{\RqFourCSecDeiRRlo}{0.112}
\newcommand{\RqFourCSecDeiRRhi}{0.306}
\newcommand{\RqFourCSecNoDefenseHit}{81}
\newcommand{\RqFourCSecNoDefenseN}{150}
\newcommand{\RqFourCSecNoDefenseASR}{0.540}
\newcommand{\RqFourCSecNoDefenseLo}{0.460}
\newcommand{\RqFourCSecNoDefenseHi}{0.618}
\newcommand{\RqFourCSecBOneSuiteBankingHit}{0}
\newcommand{\RqFourCSecBOneSuiteBankingN}{30}
\newcommand{\RqFourCSecBOneSuiteBankingASR}{0.000}
\newcommand{\RqFourCSecBOneSuiteBankingLo}{0.000}
\newcommand{\RqFourCSecBOneSuiteBankingHi}{0.114}
\newcommand{\RqFourCSecDeiSuiteBankingHit}{0}
\newcommand{\RqFourCSecDeiSuiteBankingN}{30}
\newcommand{\RqFourCSecDeiSuiteBankingASR}{0.000}
\newcommand{\RqFourCSecDeiSuiteBankingLo}{0.000}
\newcommand{\RqFourCSecDeiSuiteBankingHi}{0.114}
\newcommand{\RqFourCSecNoDefenseSuiteBankingHit}{12}
\newcommand{\RqFourCSecNoDefenseSuiteBankingN}{30}
\newcommand{\RqFourCSecNoDefenseSuiteBankingASR}{0.400}
\newcommand{\RqFourCSecNoDefenseSuiteBankingLo}{0.246}
\newcommand{\RqFourCSecNoDefenseSuiteBankingHi}{0.577}
\newcommand{\RqFourCSecBOneSuiteSlackHit}{5}
\newcommand{\RqFourCSecBOneSuiteSlackN}{30}
\newcommand{\RqFourCSecBOneSuiteSlackASR}{0.167}
\newcommand{\RqFourCSecBOneSuiteSlackLo}{0.073}
\newcommand{\RqFourCSecBOneSuiteSlackHi}{0.336}
\newcommand{\RqFourCSecDeiSuiteSlackHit}{6}
\newcommand{\RqFourCSecDeiSuiteSlackN}{30}
\newcommand{\RqFourCSecDeiSuiteSlackASR}{0.200}
\newcommand{\RqFourCSecDeiSuiteSlackLo}{0.095}
\newcommand{\RqFourCSecDeiSuiteSlackHi}{0.373}
\newcommand{\RqFourCSecNoDefenseSuiteSlackHit}{26}
\newcommand{\RqFourCSecNoDefenseSuiteSlackN}{30}
\newcommand{\RqFourCSecNoDefenseSuiteSlackASR}{0.867}
\newcommand{\RqFourCSecNoDefenseSuiteSlackLo}{0.703}
\newcommand{\RqFourCSecNoDefenseSuiteSlackHi}{0.947}
\newcommand{\RqFourCSecBOneSuiteTravelHit}{7}
\newcommand{\RqFourCSecBOneSuiteTravelN}{30}
\newcommand{\RqFourCSecBOneSuiteTravelASR}{0.233}
\newcommand{\RqFourCSecBOneSuiteTravelLo}{0.118}
\newcommand{\RqFourCSecBOneSuiteTravelHi}{0.409}
\newcommand{\RqFourCSecDeiSuiteTravelHit}{9}
\newcommand{\RqFourCSecDeiSuiteTravelN}{30}
\newcommand{\RqFourCSecDeiSuiteTravelASR}{0.300}
\newcommand{\RqFourCSecDeiSuiteTravelLo}{0.167}
\newcommand{\RqFourCSecDeiSuiteTravelHi}{0.479}
\newcommand{\RqFourCSecNoDefenseSuiteTravelHit}{26}
\newcommand{\RqFourCSecNoDefenseSuiteTravelN}{30}
\newcommand{\RqFourCSecNoDefenseSuiteTravelASR}{0.867}
\newcommand{\RqFourCSecNoDefenseSuiteTravelLo}{0.703}
\newcommand{\RqFourCSecNoDefenseSuiteTravelHi}{0.947}
\newcommand{\RqFourCSecBOneSuiteWorkspaceHit}{0}
\newcommand{\RqFourCSecBOneSuiteWorkspaceN}{60}
\newcommand{\RqFourCSecBOneSuiteWorkspaceASR}{0.000}
\newcommand{\RqFourCSecBOneSuiteWorkspaceLo}{0.000}
\newcommand{\RqFourCSecBOneSuiteWorkspaceHi}{0.060}
\newcommand{\RqFourCSecDeiSuiteWorkspaceHit}{0}
\newcommand{\RqFourCSecDeiSuiteWorkspaceN}{60}
\newcommand{\RqFourCSecDeiSuiteWorkspaceASR}{0.000}
\newcommand{\RqFourCSecDeiSuiteWorkspaceLo}{0.000}
\newcommand{\RqFourCSecDeiSuiteWorkspaceHi}{0.060}
\newcommand{\RqFourCSecNoDefenseSuiteWorkspaceHit}{17}
\newcommand{\RqFourCSecNoDefenseSuiteWorkspaceN}{60}
\newcommand{\RqFourCSecNoDefenseSuiteWorkspaceASR}{0.283}
\newcommand{\RqFourCSecNoDefenseSuiteWorkspaceLo}{0.185}
\newcommand{\RqFourCSecNoDefenseSuiteWorkspaceHi}{0.408}
\newcommand{\RqFourCSecBOneTauTuningHit}{0}
\newcommand{\RqFourCSecBOneTauTuningN}{32}
\newcommand{\RqFourCSecBOneTauTuningASR}{0.000}
\newcommand{\RqFourCSecBOneTauTuningLo}{0.000}
\newcommand{\RqFourCSecBOneTauTuningHi}{0.107}
\newcommand{\RqFourCSecDeiTauTuningHit}{0}
\newcommand{\RqFourCSecDeiTauTuningN}{32}
\newcommand{\RqFourCSecDeiTauTuningASR}{0.000}
\newcommand{\RqFourCSecDeiTauTuningLo}{0.000}
\newcommand{\RqFourCSecDeiTauTuningHi}{0.107}
\newcommand{\RqFourCSecNoDefenseTauTuningHit}{8}
\newcommand{\RqFourCSecNoDefenseTauTuningN}{32}
\newcommand{\RqFourCSecNoDefenseTauTuningASR}{0.250}
\newcommand{\RqFourCSecNoDefenseTauTuningLo}{0.133}
\newcommand{\RqFourCSecNoDefenseTauTuningHi}{0.421}
\newcommand{\RqFourCSecBOneTauUntouchedHit}{0}
\newcommand{\RqFourCSecBOneTauUntouchedN}{28}
\newcommand{\RqFourCSecBOneTauUntouchedASR}{0.000}
\newcommand{\RqFourCSecBOneTauUntouchedLo}{0.000}
\newcommand{\RqFourCSecBOneTauUntouchedHi}{0.121}
\newcommand{\RqFourCSecDeiTauUntouchedHit}{0}
\newcommand{\RqFourCSecDeiTauUntouchedN}{28}
\newcommand{\RqFourCSecDeiTauUntouchedASR}{0.000}
\newcommand{\RqFourCSecDeiTauUntouchedLo}{0.000}
\newcommand{\RqFourCSecDeiTauUntouchedHi}{0.121}
\newcommand{\RqFourCSecNoDefenseTauUntouchedHit}{9}
\newcommand{\RqFourCSecNoDefenseTauUntouchedN}{28}
\newcommand{\RqFourCSecNoDefenseTauUntouchedASR}{0.321}
\newcommand{\RqFourCSecNoDefenseTauUntouchedLo}{0.179}
\newcommand{\RqFourCSecNoDefenseTauUntouchedHi}{0.507}
\newcommand{\RqFourCUtilBOneFail}{39}
\newcommand{\RqFourCUtilBOneN}{56}
\newcommand{\RqFourCUtilBOneRate}{30.4\%}
\newcommand{\RqFourCUtilBOneBlockRate}{0.696}
\newcommand{\RqFourCUtilBOneBlockLo}{0.567}
\newcommand{\RqFourCUtilBOneBlockHi}{0.801}
\newcommand{\RqFourCUtilBOneQuarEpisodes}{0}
\newcommand{\RqFourCUtilBOneQuarRate}{0.000}
\newcommand{\RqFourCUtilBOneQuarLo}{0.000}
\newcommand{\RqFourCUtilBOneQuarHi}{0.064}
\newcommand{\RqFourCUtilBOneFailDefense}{19}
\newcommand{\RqFourCUtilBOneFailModel}{20}
\newcommand{\RqFourCUtilBOneFailDegenerate}{0}
\newcommand{\RqFourCUtilBOneFisherP}{0.00}
\newcommand{\RqFourCUtilBOneMcNemarP}{0.000}
\newcommand{\RqFourCUtilBOneDiscordantB}{20}
\newcommand{\RqFourCUtilBOneDiscordantC}{3}
\newcommand{\RqFourCUtilBOnePairedN}{56}
\newcommand{\RqFourCUtilBOnePairedRD}{-0.304}
\newcommand{\RqFourCUtilBOnePairedRDlo}{-0.436}
\newcommand{\RqFourCUtilBOnePairedRDhi}{-0.149}
\newcommand{\RqFourCUtilBOneClusterN}{56}
\newcommand{\RqFourCUtilBOneClusterRDlo}{-0.446}
\newcommand{\RqFourCUtilBOneClusterRDhi}{-0.161}
\newcommand{\RqFourCUtilNoDefenseFail}{22}
\newcommand{\RqFourCUtilNoDefenseN}{56}
\newcommand{\RqFourCUtilNoDefenseRate}{60.7\%}
\newcommand{\RqFourCUtilNoDefenseBlockRate}{0.393}
\newcommand{\RqFourCUtilNoDefenseBlockLo}{0.276}
\newcommand{\RqFourCUtilNoDefenseBlockHi}{0.524}
\newcommand{\RqFourCUtilNoDefenseQuarEpisodes}{0}
\newcommand{\RqFourCUtilNoDefenseQuarRate}{0.000}
\newcommand{\RqFourCUtilNoDefenseQuarLo}{0.000}
\newcommand{\RqFourCUtilNoDefenseQuarHi}{0.064}
\newcommand{\RqFourCUtilNoDefenseFailDefense}{0}
\newcommand{\RqFourCUtilNoDefenseFailModel}{22}
\newcommand{\RqFourCUtilNoDefenseFailDegenerate}{0}
\newcommand{\RqFourCUtilDeiFail}{35}
\newcommand{\RqFourCUtilDeiN}{56}
\newcommand{\RqFourCUtilDeiRate}{37.5\%}
\newcommand{\RqFourCUtilDeiBlockRate}{0.625}
\newcommand{\RqFourCUtilDeiBlockLo}{0.494}
\newcommand{\RqFourCUtilDeiBlockHi}{0.740}
\newcommand{\RqFourCUtilDeiQuarEpisodes}{0}
\newcommand{\RqFourCUtilDeiQuarRate}{0.000}
\newcommand{\RqFourCUtilDeiQuarLo}{0.000}
\newcommand{\RqFourCUtilDeiQuarHi}{0.064}
\newcommand{\RqFourCUtilDeiFailDefense}{16}
\newcommand{\RqFourCUtilDeiFailModel}{19}
\newcommand{\RqFourCUtilDeiFailDegenerate}{0}
\newcommand{\RqFourCUtilDeiFisherP}{0.02}
\newcommand{\RqFourCUtilDeiMcNemarP}{0.011}
\newcommand{\RqFourCUtilDeiDiscordantB}{18}
\newcommand{\RqFourCUtilDeiDiscordantC}{5}
\newcommand{\RqFourCUtilDeiPairedN}{56}
\newcommand{\RqFourCUtilDeiPairedRD}{-0.232}
\newcommand{\RqFourCUtilDeiPairedRDlo}{-0.375}
\newcommand{\RqFourCUtilDeiPairedRDhi}{-0.071}
\newcommand{\RqFourCUtilDeiClusterN}{56}
\newcommand{\RqFourCUtilDeiClusterRDlo}{-0.393}
\newcommand{\RqFourCUtilDeiClusterRDhi}{-0.071}
\newcommand{\RqFourCUtilBOneProtectedUseFail}{21}
\newcommand{\RqFourCUtilBOneProtectedUseN}{25}
\newcommand{\RqFourCUtilBOneProtectedUseBlockRate}{0.840}
\newcommand{\RqFourCUtilBOneProtectedUseBlockLo}{0.653}
\newcommand{\RqFourCUtilBOneProtectedUseBlockHi}{0.936}
\newcommand{\RqFourCUtilDeiProtectedUseFail}{19}
\newcommand{\RqFourCUtilDeiProtectedUseN}{27}
\newcommand{\RqFourCUtilDeiProtectedUseBlockRate}{0.704}
\newcommand{\RqFourCUtilDeiProtectedUseBlockLo}{0.515}
\newcommand{\RqFourCUtilDeiProtectedUseBlockHi}{0.841}
\newcommand{\RqFourCUtilNoDefenseProtectedUseFail}{5}
\newcommand{\RqFourCUtilNoDefenseProtectedUseN}{28}
\newcommand{\RqFourCUtilNoDefenseProtectedUseBlockRate}{0.179}
\newcommand{\RqFourCUtilNoDefenseProtectedUseBlockLo}{0.079}
\newcommand{\RqFourCUtilNoDefenseProtectedUseBlockHi}{0.356}
\newcommand{\RqFourCUtilBOneNoProtectedUseFail}{18}
\newcommand{\RqFourCUtilBOneNoProtectedUseN}{31}
\newcommand{\RqFourCUtilBOneNoProtectedUseBlockRate}{0.581}
\newcommand{\RqFourCUtilBOneNoProtectedUseBlockLo}{0.408}
\newcommand{\RqFourCUtilBOneNoProtectedUseBlockHi}{0.736}
\newcommand{\RqFourCUtilDeiNoProtectedUseFail}{16}
\newcommand{\RqFourCUtilDeiNoProtectedUseN}{29}
\newcommand{\RqFourCUtilDeiNoProtectedUseBlockRate}{0.552}
\newcommand{\RqFourCUtilDeiNoProtectedUseBlockLo}{0.375}
\newcommand{\RqFourCUtilDeiNoProtectedUseBlockHi}{0.716}
\newcommand{\RqFourCUtilNoDefenseNoProtectedUseFail}{17}
\newcommand{\RqFourCUtilNoDefenseNoProtectedUseN}{28}
\newcommand{\RqFourCUtilNoDefenseNoProtectedUseBlockRate}{0.607}
\newcommand{\RqFourCUtilNoDefenseNoProtectedUseBlockLo}{0.424}
\newcommand{\RqFourCUtilNoDefenseNoProtectedUseBlockHi}{0.764}
\newcommand{\RqFourCUtilBOneDevelopmentUntouchedFail}{39}
\newcommand{\RqFourCUtilBOneDevelopmentUntouchedN}{56}
\newcommand{\RqFourCUtilBOneDevelopmentUntouchedBlockRate}{0.696}
\newcommand{\RqFourCUtilBOneDevelopmentUntouchedBlockLo}{0.567}
\newcommand{\RqFourCUtilBOneDevelopmentUntouchedBlockHi}{0.801}
\newcommand{\RqFourCUtilDeiDevelopmentUntouchedFail}{35}
\newcommand{\RqFourCUtilDeiDevelopmentUntouchedN}{56}
\newcommand{\RqFourCUtilDeiDevelopmentUntouchedBlockRate}{0.625}
\newcommand{\RqFourCUtilDeiDevelopmentUntouchedBlockLo}{0.494}
\newcommand{\RqFourCUtilDeiDevelopmentUntouchedBlockHi}{0.740}
\newcommand{\RqFourCUtilNoDefenseDevelopmentUntouchedFail}{22}
\newcommand{\RqFourCUtilNoDefenseDevelopmentUntouchedN}{56}
\newcommand{\RqFourCUtilNoDefenseDevelopmentUntouchedBlockRate}{0.393}
\newcommand{\RqFourCUtilNoDefenseDevelopmentUntouchedBlockLo}{0.276}
\newcommand{\RqFourCUtilNoDefenseDevelopmentUntouchedBlockHi}{0.524}
\newcommand{\RqFourCSweepTauThirtySecAdmit}{18}
\newcommand{\RqFourCSweepTauThirtySecN}{97}
\newcommand{\RqFourCSweepTauThirtySecRate}{0.186}
\newcommand{\RqFourCSweepTauThirtyBenignAdmit}{11}
\newcommand{\RqFourCSweepTauThirtyBenignN}{27}
\newcommand{\RqFourCSweepTauThirtyBenignRate}{0.407}
\newcommand{\RqFourCSweepTauFortyFiveSecAdmit}{18}
\newcommand{\RqFourCSweepTauFortyFiveSecN}{97}
\newcommand{\RqFourCSweepTauFortyFiveSecRate}{0.186}
\newcommand{\RqFourCSweepTauFortyFiveBenignAdmit}{11}
\newcommand{\RqFourCSweepTauFortyFiveBenignN}{27}
\newcommand{\RqFourCSweepTauFortyFiveBenignRate}{0.407}
\newcommand{\RqFourCSweepTauFiftySecAdmit}{18}
\newcommand{\RqFourCSweepTauFiftySecN}{97}
\newcommand{\RqFourCSweepTauFiftySecRate}{0.186}
\newcommand{\RqFourCSweepTauFiftyBenignAdmit}{11}
\newcommand{\RqFourCSweepTauFiftyBenignN}{27}
\newcommand{\RqFourCSweepTauFiftyBenignRate}{0.407}
\newcommand{\RqFourCSweepTauSixtySecAdmit}{0}
\newcommand{\RqFourCSweepTauSixtySecN}{97}
\newcommand{\RqFourCSweepTauSixtySecRate}{0.000}
\newcommand{\RqFourCSweepTauSixtyBenignAdmit}{0}
\newcommand{\RqFourCSweepTauSixtyBenignN}{27}
\newcommand{\RqFourCSweepTauSixtyBenignRate}{0.000}
\newcommand{\RqFourCSweepTauSixtySevenSecAdmit}{0}
\newcommand{\RqFourCSweepTauSixtySevenSecN}{97}
\newcommand{\RqFourCSweepTauSixtySevenSecRate}{0.000}
\newcommand{\RqFourCSweepTauSixtySevenBenignAdmit}{0}
\newcommand{\RqFourCSweepTauSixtySevenBenignN}{27}
\newcommand{\RqFourCSweepTauSixtySevenBenignRate}{0.000}
\newcommand{\RqFourCSweepEllAllTrustedSecAdmit}{10}
\newcommand{\RqFourCSweepEllAllTrustedSecN}{97}
\newcommand{\RqFourCSweepEllAllTrustedSecRate}{0.103}
\newcommand{\RqFourCSweepEllAllTrustedBenignAdmit}{11}
\newcommand{\RqFourCSweepEllAllTrustedBenignN}{27}
\newcommand{\RqFourCSweepEllAllTrustedBenignRate}{0.407}
  % tool_knowledge replication
% AUTO-GENERATED by analyze.py from out/rq4b_workspace.jsonl; primary (degenerate retained); do not edit.
% Voided: 0; marked degenerate in source: 3; bootstrap seed: 20260814.
\newcommand{\RqFourBEpisodes}{120}
\newcommand{\RqFourBVoided}{0}
\newcommand{\RqFourBDegenerate}{3}
\newcommand{\RqFourBAnalysisMode}{primary}
\newcommand{\RqFourBSecDeiHit}{0}
\newcommand{\RqFourBSecDeiN}{60}
\newcommand{\RqFourBSecDeiASR}{0.000}
\newcommand{\RqFourBSecDeiLo}{0.000}
\newcommand{\RqFourBSecDeiHi}{0.060}
\newcommand{\RqFourBSecDeiFisherP}{1.0e+00}
\newcommand{\RqFourBSecDeiRD}{-0.017}
\newcommand{\RqFourBSecDeiRDlo}{-0.049}
\newcommand{\RqFourBSecDeiRDhi}{0.016}
\newcommand{\RqFourBSecDeiRR}{0.333}
\newcommand{\RqFourBSecDeiRRlo}{0.014}
\newcommand{\RqFourBSecDeiRRhi}{8.023}
\newcommand{\RqFourBSecNoDefenseHit}{1}
\newcommand{\RqFourBSecNoDefenseN}{60}
\newcommand{\RqFourBSecNoDefenseASR}{0.017}
\newcommand{\RqFourBSecNoDefenseLo}{0.003}
\newcommand{\RqFourBSecNoDefenseHi}{0.089}
\newcommand{\RqFourBSecDeiSuiteWorkspaceHit}{0}
\newcommand{\RqFourBSecDeiSuiteWorkspaceN}{60}
\newcommand{\RqFourBSecDeiSuiteWorkspaceASR}{0.000}
\newcommand{\RqFourBSecDeiSuiteWorkspaceLo}{0.000}
\newcommand{\RqFourBSecDeiSuiteWorkspaceHi}{0.060}
\newcommand{\RqFourBSecNoDefenseSuiteWorkspaceHit}{1}
\newcommand{\RqFourBSecNoDefenseSuiteWorkspaceN}{60}
\newcommand{\RqFourBSecNoDefenseSuiteWorkspaceASR}{0.017}
\newcommand{\RqFourBSecNoDefenseSuiteWorkspaceLo}{0.003}
\newcommand{\RqFourBSecNoDefenseSuiteWorkspaceHi}{0.089}
\newcommand{\RqFourBSecDeiTauTuningHit}{0}
\newcommand{\RqFourBSecDeiTauTuningN}{32}
\newcommand{\RqFourBSecDeiTauTuningASR}{0.000}
\newcommand{\RqFourBSecDeiTauTuningLo}{0.000}
\newcommand{\RqFourBSecDeiTauTuningHi}{0.107}
\newcommand{\RqFourBSecNoDefenseTauTuningHit}{1}
\newcommand{\RqFourBSecNoDefenseTauTuningN}{32}
\newcommand{\RqFourBSecNoDefenseTauTuningASR}{0.031}
\newcommand{\RqFourBSecNoDefenseTauTuningLo}{0.006}
\newcommand{\RqFourBSecNoDefenseTauTuningHi}{0.157}
\newcommand{\RqFourBSecDeiTauUntouchedHit}{0}
\newcommand{\RqFourBSecDeiTauUntouchedN}{28}
\newcommand{\RqFourBSecDeiTauUntouchedASR}{0.000}
\newcommand{\RqFourBSecDeiTauUntouchedLo}{0.000}
\newcommand{\RqFourBSecDeiTauUntouchedHi}{0.121}
\newcommand{\RqFourBSecNoDefenseTauUntouchedHit}{0}
\newcommand{\RqFourBSecNoDefenseTauUntouchedN}{28}
\newcommand{\RqFourBSecNoDefenseTauUntouchedASR}{0.000}
\newcommand{\RqFourBSecNoDefenseTauUntouchedLo}{0.000}
\newcommand{\RqFourBSecNoDefenseTauUntouchedHi}{0.121}
\newcommand{\RqFourBSweepTauThirtySecAdmit}{12}
\newcommand{\RqFourBSweepTauThirtySecN}{25}
\newcommand{\RqFourBSweepTauThirtySecRate}{0.480}
\newcommand{\RqFourBSweepTauThirtyBenignAdmit}{0}
\newcommand{\RqFourBSweepTauThirtyBenignN}{0}
\newcommand{\RqFourBSweepTauThirtyBenignRate}{--}
\newcommand{\RqFourBSweepTauFortyFiveSecAdmit}{12}
\newcommand{\RqFourBSweepTauFortyFiveSecN}{25}
\newcommand{\RqFourBSweepTauFortyFiveSecRate}{0.480}
\newcommand{\RqFourBSweepTauFortyFiveBenignAdmit}{0}
\newcommand{\RqFourBSweepTauFortyFiveBenignN}{0}
\newcommand{\RqFourBSweepTauFortyFiveBenignRate}{--}
\newcommand{\RqFourBSweepTauFiftySecAdmit}{12}
\newcommand{\RqFourBSweepTauFiftySecN}{25}
\newcommand{\RqFourBSweepTauFiftySecRate}{0.480}
\newcommand{\RqFourBSweepTauFiftyBenignAdmit}{0}
\newcommand{\RqFourBSweepTauFiftyBenignN}{0}
\newcommand{\RqFourBSweepTauFiftyBenignRate}{--}
\newcommand{\RqFourBSweepTauSixtySecAdmit}{0}
\newcommand{\RqFourBSweepTauSixtySecN}{25}
\newcommand{\RqFourBSweepTauSixtySecRate}{0.000}
\newcommand{\RqFourBSweepTauSixtyBenignAdmit}{0}
\newcommand{\RqFourBSweepTauSixtyBenignN}{0}
\newcommand{\RqFourBSweepTauSixtyBenignRate}{--}
\newcommand{\RqFourBSweepTauSixtySevenSecAdmit}{0}
\newcommand{\RqFourBSweepTauSixtySevenSecN}{25}
\newcommand{\RqFourBSweepTauSixtySevenSecRate}{0.000}
\newcommand{\RqFourBSweepTauSixtySevenBenignAdmit}{0}
\newcommand{\RqFourBSweepTauSixtySevenBenignN}{0}
\newcommand{\RqFourBSweepTauSixtySevenBenignRate}{--}
\newcommand{\RqFourBSweepEllAllTrustedSecAdmit}{2}
\newcommand{\RqFourBSweepEllAllTrustedSecN}{25}
\newcommand{\RqFourBSweepEllAllTrustedSecRate}{0.080}
\newcommand{\RqFourBSweepEllAllTrustedBenignAdmit}{0}
\newcommand{\RqFourBSweepEllAllTrustedBenignN}{0}
\newcommand{\RqFourBSweepEllAllTrustedBenignRate}{--}
  % open-weight replication with limited susceptibility
% AUTO-GENERATED by analyze.py from out/agentdyn_all.jsonl; primary (degenerate retained); do not edit.
% Voided: 30; marked degenerate in source: 0; bootstrap seed: 20260814.
\newcommand{\RqAgentDynEpisodes}{1120}
\newcommand{\RqAgentDynVoided}{30}
\newcommand{\RqAgentDynDegenerate}{0}
\newcommand{\RqAgentDynAnalysisMode}{primary}
\newcommand{\RqAgentDynSecDeiHit}{133}
\newcommand{\RqAgentDynSecDeiN}{560}
\newcommand{\RqAgentDynSecDeiASR}{0.237}
\newcommand{\RqAgentDynSecDeiLo}{0.204}
\newcommand{\RqAgentDynSecDeiHi}{0.274}
\newcommand{\RqAgentDynSecDeiFisherP}{4.8e-16}
\newcommand{\RqAgentDynSecDeiRD}{-0.232}
\newcommand{\RqAgentDynSecDeiRDlo}{-0.286}
\newcommand{\RqAgentDynSecDeiRDhi}{-0.178}
\newcommand{\RqAgentDynSecDeiRR}{0.506}
\newcommand{\RqAgentDynSecDeiRRlo}{0.426}
\newcommand{\RqAgentDynSecDeiRRhi}{0.601}
\newcommand{\RqAgentDynSecNoDefenseHit}{263}
\newcommand{\RqAgentDynSecNoDefenseN}{560}
\newcommand{\RqAgentDynSecNoDefenseASR}{0.470}
\newcommand{\RqAgentDynSecNoDefenseLo}{0.429}
\newcommand{\RqAgentDynSecNoDefenseHi}{0.511}
\newcommand{\RqAgentDynSecDeiSuiteDailylifeHit}{60}
\newcommand{\RqAgentDynSecDeiSuiteDailylifeN}{200}
\newcommand{\RqAgentDynSecDeiSuiteDailylifeASR}{0.300}
\newcommand{\RqAgentDynSecDeiSuiteDailylifeLo}{0.241}
\newcommand{\RqAgentDynSecDeiSuiteDailylifeHi}{0.367}
\newcommand{\RqAgentDynSecNoDefenseSuiteDailylifeHit}{167}
\newcommand{\RqAgentDynSecNoDefenseSuiteDailylifeN}{200}
\newcommand{\RqAgentDynSecNoDefenseSuiteDailylifeASR}{0.835}
\newcommand{\RqAgentDynSecNoDefenseSuiteDailylifeLo}{0.777}
\newcommand{\RqAgentDynSecNoDefenseSuiteDailylifeHi}{0.880}
\newcommand{\RqAgentDynSecDeiSuiteGithubHit}{37}
\newcommand{\RqAgentDynSecDeiSuiteGithubN}{180}
\newcommand{\RqAgentDynSecDeiSuiteGithubASR}{0.206}
\newcommand{\RqAgentDynSecDeiSuiteGithubLo}{0.153}
\newcommand{\RqAgentDynSecDeiSuiteGithubHi}{0.270}
\newcommand{\RqAgentDynSecNoDefenseSuiteGithubHit}{42}
\newcommand{\RqAgentDynSecNoDefenseSuiteGithubN}{180}
\newcommand{\RqAgentDynSecNoDefenseSuiteGithubASR}{0.233}
\newcommand{\RqAgentDynSecNoDefenseSuiteGithubLo}{0.178}
\newcommand{\RqAgentDynSecNoDefenseSuiteGithubHi}{0.300}
\newcommand{\RqAgentDynSecDeiSuiteShoppingHit}{36}
\newcommand{\RqAgentDynSecDeiSuiteShoppingN}{180}
\newcommand{\RqAgentDynSecDeiSuiteShoppingASR}{0.200}
\newcommand{\RqAgentDynSecDeiSuiteShoppingLo}{0.148}
\newcommand{\RqAgentDynSecDeiSuiteShoppingHi}{0.264}
\newcommand{\RqAgentDynSecNoDefenseSuiteShoppingHit}{54}
\newcommand{\RqAgentDynSecNoDefenseSuiteShoppingN}{180}
\newcommand{\RqAgentDynSecNoDefenseSuiteShoppingASR}{0.300}
\newcommand{\RqAgentDynSecNoDefenseSuiteShoppingLo}{0.238}
\newcommand{\RqAgentDynSecNoDefenseSuiteShoppingHi}{0.371}
\newcommand{\RqAgentDynSweepTauThirtySecAdmit}{14}
\newcommand{\RqAgentDynSweepTauThirtySecN}{209}
\newcommand{\RqAgentDynSweepTauThirtySecRate}{0.067}
\newcommand{\RqAgentDynSweepTauThirtyBenignAdmit}{0}
\newcommand{\RqAgentDynSweepTauThirtyBenignN}{0}
\newcommand{\RqAgentDynSweepTauThirtyBenignRate}{--}
\newcommand{\RqAgentDynSweepTauFortyFiveSecAdmit}{14}
\newcommand{\RqAgentDynSweepTauFortyFiveSecN}{209}
\newcommand{\RqAgentDynSweepTauFortyFiveSecRate}{0.067}
\newcommand{\RqAgentDynSweepTauFortyFiveBenignAdmit}{0}
\newcommand{\RqAgentDynSweepTauFortyFiveBenignN}{0}
\newcommand{\RqAgentDynSweepTauFortyFiveBenignRate}{--}
\newcommand{\RqAgentDynSweepTauFiftySecAdmit}{14}
\newcommand{\RqAgentDynSweepTauFiftySecN}{209}
\newcommand{\RqAgentDynSweepTauFiftySecRate}{0.067}
\newcommand{\RqAgentDynSweepTauFiftyBenignAdmit}{0}
\newcommand{\RqAgentDynSweepTauFiftyBenignN}{0}
\newcommand{\RqAgentDynSweepTauFiftyBenignRate}{--}
\newcommand{\RqAgentDynSweepTauSixtySecAdmit}{0}
\newcommand{\RqAgentDynSweepTauSixtySecN}{209}
\newcommand{\RqAgentDynSweepTauSixtySecRate}{0.000}
\newcommand{\RqAgentDynSweepTauSixtyBenignAdmit}{0}
\newcommand{\RqAgentDynSweepTauSixtyBenignN}{0}
\newcommand{\RqAgentDynSweepTauSixtyBenignRate}{--}
\newcommand{\RqAgentDynSweepTauSixtySevenSecAdmit}{0}
\newcommand{\RqAgentDynSweepTauSixtySevenSecN}{209}
\newcommand{\RqAgentDynSweepTauSixtySevenSecRate}{0.000}
\newcommand{\RqAgentDynSweepTauSixtySevenBenignAdmit}{0}
\newcommand{\RqAgentDynSweepTauSixtySevenBenignN}{0}
\newcommand{\RqAgentDynSweepTauSixtySevenBenignRate}{--}
\newcommand{\RqAgentDynSweepEllAllTrustedSecAdmit}{14}
\newcommand{\RqAgentDynSweepEllAllTrustedSecN}{209}
\newcommand{\RqAgentDynSweepEllAllTrustedSecRate}{0.067}
\newcommand{\RqAgentDynSweepEllAllTrustedBenignAdmit}{0}
\newcommand{\RqAgentDynSweepEllAllTrustedBenignN}{0}
\newcommand{\RqAgentDynSweepEllAllTrustedBenignRate}{--}
  % AgentDyn dynamic-suite replication

\title{Epistemic Security: Enforceable Admission and the Limits of Counterfactual Defense}

\author{\IEEEauthorblockN{Anonymous Author(s)}
\IEEEauthorblockA{Submission to the IEEE Symposium on Security and Privacy}}

\maketitle

\begin{abstract}
Agentic AI systems act on beliefs derived from adversarially writable
content. We formalize \emph{decision-level epistemic integrity} (DEI)
as five invariants and prove that the first four are
execution-monitor-enforceable safety properties, while the fifth,
counterfactual dependence, is a 2-safety hyperproperty that no
single-trace monitor can evaluate. We instantiate the enforceable
fragment in CZTA, a zero-trust MCP gateway, as the physical realization
of Algorithm~1 and Proposition~1, and measure its cost: mean p50
\RqFiveATwoTwentyFiveMedian\,ms at 25 req/s and A2$-$A0 overhead of
\RqFiveOverheadRange{} at 5--100 req/s. The admission result is
\Ifour, not full-DEI superiority. Deployable $\delta_3$ catches
\RqTwoSeparableDeltaThreeK/\RqTwoSeparableDeltaThreeN{} separable
commitments but misses \RqTwoDeltaThreeFn{} oracle-positive cases
(miss rate \RqTwoDeltaThreeMissRate) and falsely quarantines
\RqTwoHonestDeltaThreeK/\RqTwoHonestDeltaThreeN{} actionable controls.
DEI+gate changes completion by \RqThreeUtilDeiGatePairedRD{} relative
to no defense, but outperforms B2strict by
\RqThreeUtilDeiGateVsBTwoStrictPairedRD{} and B1 by
\RqThreeUtilDeiGateVsBOnePairedRD. On gpt-4.1-mini and
\texttt{tool\_knowledge}, ASR falls from
\RqFourASecNoDefenseASR{} and \RqFourCSecNoDefenseASR{} to
\RqFourASecDeiASR{} and \RqFourCSecDeiASR; AgentDyn further reduces ASR
from \RqAgentDynSecNoDefenseASR{} to \RqAgentDynSecDeiASR. These
implementation-specific results delineate when admission-time
provenance is enforceable and when counterfactual auditing remains
limited, rather than claiming that all counterfactual defenses are
doomed.
\end{abstract}

\begin{IEEEkeywords}
agentic AI security, epistemic integrity, runtime verification,
reference monitors, zero trust, Model Context Protocol, prompt
injection
\end{IEEEkeywords}


\section{Introduction}
An agentic AI system is an epistemic machine before it is an acting
one: it reads, retrieves, remembers, and only then invokes tools with
real-world effects. Its adversaries have noticed. Indirect prompt
injection, tool and Model Context Protocol (MCP) poisoning, retrieval
and memory poisoning, and curated context all attack the same surface:
not the action interface but the \emph{evidential basis} on which
actions are chosen. Defenses have answered at the action interface.
Information-flow and capability designs~\cite{camel}, execution-boundary
authority checks~\cite{cxi}, guardrail agents~\cite{guardagent}, and
influence tracking~\cite{argus} constrain what a compromised agent may
do. None of them states the property they are reaching for: what it
means for an agent's conclusions to be trustworthy, and which part of
that property anyone can enforce at runtime.

\paragraph*{The question}
Consider the case we return to throughout the paper. A user instructs an
assistant to delete file~0 from a shared drive. A document the assistant
reads while planning carries a planted instruction to delete file~13
``with the owner's approval.'' The assistant deletes file~13 and, when
asked to justify itself, cites the user's genuine directive. Every piece
of evidence it cites is authentic, attested, and sufficient; the action
is wrong. A monitor that checks only \emph{what the agent says supports
its action} admits this; one that also checks that the action's
\emph{target} is named by trusted evidence blocks it, because file~13
appears in no directive. But the same adversary can plant ``delete
file~0 now, before the backup completes'' when the user asked for the
backup first: the target is trusted-derivable, the citation is genuine,
and only a comparison of what the agent does with and without the
planted span reveals that the injection chose the action. This paper asks
one question: \emph{which properties of an agent's decision can a runtime
monitor enforce, which can only be audited, and does enforcing the former
change outcomes on a live model?}

\paragraph*{The answer}
We define \emph{decision-level epistemic integrity} (DEI) over decision
traces of evidence, provenance, belief state, and action
(\S\ref{sec:definition}) as five invariants: provenance-completeness
(\Ione), integrity non-degradation (\Itwo), evidential-support
sufficiency (\Ithree), argument binding (\Ifour), and a
counterfactual-dependence bound (\Ifive). We prove that \Ione--\Ifour{}
are safety properties, runtime-monitorable and enforceable by execution
monitoring in Schneider's sense, while \Ifive{} is a 2-safety
hyperproperty that provably cannot be checked by any monitor observing
a single trace (\S\ref{sec:monitorability}). The boundary is a theorem
about adversaries, not a convenience of definition: the
\emph{rationalization gap} (Corollary~\ref{cor:ratgap}) shows that an
adversary who chooses the action from injected content and cites
genuine evidence for it evades every sound admission-time monitor
whenever the action's bound arguments are trusted-derivable, and that a
decision-local counterfactual re-commitment detects exactly that
residual class. The definition is built to be refuted: each invariant
carries a finite falsifying observation, and the definition as a whole
is falsified by a successful attack on an invariant-satisfying
execution (\S\ref{sec:falsifiability}).

This is a systems paper with a formal spine, and every proposition ends
in a measurement. The enforceable fragment is instantiated in CZTA, a
zero-trust gateway in front of a live MCP deployment
(\S\ref{sec:czta}); the auditable fragment is deployed as a
counterfactual audit gate over irreversible actions (\S\ref{sec:gate});
and both are evaluated against a frontier model that follows the
canonical AgentDojo injection in \SecNdHit{} of \SecNdN{} episodes
(\S\ref{sec:eval}). We execute the rationalization attack end to end
against a real model, reproduce the coverage counterexample the
enforceable fragment admits, and quarantine it with the gate.

\paragraph*{Contributions}
\begin{enumerate}[leftmargin=1.3em,itemsep=1pt,label=\textbf{C\arabic*}]
\item \textbf{A formal, falsifiable definition with a proved
enforceability boundary.} DEI, stated relative to an explicit and
minimized epistemic TCB; the decomposition into four
execution-monitor-enforceable invariants and one 2-safety invariant; the
rationalization-gap theorem that locates the boundary; a soundness
theorem for the mapping from attack classes to invariants; and a
mechanization of the safety core in Lean~4.
\item \textbf{I4-centered admission and a conditional counterfactual gate.}
CZTA enforces the safety fragment on a live MCP deployment; the
measured replacement gateway adds \RqFiveOverheadRange{} p50 overhead
at 5--100 req/s (mean p50 \RqFiveATwoTwentyFiveMedian\,ms at 25 req/s); the gate
re-commits irreversible actions under pre-registered perturbations and
quarantines those whose commitment does not survive. We give the
architecture, the end-to-end workflow on the running example, and a
security analysis of the gateway's own attack surface, with miss,
false-quarantine, and utility costs reported as first-class results.
\item \textbf{A live-model evaluation under a registered design.}
Against gemini-2.5-flash on AgentDojo, a pilot shows the attack family
succeeds in \SecNdASR{} of episodes and that a protected-tool filter
already reduces it to \SecGateASR{} ($p{=}\GateVsNdFisherP$) at a utility
cost $n{=}\UtilGateN$ per arm cannot resolve; the registered campaign
separates the cases used to develop the gateway from the held-out cases
used to score it, measures admission against that filter, against
argument--instruction matching, and against \Ifour{} alone, measures the
gate where the enforceable fragment is blind, and powers the utility
comparison on a pre-registered minimal important difference.
The registered study ties admission and B1 at ASR
\SpSecDeiASR{} while no defense reaches \SpSecNoDefenseASR, localizing
the security result to \Ifour. Deployable $\delta_3$ catches
\RqTwoSeparableDeltaThreeK/\RqTwoSeparableDeltaThreeN{} separable
commitments but has miss rate \RqTwoDeltaThreeMissRate{} and false
quarantine rate \RqTwoHonestDeltaThreeRate. DEI+gate changes completion
by \RqThreeUtilDeiGatePairedRD{} versus no defense, but improves it by
\RqThreeUtilDeiGateVsBTwoStrictPairedRD{} over B2strict and by
\RqThreeUtilDeiGateVsBOnePairedRD{} over B1. On gpt-4.1-mini and
\texttt{tool\_knowledge}, the defensive direction replicates.
\end{enumerate}

\paragraph*{Scope of the novelty claim}
Within a registered search through 17~July~2026 (\S\ref{sec:related}),
no prior work gives a formal, falsifiable definition of decision-level
epistemic integrity that separates runtime-enforceable from
auditable-only invariants with supporting theory, proves where the
boundary lies, and instantiates the enforceable fragment in a working
monitor. We adopt the invariance condition of~\cite{qu} as \Ifive{} and
claim neither first use of the phrase~\cite{shen} nor first statement
of decision invariance.


\section{Setting, Threat Model, and Positioning}
\label{sec:setting}

\subsection{The epistemic attack surface}
\label{sec:taxonomy}
We organize attacks on the evidential basis of agent decisions by a
ten-class taxonomy of primary attack or failure mechanisms:
misinformation and disinformation; hallucination and factuality
failure; prompt injection and jailbreak; provenance or identity
tampering; retrieval and knowledge poisoning; data or model
poisoning and backdoors; confidence or trust manipulation; tool or
agent hijacking (including MCP tool-description poisoning); drift,
staleness, or benign reliability failure; and a non-attack class for
defense-only, benchmark-only, and evaluation records. The taxonomy was
produced by a registered PRISMA-ScR scoping review~\cite{prisma} of
published work on this surface (6,512 records across six databases;
552 retained as core evidence) whose automated labeling was audited
against two blinded domain experts and an adjudicator on a
hash-registered 20\,\% sample ($n{=}114$): class agreement $\kappa =
0.63$ / Gwet's $AC_1 = 0.65$ between experts, and $\kappa = 0.48$
between experts and pipeline against the pipeline's self-agreement of
0.83. Protocol, ledger, codebook, instruments, hashes, and amendment
log ship with the artifact~\cite{artifact}. The review contributes exactly two things to
this paper: the classes that the crosswalk of
\S\ref{sec:falsifiability} maps onto invariants, and the corpus
prevalence of each class (Appendix~\ref{app:crosswalk}), which is what lets
us say which invariants matter most often. Nothing else in the paper
depends on it.

Categories are mutually exclusive by fiat: where mechanisms co-occur,
the record's own framing decides. One registered disambiguation matters
for the formalization: \emph{retrieval and knowledge poisoning} requires
that corrupted content \emph{enter} the corpus, index, or memory, which
distinguishes substrate corruption (a provenance failure, \Ione) from
manipulation of which authentic evidence is selected, ranked, or
surfaced from an intact substrate (a sufficiency failure, \Ithree); the
latter is carried as a boundary case in the crosswalk.

\subsection{Threat model}
\label{sec:threat}
The adversary controls the untrusted channels: they may inject
instructions or content into retrieved documents and tool results
(indirect prompt injection), poison tool descriptions or servers (MCP
tool poisoning), write to persistent memory or retrieval stores
(memory and RAG poisoning), and shape which evidence the agent sees
(context-selection manipulation). The adversary does not control the
epistemic TCB of \S\ref{sec:tcbdef} or hold its keys. Consistent with
the zero-trust posture now advocated for internal agent
deployments~\cite{controlroadmap}, the model itself is treated as
potentially compromised: a successful injection makes the model an
adversarial component, and every guarantee below must survive that.

Two attack families deserve explicit placement because they live inside
the model rather than on a channel. \emph{Reasoning manipulation and
chain-of-thought poisoning}, in which injected content steers the
model's intermediate reasoning without appearing in any cited evidence,
are, in our terms, attacks on an untrusted component: DEI makes no claim
about the model's internal process, so such attacks surface either as
violations of the counterfactual bound (\Ifive, when removing the
injected span changes the decision) or not at all (when the decision
would have been the same). \emph{Parameter poisoning and backdoors}
triggered by trusted input are declared out of scope: a backdoor fired
by a legitimate directive yields an invariant-satisfying wrong action.
The gateway's own attack surface---what an adversary can do \emph{to}
the monitor rather than around it---is analyzed in \S\ref{sec:security}.

\subsection{Positioning}
\label{sec:positioning}
Out-of-band defenses in the CaMeL lineage~\cite{camel,oob} and
sink-authorization gates such as CXI~\cite{cxi} constrain what a
proposal may \emph{do}; we constrain what may count as a \emph{reason}
for it and prove which part of that constraint a monitor can enforce.
Intent-to-execution integrity~\cite{qu} states the invariance
condition we adopt as \Ifive{} but stops short of a monitor, a
monitorability analysis, or a falsification protocol; agent epistemic
integrity~\cite{shen} prescribes an architecture. Table~\ref{tab:positioning}
records the comparison; \S\ref{sec:related} gives the landscape.

\begin{table}[t]
\centering
\caption{Positioning against the closest prior work.
$\tick$ = provided; $\partl$ = partial; --- = absent.
\emph{Split} = runtime-enforceable separated from auditable-only
guarantees, with supporting theory. All rows verified against the
public record on 17 July 2026.}
\label{tab:positioning}
\scriptsize
\setlength{\tabcolsep}{3pt}
\begin{tabularx}{\columnwidth}{@{}>{\raggedright\arraybackslash}X c c c c c >{\centering\arraybackslash}p{1.05cm}@{}}
\toprule
& Def. & Falsif. & Monit. & Split & Runtime & Dec.-layer \\
\midrule
CaMeL~\cite{camel} & \absent & \absent & \absent & \absent & \tick & \absent \\
GuardAgent~\cite{guardagent} & \absent & \absent & \absent & \absent & \tick & \partl \\
Intent-to-exec.~\cite{qu} & \partl & \absent & \absent & \absent & \absent & \tick \\
Agent epist.\ int.~\cite{shen} & \partl & \absent & \absent & \absent & \absent & \tick \\
CXI~\cite{cxi} & \tick & \absent & \absent & \partl & \tick & \absent \\
ARGUS~\cite{argus} & \absent & \absent & \absent & \absent & \tick & \partl \\
OOB systematization~\cite{oob} & \absent & \partl & \absent & \absent & \absent & \partl \\
\midrule
\textbf{This work} & \tick & \tick & \tick & \tick & \tick & \tick \\
\bottomrule
\end{tabularx}
\end{table}


\section{Decision-Level Epistemic Integrity}
\label{sec:definition}
Prior work formalizes what an agent may \emph{do}: which tools it may
invoke, which flows may reach which sinks~\cite{camel,cxi,argus}. We
formalize what an agent may \emph{conclude}: the conditions under which
the beliefs that support a consequential action can be trusted, and
which of those conditions a runtime monitor can actually enforce. Every
object the invariants quantify over is declared once, in
Definition~\ref{def:policy}, before any invariant uses it; the running
example of \S\ref{sec:setting} is instantiated as soon as the objects
exist and is revisited after each result.

\subsection{Policy Objects and Decision Traces}
\label{sec:prelim}

\begin{definition}[Integrity lattice, labels, and policy objects]
\label{def:policy}
Let $(L, \sqsubseteq, \sqcap, \sqcup)$ be a finite lattice of integrity
labels with top $\top$ (maximally trusted) and bottom $\bot$
(untrusted), in the sense of Biba~\cite{biba} and
Denning~\cite{denning}. Let $\mathit{Ch}$ be the finite set of ingress
channels (user turn, tool result, retrieved document, memory read, model
output), $C$ the set of action classes, and $\mathit{Slots}$ the set of
named argument positions. A \emph{policy} $\Pi$ fixes:
\begin{itemize}[leftmargin=1.15em,itemsep=1pt]
\item an \emph{ingress labeling} $\lambda : \mathit{Ch} \to L$ and a
\emph{floor} $\lmin \in L$; model-generated and external content default
to low labels;
\item \emph{protected classes} $\Cprot \subseteq C$ (irreversible,
external, or safety-relevant effects), a \emph{required label}
$\lreq : \Cprot \to L$, and a \emph{sufficiency threshold}
$\tau : \Cprot \to [0,1]$ (written $\tc$ for $\tau(c)$);
\item a \emph{binding policy} $\bindpol$ assigning to each
$c \in \Cprot$ a set $\bindpol(c) \subseteq \mathit{Slots}$ of bound
argument slots and, for each $s \in \bindpol(c)$, a required label
$\ell_{\bindpol}(c, s) \in L$;
\item an \emph{endorsement policy} $\pi$, a finite set of authorized
label-raising rules (e.g., $k$-source corroboration, verifier-tool
confirmation, human sign-off), each with a predicate
$\mathsf{Authorized}(\rho, b, \ell')$ stating when rule $\rho \in \pi$
may raise belief $b$ to label $\ell'$;
\item a \emph{tolerance} $\varepsilon : \Cprot \to \mathbb{R}_{\ge 0}$
and a class $\Dirrel(q)$ of task-irrelevant perturbations
(Def.~\ref{def:i5}).
\end{itemize}
Every evidence atom and belief carries a label $\ell(\cdot) \in L$;
for an atom $e$ ingested on channel $\mathit{ch}$, $\ell(e) :=
\lambda(\mathit{ch})$. The concrete lattice used throughout the paper
is $\mathrm{UNTRUSTED} \sqsubset \mathrm{MEDIUM} \sqsubset
\mathrm{TRUSTED}$, with $\lambda(\text{user turn}) = \mathrm{TRUSTED}$
(session-signed), $\lambda(\text{tool result}) = \lambda(\text{memory})
= \mathrm{MEDIUM}$ (channel-attested), and $\lambda(\text{model
output}) = \mathrm{UNTRUSTED}$; the deployed values of the remaining
objects are given in \S\ref{sec:evalsetup}.
\end{definition}

\begin{definition}[Evidence and attestation]
\label{def:evidence}
An \emph{evidence atom} $e \in E$ is an immutable record of ingested
content with label $\ell(e)$, channel $\mathit{ch}(e)$, source metadata,
and an optional attestation. The predicate $\att(e)$ holds iff $e$
carries a signature chain that verifies against a trusted root and
whose validity window covers the current time. The freshness clause
subsumes long-horizon staleness: expired attestations fail $\att$, so
drift over long deployments surfaces as a provenance-completeness
violation rather than a silent decay.
\end{definition}

\paragraph*{What an attestation is, concretely}
In CZTA (\S\ref{sec:czta}) an attestation is the tuple (channel
credential, content digest, issue time, validity window, signature).
For user turns the credential is the session's RS256 token, verified
against the identity provider's published JWKS; for tool results and
memory reads it is the gateway's Ed25519 receipt over the content
digest, issued at ingress and chained into the ledger. The trusted root
is thus the realm key for user channels and the ledger key otherwise;
windows are the token expiry and a per-channel time-to-live; $\att$
fails on a bad signature, an unknown signer, a digest mismatch, or a
lapsed window. The definition does not depend on this choice; it is
stated so that the TCB of \S\ref{sec:tcbdef} is concrete.

\begin{definition}[Beliefs, provenance, updates, endorsements]
\label{def:beliefs}
A \emph{belief} $b \in B$ is a labeled assertion the agent system holds
at runtime, with dependency set $\deps(b) \subseteq E \cup B$ and a
finite set $\vals(b) \subseteq V$ of the values it names (identifiers,
principals, amounts, file names), extracted at creation time by the
recorder from the belief's content. The \emph{provenance graph} $P$ is
the directed acyclic graph whose nodes are $E \cup B$ and whose edges
run from each $b$ to every element of $\deps(b)$; it is recorded in an
append-only, hash-chained log (we use the W3C~PROV vocabulary for the
edge kinds~\cite{prov}, a recording convention that carries no semantic
commitment). Two event kinds modify beliefs:
\begin{itemize}[leftmargin=1.15em,itemsep=1pt]
\item an \emph{update} $\ev{upd}(b', X)$ creates $b'$ with
$\deps(b') = X$ and default label
$\ell(b') := \bigsqcap_{x \in X} \ell(x)$. This is the Biba derivation
rule~\cite{biba}; its guarantee is that a derived belief's label is a
certified lower bound on the trust of everything it was derived from,
so no derivation can manufacture trust;
\item an \emph{endorsement} $\ev{end}(b, \ell', \rho)$ raises
$\ell(b)$ to $\ell'$ under a rule $\rho \in \pi$ with
$\mathsf{Authorized}(\rho, b, \ell')$, in the sense of endorsement in
decentralized IFC~\cite{myersliskov}. Endorsements are logged,
attributable, and rate-limitable; they are the only lawful label
increases.
\end{itemize}
\end{definition}

\begin{definition}[Actions, arguments, and decision traces]
\label{def:actions}
An \emph{action event} $\ev{act}(a)$ has class $c(a) \in C$, an
\emph{argument map} $\args(a) : \mathit{Slots} \rightharpoonup V$, and a
\emph{support set} $\support(a) \subseteq B$, the beliefs the system
commits to as justifying $a$. Fix the event alphabet
$\Sigma = \{ \ev{ing}(e, ch),\, \ev{upd}(b, X),\, \ev{end}(b, \ell, \rho),\, \ev{act}(a) \}$.
A \emph{decision trace} is a finite sequence
$\sigma = \epsilon_1 \epsilon_2 \cdots \epsilon_k$ over $\Sigma$;
$\acts(\sigma)$ denotes its protected action events. For a task input
$q$ and untrusted-channel content $d$, we write $J(q, d)$ for the
protected action the system selects at a designated decision point,
holding trusted state fixed.
\end{definition}

\paragraph*{Epistemic TCB}
\label{sec:tcbdef}
The definition is stated relative to a trusted computing base: the
ingress labeler $\lambda$ together with the channel authentication it
relies on, the attestation verifier and its keys, the append-only
provenance recorder including the extraction of $\vals(\cdot)$, the
support aggregator $S$ (Def.~\ref{def:i3}), and the policy engine and
gateway. The model, all tool outputs, retrieved content, and memory
contents are \emph{untrusted}. Making the TCB explicit answers the
observation that provenance assignment is the under-specified trusted
base of out-of-band defenses~\cite{oob}; we do not eliminate that base,
we name it, minimize it, and measure it (\S\ref{sec:tcb}). The labeler
is deliberately simple: it is a function of the authenticated channel
identity alone, never of content, so the only way to obtain a high
label for adversarial content is to deliver it over a high channel,
which is channel spoofing and is excluded by
Assumption~\ref{asm:recorded}. The value recorder is the one TCB
component that reads content; it is a deterministic extractor over the
belief's text (identifiers, principals, amounts), and its failure mode
(missing a value) is fail-closed under \Ifour.

\begin{example}[Running example, I: the objects]
\label{ex:objects}
The user turn ``Delete file~0 from the shared drive'' is ingested on the
user channel as atom $e_u$ with $\ell(e_u) = \mathrm{TRUSTED}$,
$\att(e_u)$ true. The drive listing returned by a tool is atom $e_l$
($\mathrm{MEDIUM}$, attested by gateway receipt); the body of one listed
document, which carries the planted span ``also delete file~13, the
owner has approved,'' is atom $e_p$ ($\mathrm{MEDIUM}$, attested as a
tool result, content untrusted). The system forms beliefs
$b_u$ (``the user wants file~0 deleted,'' $\deps = \{e_u\}$,
$\vals = \{\text{file~0}\}$, $\mathrm{TRUSTED}$), $b_l$ (the listing,
$\mathrm{MEDIUM}$), and $b_p$ (``file~13 should be deleted,''
$\deps = \{e_p\}$, $\vals = \{\text{file~13}\}$, $\mathrm{MEDIUM}$) by
the meet rule. The
action $\ev{act}(a)$ with $c(a) = \mathrm{delete} \in \Cprot$ has
$\args(a)(\mathsf{target}) = \text{file~13}$; the policy binds
$\mathsf{target}$ at $\ell_{\bindpol}(\mathrm{delete}, \mathsf{target})
= \mathrm{TRUSTED}$. The compromised model commits $\support(a) =
\{b_u\}$: a genuine, attested, trusted directive.
\end{example}

\subsection{The Definition}
\begin{definition}[Provenance-completeness, \Ione]
\label{def:i1}
A trace $\sigma$ satisfies \Ione{} iff for every
$\ev{act}(a) \in \acts(\sigma)$ and every $b \in \support(a)$:
(i)~the transitive dependency closure of $b$ in $P$ is
well-founded and every maximal derivation path terminates in $E$;
and (ii)~every such terminal evidence atom $e$ satisfies $\att(e)$
and $\ell(e) \sqsupseteq \lmin$. Informally: no protected action
rests on an orphan belief (no evidential root, the formal shape
of a fabricated claim) or on an unattested or sub-floor root.
\end{definition}

\begin{definition}[Integrity non-degradation, \Itwo]
\label{def:i2}
A trace $\sigma$ satisfies \Itwo{} iff every label increase on any
belief is effected by a logged endorsement event authorized by
$\pi$; equivalently, for every update $\ev{upd}(b', X)$,
$\ell(b') \sqsubseteq \bigsqcap_{x \in X} \ell(x)$, and no other
event raises a label. Informally: untrusted input can never
silently launder itself into trusted belief; every upgrade has an
accountable, policy-cited justification.
\end{definition}

\begin{definition}[Evidential-support sufficiency, \Ithree]
\label{def:i3}
Fix an audited aggregator $S : 2^{B} \times L \to [0,1]$ scoring a
support set against a required label, subject to three stated
requirements: \emph{(A1) qualification}: $S(U, \ell)$ depends only
on evidence at or above $\ell$ that is attested;
\emph{(A2) monotonicity}: adding qualifying evidence never
decreases the score; \emph{(A3) auditability}: $S$ is fixed,
computable from the recorded prefix state, and inside the TCB.
$S$ is instantiable as Dempster--Shafer belief $\mathrm{Bel}(\cdot)$
over the qualifying evidence~\cite{shafer}, or as an
accepted-argument criterion in an argumentation
framework~\cite{dung}. A trace satisfies \Ithree{} iff for every
$\ev{act}(a) \in \acts(\sigma)$,
$S(\support(a), \lreq(c(a))) \geq \tau_{c(a)}$. Informally:
consequential actions require enough qualifying evidence, with
per-class thresholds.
\end{definition}

\begin{definition}[Argument binding, \Ifour]
\label{def:i4}
A trace $\sigma$ satisfies \Ifour{} iff for every
$\ev{act}(a) \in \acts(\sigma)$ and every bound slot
$s \in \bindpol(c(a))$ with $\args(a)(s)$ defined, there exists
$b \in \support(a)$ with $\args(a)(s) \in \vals(b)$ and
$\ell(b) \sqsupseteq \ell_{\bindpol}(c(a), s)$. Informally: every
argument that determines what the action does to the world is named by
a cited belief of sufficient standing. The model may choose the action;
it may not choose its targets. \Ifour{} is the decision-layer
counterpart of the authority binding that context-to-execution
integrity enforces at the sink~\cite{cxi}: CXI binds a field value to
the manifest that authorizes its effect, \Ifour{} binds an argument
value to the evidence that justifies naming it. An argument slot left
unbound by $\bindpol$ (free-text message bodies, say) is unconstrained
by \Ifour{} and falls to \Ifive.
\end{definition}

\begin{definition}[Counterfactual-dependence bound, \Ifive]
\label{def:i5}
Let $\Dirrel(q)$ be a class of task-irrelevant perturbations of
untrusted content (perturbations that preserve the task-relevant
semantics of $d$; operationalized by a generator suite,
\S\ref{sec:falsifiability}), and let $\dist$ be a metric on
protected actions (0/1 for discrete choices). A system satisfies
\Ifive{} iff for all $\delta \in \Dirrel(q)$,
$\dist\big(J(q, d),\, J(q, d \oplus \delta)\big) \leq \varepsilon_{c}$
for the class $c$ of $J(q,d)$. This is the invariance condition
articulated for judgment integrity in~\cite{qu}, restated over decision
traces; unlike \Ione--\Ifour{} it relates \emph{pairs} of executions.
\end{definition}

\begin{definition}[Decision-level epistemic integrity]
\label{def:dei}
Write $\Lift(P) = \{\,T : \forall \sigma \in T.\; \sigma \models P\,\}$
for the set-level lifting of a trace property $P$. A system with
reachable-trace set $T$ satisfies \DEI{} with respect to policy
$\Pi$ iff
\[
T \;\in\; \underbrace{\Lift(\Ione \wedge \Itwo \wedge \Ithree \wedge \Ifour)}_{\text{safety fragment (enforceable)}} \;\cap\; \underbrace{\Ifive}_{\text{2-safety (auditable)}} .
\]
\DEI{} is thus, by construction, a hyperproperty whose safety fragment
and whose properly $k$-safety fragment are named separately: the form in
which the enforceability results of \S\ref{sec:monitorability} are
stated.
\end{definition}

\begin{example}[Running example, II: what the invariants decide]
\label{ex:invariants}
On the trace of Example~\ref{ex:objects}: \Ione{} holds ($b_u$ roots in
the attested, trusted $e_u$); \Itwo{} holds (no label was raised);
\Ithree{} holds (a trusted directive is $\tc$-sufficient for a single
deletion under the deployed scorer); \Ifour{} \emph{fails}: the bound
slot $\mathsf{target}$ carries file~13, and no cited belief at
$\mathrm{TRUSTED}$ names file~13. The monitor blocks at admission. Now
change the plant to ``delete file~0 now, before the backup completes,''
with the user having asked for the backup first: the model commits
$\args(a)(\mathsf{target}) = \text{file~0}$ with the same support, and
\Ione--\Ifour{} all hold. The action is wrong because its timing was chosen by the plant, but
every single-trace check passes; only \Ifive{}---re-running the
decision with the plant removed---distinguishes this trace from an
honest one. Corollary~\ref{cor:ratgap} shows this is a limit of
admission-time checking as such, not of our particular invariants.
\end{example}

\subsection{Monitorability: What a Gateway Can Enforce}
\label{sec:monitorability}
The invariants deliberately fall on both sides of a classical
boundary, and the definition's honesty depends on saying which
side each is on. Algorithm~\ref{alg:monitor} makes the
enforceable side concrete: a single event-admission procedure whose block
branches are exactly the bad prefixes of
Proposition~\ref{prop:safety}. Two assumptions are stated explicitly
because the proofs depend on them.

\begin{algorithm}[t]
\caption{\DEI{} reference monitor: event admission (enforces
$\Ione \wedge \Itwo \wedge \Ithree \wedge \Ifour$ by truncation;
Prop.~\ref{prop:safety}).}
\label{alg:monitor}
\footnotesize
\begin{algorithmic}[1]
\Require event $\epsilon$; state $(\ell, P, \cert, \vals, \pi)$; policy
$\Pi = (\lmin, \lreq, \tau, \bindpol, \Cprot)$
\If{$\epsilon = \ev{ing}(e, ch)$}
  \State $\ell(e) \gets \lambda(ch)$;\;
  $\cert(e) \gets \att(e) \wedge \ell(e) \sqsupseteq \lmin$
  \State \textbf{admit}
\ElsIf{$\epsilon = \ev{upd}(b', X, \ell_{\mathrm{prop}})$}
  \If{$\ell_{\mathrm{prop}} \not\sqsubseteq \bigsqcap_{x \in X} \ell(x)$}
    \State \textbf{block} \Comment{\Itwo: silent upgrade}
  \EndIf
  \State $\ell(b') \gets \ell_{\mathrm{prop}}$;\;
  $\cert(b') \gets \bigwedge_{x \in X} \cert(x)$;\; $\vals(b') \gets \mathsf{extract}(b')$ \Comment{$O(|X|)$, memoized}
  \State \textbf{admit}
\ElsIf{$\epsilon = \ev{end}(b, \ell', \rho)$}
  \If{$\rho \notin \pi$ or $\neg\,\mathsf{Authorized}(\rho, b, \ell')$}
    \State \textbf{block} \Comment{\Itwo: unauthorized endorsement}
  \EndIf
  \State $\ell(b) \gets \ell'$;\; $\mathsf{log}(b, \ell', \rho, \mathrm{principal})$;\; \textbf{admit}
\ElsIf{$\epsilon = \ev{act}(a)$}
  \If{$c(a) \notin \Cprot$} \State \textbf{admit}
  \EndIf
  \ForAll{$b \in \support(a)$}
    \If{$\neg\,\cert(b)$ or $\neg\,\mathsf{WindowsValid}(b)$}
      \State \textbf{block} \Comment{\Ione; freshness per Def.~\ref{def:evidence}}
    \EndIf
  \EndFor
  \If{$S(\support(a), \lreq(c(a))) < \tau_{c(a)}$}
    \State \textbf{block} \Comment{\Ithree: insufficient support}
  \EndIf
  \ForAll{$s \in \bindpol(c(a))$ with $\args(a)(s)$ defined}
    \If{$\neg\,\exists b \in \support(a):\; \args(a)(s) \in \vals(b) \wedge \ell(b) \sqsupseteq \ell_{\bindpol}(c(a), s)$}
      \State \textbf{block} \Comment{\Ifour: unbound argument}
    \EndIf
  \EndFor
  \State \textbf{admit}
\EndIf
\Statex \textbf{block} $\equiv$ suppress $\epsilon$ and log the
verdict (truncation); \textbf{admit} $\equiv$ append $\epsilon$ to
the trace.
\end{algorithmic}
\end{algorithm}

\begin{assumption}[Complete mediation]
\label{asm:mediation}
The gateway mediates every $\Sigma$-event that can affect a
protected action; there is no side channel from untrusted content
to belief state or action that bypasses admission.
\end{assumption}

\begin{assumption}[Trusted TCB]
\label{asm:recorded}
The epistemic TCB of \S\ref{sec:tcbdef} behaves as specified:
(i)~channel identity is authenticated, so $\lambda$ assigns to every
atom the label its channel carries by policy; (ii)~the recorder
extracts $\vals(\cdot)$ faithfully and the log is append-only;
(iii)~the attestation verifier, the aggregator $S$, and the policy
engine compute what their definitions state. Compromise of any
component voids all guarantees (\S\ref{sec:limitations}).
\end{assumption}

\begin{proposition}[\Ione--\Ifour{} are enforceable safety properties]
\label{prop:safety}
Over decision traces, \Ione, \Itwo, \Ithree, and \Ifour{} are safety
properties in the sense of Alpern--Schneider~\cite{alpern}: every
violating trace has a finite bad prefix (the offending action or
update event) after which no extension repairs the violation.
Consequently (i)~they are monitorable, and an
$\mathrm{LTL}_3$-style monitor~\cite{bauer,pnueli} issues the
verdict $\bot$ exactly at the bad prefix; and (ii)~they lie in
the class of policies enforceable by execution
monitoring~\cite{schneider}: a reference monitor that truncates
the trace at the bad prefix (blocking the action or
update) enforces them, not merely detects them.
\end{proposition}

\begin{proof}[Proof sketch]
For \Ione, \Ithree, and \Ifour{} the candidate bad event is a protected
action; its support set, the relevant provenance subgraph, the
aggregator score, the cited beliefs' value sets, and the action's
argument map are all determined by the finite prefix at that event, so
violation is a property of the prefix and is irremediable once the
action executes. For \Itwo{} the bad event is a label-raising
transition without a matching authorized endorsement in the same
prefix. Quantification over runtime objects (beliefs, evidence) is
handled by parametric monitoring~\cite{mop}: the object universe at
any prefix is the finite set of objects observed so far, so each
instantiated monitor is a finite-state safety monitor. Enforcement
follows because blocking the bad event yields a legal trace, i.e., the
monitor is a truncation automaton in Schneider's
sense~\cite{schneider}. The full proof, including the two
legality-preservation facts, is Appendix~\ref{app:proofs}. The
\Ione--\Ifour{} clauses and the two legality facts are mechanized in
Lean~4 (finite-trace form; kernel-checked; the only axiom is
\texttt{propext}, on two lemmas), re-checkable in the artifact via
\texttt{formal/check\_formal.sh}.
\end{proof}

\begin{remark}[Online cost]
With per-belief \emph{root certificates} (a memoized, hash-chained
receipt that a belief's evidential roots are attested and
$\sqsupseteq \lmin$, established when the belief is created) the
\Ione{} check at an action is $O(|\support(a)|)$ constant-time receipt
verifications; \Itwo{} is one lattice meet per update; \Ithree{} is one
aggregator evaluation; \Ifour{} is $O(|\bindpol(c(a))| \cdot
|\support(a)|)$ set-membership tests on recorded value sets. These are
the operations whose cost CZTA measures (Table~\ref{tab:overhead}); the
reference implementation confirms the asymptotics, with memoized
\Ione{} checks flat across a $256\times$ range of provenance depth
while the unmemoized check grows linearly
(Appendix~\ref{app:refsem}).
\end{remark}

\begin{proposition}[\Ifive{} is auditable but not single-trace monitorable]
\label{prop:hyper}
\Ifive{} is a 2-safety hyperproperty~\cite{hyperprops}: its
violation is witnessed only by a pair of finite traces
$(J(q, d), J(q, d \oplus \delta))$. No black-box monitor observing
a single production trace can detect a violation, since
satisfaction is not a property of any individual
trace~\cite{hyperprops}; runtime checking requires either stored
comparison executions or a model of the system (gray-box
monitoring)~\cite{finkbeiner,stucki}. \Ifive{} is therefore
classified \emph{auditable}: falsifiable offline by exhibiting a
witness pair, testable by a differential harness, and checkable before
dispatch only by re-executing the decision (\S\ref{sec:gate}), never by
inspecting the production trace alone.
\end{proposition}

\begin{corollary}[The DEI decomposition]
\label{cor:decomp}
$\DEI = \Lift(\Ione \wedge \Itwo \wedge \Ithree \wedge \Ifour) \cap \Ifive$
(Def.~\ref{def:dei}) is the conjunction of an EM-enforceable
safety fragment (Prop.~\ref{prop:safety}) and a 2-safety fragment
that no single-trace monitor enforces (Prop.~\ref{prop:hyper}).
Any monitor claiming to enforce \DEI{} wholesale is unsound.
\end{corollary}

\begin{remark}[The boundary, stated plainly]
Full noninterference between untrusted channels and decisions is a
general hyperproperty~\cite{hyperprops,goguen} and inherits the same
non-monitorability; we neither claim nor need it. Complete
transformer-level noninterference is explicitly a non-goal even for
sink-authorization systems~\cite{cxi}, and there are principled
arguments that in-band immunity to injection may be
unattainable~\cite{abdelnabi}. Likewise, no monitor decides whether a
belief is \emph{true}; \DEI{} constrains the evidential and provenance
structure of decisions, not their semantic veracity. Concurrent work
characterizes generically what a deterministic pre-execution agent gate
can enforce: the safety policies whose good prefixes its state model
recognizes~\cite{ray2026enforced}. Prop.~\ref{prop:hyper} names a
specific, security-relevant property outside that class for
\emph{every} state model, and \S\ref{sec:gate} builds the
decision-local gate that audits exactly that property.
\end{remark}

\begin{proposition}[Mutual independence of the aggregator axioms]
\label{prop:indep}
There exist aggregators $S_{\neg\mathrm{A1}}$, $S_{\neg\mathrm{A2}}$,
$S_{\neg\mathrm{A3}}$, each violating exactly one axiom of
Def.~\ref{def:i3} while satisfying the other two: no axiom is
derivable from the rest, and dropping any one admits a named attack.
\end{proposition}
\noindent\emph{Proof:} the artifact, one explicit violating aggregator
per axiom with codomain and monotonicity checks.

\subsection{The Rationalization Gap}
\label{sec:guarantees}
The crosswalk of Appendix~\ref{app:crosswalk} maps each attack class of
\S\ref{sec:taxonomy} to the invariant it must violate.
Theorem~\ref{thm:crosswalk} shows the mapping is sound against the
enforceable fragment; Corollary~\ref{cor:ratgap} shows why that
fragment is nevertheless insufficient, and exactly which attacks
survive it.

\begin{theorem}[Crosswalk soundness]
\label{thm:crosswalk}
Under Assumptions~\ref{asm:mediation} and~\ref{asm:recorded}: for
every charting class $K$ with machine-fired primary $j(K)$ and every
trace $\sigma$ with $\mathrm{Mech}_K(\sigma)$, there is a finite prefix
$\sigma[..k]$ with $\mathit{viol}_{j(K)}(\sigma[..k])$, and
Algorithm~\ref{alg:monitor} blocks at or before $\epsilon_k$. Conversely
(anti-triviality), control/evaluation instances satisfy $\neg
\mathit{viol}_j$ at every prefix for all $j$ and are admitted.
\end{theorem}
\noindent\emph{Proof:} Appendix~\ref{app:t1c2} (finite case analysis
over the machine-fired rows, anti-triviality, and a delimitation of the
proved implication from its empirical support). The per-row predicates
$\mathrm{Mech}_K$ (Definition~\ref{def:mech}) are stated there.

\begin{corollary}[Rationalization gap: choice laundering is admission-undetectable]
\label{cor:ratgap}
Let $M$ be any sound monitor enforcing $\Ione \wedge \Itwo \wedge
\Ithree \wedge \Ifour$ at admission over traces where $\support(a)$ is
committed by an untrusted component (the model; a compromised model is
an adversarial component). Let $S_0$ be the support set of any honestly
admitted protected action of class $c_0$, and let $A_{\bindpol}(S_0)$
be the set of class-$c_0$ actions every bound argument of which is named
by some belief in $S_0$ at the label $\bindpol$ requires. Then for every
$a^{*} \in A_{\bindpol}(S_0)$ there exists an adversarial model
strategy---select $a^{*}$ from injected content; commit $S_0$---such that
on the resulting trace $\mathit{viol}_j$ holds at no prefix for any
$j \in \{1,2,3,4\}$ while the action's actual basis is the injection.
Hence (i)~the strategy is undetectable at admission by any sound
$\Ione$--$\Ifour$ monitor at \emph{any} provenance granularity; and
(ii)~an $\Ifive$-side counterfactual check is necessary, not optional:
no refinement of the enforceable fragment detects it.
\end{corollary}
\noindent\emph{Proof:} Appendix~\ref{app:t1c2} (existence via the
full-support commitment $S_0$ restricted to $A_{\bindpol}(S_0)$;
necessity via the $S^{\pm}$ construction of Prop.~\ref{prop:hyper}).

\begin{remark}[What \Ifour{} catches and what it leaves]
\label{rem:residual}
The corollary quantifies over $A_{\bindpol}(S_0)$, which is the
precise residual: an injection can still choose \emph{which} of the
trusted-derivable actions is taken, \emph{whether} it is taken, and
\emph{when}, but it can no longer choose a target that trusted evidence
does not name. The laundering probes executed in \S\ref{sec:rq2}
(delete file~13 for file~0; share with an attacker-chosen recipient) are
\emph{not} instances of the corollary's strategy: they violate \Ifour{}
and are blocked at admission once \Ifour{} is enforced. The gate's
incremental value must therefore be demonstrated on choice-laundering
probes, and \S\ref{sec:rq2} is designed accordingly. This is also why
\Ifour{} is a separate invariant rather than a clause of \Ithree: it
changes the adversary's residual strategy space, and the corollary
states that space exactly.
\end{remark}

\subsection{Falsifiability Protocol}
\label{sec:falsifiability}
Each invariant carries a refutation condition: the finite observation
that falsifies the claim that a given system enforces or satisfies it.
A single logged protected action admitted with an unattested or
sub-floor evidential root refutes ``the deployment enforces \Ione'';
a single silent label upgrade refutes \Itwo; a single admitted action
below threshold refutes \Ithree; a single admitted action with a bound
argument named by no qualifying cited belief refutes \Ifour; and one
$\delta \in \Dirrel$ with $\dist(J(q,d), J(q,d \oplus \delta)) >
\varepsilon_c$ refutes \Ifive. The first four are checkable by any
auditor holding the log, which is why the log is hash-chained and
receipt-bearing; the fifth requires the pre-registered differential
harness.

\paragraph*{Falsifying the definition itself (coverage)}
Let $\mathcal{A}$ be the attack taxonomy of \S\ref{sec:taxonomy} with
the per-instance success predicate of the evaluation harness. Say
\DEI{} \emph{covers} $\mathcal{A}$ on a corpus if every successful
attack instance violates at least one of \Ione--\Ifour{} on its trace
or \Ifive{} on its witness pair. A successful epistemic attack on an
invariant-satisfying execution falsifies coverage and shows the
definition incomplete: a reportable outcome, not an embarrassment. The
crosswalk (Appendix~\ref{app:crosswalk}) records, per attack class, the
invariant predicted to be violated, registered before the campaigns
that test it. An adaptive adversary must therefore either violate a
monitored invariant (and be blocked or logged) or exhibit a coverage
counterexample (and falsify the definition), the adaptive-evaluation
design the systematization of out-of-band defenses calls
for~\cite{oob}.

\begin{example}[Running example, III: the protocol applied]
\label{ex:protocol}
Walk the two traces of Example~\ref{ex:invariants} through the
protocol. Trace~1 (target file~13): if the deployment admits the action,
the admission log itself refutes ``the deployment enforces
\Ifour''---an implementation failure, repaired by fixing the monitor,
with no bearing on the definition. Trace~2 (target file~0, timing chosen
by the plant): every enforceable invariant holds on the log, so no
admission record can refute anything; the differential harness re-runs
the decision with the plant removed, the action changes, and the pair
refutes \Ifive. Had the action \emph{not} changed while the attack still
succeeded---say through an unbound free-text argument the perturbation
family never touches---that would be a coverage counterexample,
reported as evidence that the definition (here $\bindpol$ or
$\Dirrel$) is incomplete. The three outcomes are distinguishable from
the artifacts alone, which is what makes the protocol executable.
\end{example}

\paragraph*{Human validation of the crosswalk}
Because the crosswalk is a judgment artifact, two blinded domain
experts independently assigned invariants to each class from the
definitions alone, with an adjudicator resolving disagreements. On the
eight classes with a machine-fired primary they agreed exactly on 7/8
(Wilson 95\,\% CI 0.53--0.98) and 10/10 on the in-scope boundary; the
adjudicated consensus matches the invariant that actually fired in the
reference monitor on 6/8. The divergences are informative: both coders
read stale evidence as insufficient support (\Ithree) where the monitor
enforces a lapsed attestation (\Ione), and read prompt injection as a
provenance failure (\Ione) where the gateway blocks the unendorsed
label raise (\Itwo). The crosswalk records what the monitor
\emph{enforces}; the mapping was frozen before any campaign was scored
against it, and the instrument and scorer are released.

\subsection{Scope and Limitations}
\label{sec:limitations}
Appendix~\ref{app:refsem} distinguishes proved, implemented, measured,
development, and uninformative evidence, so that no deferred study or
failed susceptibility check is mistaken for a positive result. Beyond
that: (1)~all guarantees are relative to the epistemic TCB; compromising
the labeler, the channel authentication it relies on, the verifier
keys, the recorder, or the aggregator voids them, and \S\ref{sec:tcb}
quantifies that base. (2)~Rationalization is a faithfulness gap, not a
granularity one. \Ione--\Ifour{} constrain the \emph{committed} support
set and the action's bound arguments, but the support set is committed
by an untrusted model, so a compromised model can select an action from
adversarial content and cite legitimate evidence for it. \Ifour{}
narrows the residual to choice laundering under trusted-derivable
arguments (Cor.~\ref{cor:ratgap}) without closing it; sub-document
laundering is a resolution limit that finer provenance
tightens~\cite{argus,cxi}, and~\cite{romanchuk2026semantic} states the
architectural version of the gap as \emph{semantic laundering}, for
which \S\ref{sec:gate} builds the runtime check. (3)~$\bindpol$ and
$\ell_{\bindpol}$ are policy with a utility cost: a target discovered
from a $\mathrm{MEDIUM}$ tool result cannot bind at $\mathrm{TRUSTED}$,
so $\bindpol$'s tightness is bounded by the attestation topology
exactly as $\tc$ is (\S\ref{sec:evalsetup}). (4)~$\Dirrel$ is itself a judgment
and the hard core of \Ifive~\cite{qu}; we pre-register the generator
suite rather than choose perturbations post hoc. (5)~\Ithree's
aggregator encodes an epistemology; we require only A1--A3.
(6)~\DEI{} constrains structure, not truth: attested, well-supported,
invariant-satisfying decisions can still be wrong.


\section{CZTA: A Reference Monitor for the Enforceable Fragment}
\label{sec:czta}
Proposition~\ref{prop:safety} says \Ione--\Ifour{} \emph{can} be
enforced by execution monitoring; this section shows how they are
enforced, at measured cost, on a real deployment. CZTA is a zero-trust
gateway interposed between an agent (a production model) and a live
MCP server, with Keycloak as policy decision point; its admission path
is Algorithm~\ref{alg:monitor} realized at the gateway.

\subsection{Architecture and Workflow}
\label{sec:arch}
Figure~\ref{fig:czta} shows the evaluated path. Every ingress channel
passes a \emph{policy enforcement point} (PEP) that authenticates the
channel, stamps $\lambda(\mathit{ch})$, digests the content, issues the
attestation receipt of Def.~\ref{def:evidence}, and appends the atom to
the ledger; the model runtime sees only labeled, receipt-bearing atoms.
Beliefs are created through the update API (meet rule, $\deps$ and
$\vals$ recorded, root certificate memoized); label raises exist only
as endorsement-API calls checked against $\pi$ and logged with their
principal. A protected tool call is committed as a structured
\textsc{support}/\textsc{action}/\textsc{args} step; the gateway
evaluates \Ione{} (receipts), \Itwo{} (per update), \Ithree{}
(aggregator), and \Ifour{} (bound arguments against the cited beliefs'
value sets), then dispatches the MCP request under a Keycloak-issued,
audience-bound token or suppresses it and logs the violated predicate.
Irreversible classes pass the audit gate of \S\ref{sec:gate} between
admission and dispatch, and a reasoning-lineage attestation (RLA) signs
the decision state at each enforcement point so that any holder of the
ledger can replay the decision.

\begin{figure}[t]
\centering
\resizebox{\columnwidth}{!}{%
\begin{tikzpicture}[font=\scriptsize,
  pipe/.style={draw, rounded corners=2pt, align=center, minimum height=7.5mm, inner sep=2.5pt, fill=white},
  pep/.style={pipe, draw=accent, thick, fill=accentlight},
  sup/.style={draw, rounded corners=2pt, align=center, inner sep=2.5pt, fill=black!4},
  arr/.style={-{Stealth[length=2.2mm]}, thick},
  darr/.style={-{Stealth[length=1.8mm]}, dashed, ledger}]
\node[pipe, text width=1.75cm] (src) at (0,0) {Sources, tools, memory};
\node[pep, text width=2.1cm] (ipep) at (2.85,0) {Ingress PEP: $\lambda$, digest, receipt};
\node[pipe, text width=1.9cm] (agent) at (5.75,0) {Model runtime (untrusted)};
\node[sup, text width=3.0cm] (api) at (1.6,-1.55) {Belief API: update (meet, $\vals$, certificate); endorse ($\pi$)};
\node[pep, text width=2.45cm] (gw) at (5.75,-1.55) {\textbf{Admission}: \Ione--\Ifour{} + gate (\S\ref{sec:gate})};
\node[pipe, text width=1.4cm] (mcp) at (8.1,-1.55) {MCP server};
\node[sup, text width=2.45cm] (kc) at (5.75,-2.95) {Keycloak PDP: identity, audience, tool scope};
\draw[arr] (src) -- (ipep); \draw[arr] (ipep) -- (agent);
\draw[arr] (agent.south) -- node[right, font=\tiny] {commit} (gw.north);
\draw[arr] (gw) -- node[above, font=\tiny] {token} (mcp);
\draw[-{Stealth[length=1.8mm]}, accent] (agent.south west) -- node[above, sloped, font=\tiny] {update / endorse} (api.north east);
\draw[-{Stealth[length=1.8mm]}, accent] (kc.north) -- (gw.south);
\node[draw=ledger, fill=ledger!10, rounded corners=2pt, align=center, inner sep=3pt, text width=8.3cm] (rla) at (4.1,-4.25)
 {\textbf{Append-only, hash-chained ledger + RLA signatures}: atoms, receipts, updates, endorsements, verdicts (Ed25519-sealed segments)};
\draw[darr] (ipep.south) -- (rla.north -| ipep);
\draw[darr] (api.south) -- (rla.north -| api);
\draw[darr] (kc.south) -- (rla.north -| kc);
\draw[darr] (gw.east) -- ++(0.35,0) |- ($(rla.north -| mcp) + (0,0)$);
\end{tikzpicture}}
\caption{The evaluated CZTA path. Blue boxes are enforcement points
inside the epistemic TCB; the model runtime and the MCP server are
untrusted. Every event that can affect a protected action crosses an
enforcement point (Assumption~\ref{asm:mediation}) and is ledgered.}
\label{fig:czta}
\end{figure}

\begin{example}[Running example, IV: the workflow end to end]
\label{ex:workflow}
(1)~The user turn arrives on the authenticated user channel; the ingress
PEP stamps $\mathrm{TRUSTED}$, digests it, issues receipt $r_u$, and
appends $e_u$. (2)~The drive listing and the document body carrying the plant
arrive on the tool channel and are stamped $\mathrm{MEDIUM}$ with
receipts $r_l, r_p$; nothing about the plant is detected here, and
nothing needs to be. (3)~The model creates $b_u$, $b_l$, $b_p$ through
the update API, each with the meet of its inputs, a root certificate
over its receipt, and its recorded values.
(4)~The compromised model commits \textsc{action} delete, \textsc{args}
$\{\mathsf{target} \mapsto \text{file~13}\}$, \textsc{support}
$\{b_u\}$. (5)~Admission: \Ione{} passes ($\cert(b_u) = \mathrm{ok}$,
window valid); \Itwo{} was never at issue; \Ithree{} passes
($S(\{b_u\}, \mathrm{TRUSTED}) \ge \tc$); \Ifour{} fails
(file~13 $\notin \vals(b_u)$). The gateway suppresses the call, logs
$\mathit{viol}_4$ with the offending slot, and the ledger entry is the
refutation record of \S\ref{sec:falsifiability}. (6)~Had the model
committed file~0 under a timing chosen by the plant, admission would
pass and the gate's re-commitment with $e_p$ struck would select a
different action, so the gateway would quarantine instead of dispatch.
\end{example}

\subsection{Instantiating the Invariants}
Table~\ref{tab:instantiation} maps each enforceable invariant to its
CZTA mechanism, the fields logged and attested to make post-hoc
refutation possible, and the enforcement point. Two engineering choices
carry the formal weight. First, \emph{root certificates}: when a belief
is created, the gateway verifies its evidential roots once and issues a
hash-chained receipt; at action time, \Ione{} reduces to
$O(|\support(a)|)$ receipt verifications, the operation whose
sign/verify latency Table~\ref{tab:overhead} reports. Second,
\emph{endorsement as an API}: label upgrades exist only as explicit,
logged gateway calls citing a policy rule, which is what makes \Itwo's
``no silent upgrade'' condition checkable by inspection of the log.

\begin{table}[t]
\centering
\caption{Instantiation of the enforceable invariants in CZTA.
Measured costs appear in Table~\ref{tab:overhead}.}
\label{tab:instantiation}
\scriptsize
\setlength{\tabcolsep}{2.5pt}
\begin{tabularx}{\columnwidth}{@{}c >{\raggedright\arraybackslash}p{2.3cm} >{\raggedright\arraybackslash}X >{\raggedright\arraybackslash}p{1.45cm}@{}}
\toprule
Inv. & CZTA mechanism & Logged/attested fields & Point \\
\midrule
\Ione & Root certificates over attested evidence; admission check on protected tool calls & belief id; root set; receipt hash chain; attestation ids; validity windows & Gateway, pre-invocation \\
\Itwo & Label propagation by lattice meet; endorsement API with policy-rule citation & update event; input labels; result label; endorsement rule id; principal & Gateway, per update \\
\Ithree & Aggregator service scoring qualifying evidence against per-class thresholds & support set; qualifying-evidence ids; score; $\tc$; verdict & Gateway, pre-invocation \\
\Ifour & Bound-slot check of $\args(a)$ against $\vals$ of cited beliefs at $\ell_{\bindpol}$ (list-valued slots element-wise); the citation contract of the pilot grid (\S\ref{sec:evalsetup}) is its weaker precursor & slot; value; citing belief id; label; verdict & Gateway, pre-invocation \\
\Ifive & Not enforced (Prop.~\ref{prop:hyper}); audited offline and, for irreversible classes, by the gate & paired traces; $\delta$ operator id; decision distance & Offline harness; gate \\
\bottomrule
\end{tabularx}
\end{table}

\subsection{Epistemic TCB Inventory and Security Analysis}
\label{sec:tcb}
\label{sec:security}
The TCB of \S\ref{sec:tcbdef} concretely comprises the ingress labeler
and the channel authentication it relies on, the attestation verifier
and its key store, the append-only hash-chained provenance log with its
value recorder, the aggregator service, and the Keycloak policy engine.
The trusted implementation is the deposited replacement
(\texttt{v9-replacement-2}), not the unavailable original v9 package:
the executable reference semantics of Algorithm~\ref{alg:monitor}
(\Ione--\Ifour{} with receipt-based root and window checks), a strict
RS256/JWKS verifier, an MCP \texttt{tools/call} decision endpoint, and a
hash-chained Ed25519-sealed ledger. Its RLA module and all other v8
modules are preserved byte-for-byte for audit; only modules actually
imported are asserted on the measured path. The Keycloak realm stamps
the claims the gateway inspects, and \texttt{/env} publishes the hash of
every trusted module imported, so each measurement is bound to its code.
The retired v8 scenario-specific policy pipeline is not used for any
number here. Benchmark, scenario, and other retired modules remain
outside the trust path. Everything else---the model, MCP servers and
their tool descriptions, retrieved content, and memory---is untrusted.

\paragraph*{What an adversary can do to the gateway}
Interposing a monitor creates a surface of its own, and the formal
guarantees say nothing about it. We enumerate the attack vectors the
design admits and how each is closed or conceded.
\emph{Channel spoofing} (adversarial content over a high channel) is
the one attack on $\lambda$, excluded by Assumption~\ref{asm:recorded};
in CZTA it reduces to forging a session token or a gateway receipt, so
its security is that of RS256 verification against JWKS and of Ed25519
receipt signatures. \emph{Endorsement abuse} requires a rule
$\rho \in \pi$ and an authorized principal; the deployed $\pi$
authorizes no post-ingress raise (\S\ref{sec:evalsetup}), and every
raise is logged with its principal. \emph{Receipt replay} is closed by
binding each receipt to a content digest checked at verification.
\emph{Recorder evasion} (a target in a form the extractor misses) fails
closed under \Ifour: an unextracted value is unbound and the action is
blocked, a utility cost rather than a security cost. \emph{Mediation bypass} (a side channel from untrusted
content to action that avoids a PEP) is Assumption~\ref{asm:mediation};
in the evaluated path the MCP transport is the only route to the
server, and a deployment that exposes a second route voids the
guarantee. \emph{Denial of service through the gate} is bounded by two
re-commitments per action and by never gating reversible classes.
Re-commitments are never dispatched and the counterfactual plane is
assembled inside the TCB from ledgered atoms, so re-execution has no
side effects; a model that commits identically with and without the
plant defeats the gate, a coverage question for the falsification
protocol rather than a silent failure (\S\ref{sec:gate}). \emph{TCB compromise}
(verifier keys, ledger key, aggregator) voids all guarantees; the
hash-chained ledger makes tampering detectable after the fact but not
preventable, which is the concession every reference monitor makes.
\emph{Rug-pull} (a tool whose description is mutated after approval)
is the one documented MCP pattern the evaluated gateway detects without
preventing, as reported below.

\subsection{Authorization-Layer Checks}
\label{sec:redteam}
The epistemic evaluation lives in \S\ref{sec:eval}; this subsection
reports a complementary check that the gateway's \emph{authorization}
layer behaves as specified, since \DEI{} enforcement composes with,
and does not replace, the sink-layer defenses of CXI-class
systems~\cite{cxi,camel}. We executed four documented MCP attack
patterns (tool poisoning, token passthrough, confused deputy, rug-pull)
plus a postmark-style output-exfiltration boundary case against the
live gateway with defenses enabled. These are MCP/OAuth binding checks
(RFC~8693/9449/9396), not tests of the coverage claim, and they were
constructed by the authors during development; they are reported
because a gateway monitor must get them right and because one is an
honest negative. Tool poisoning, token passthrough, and the confused
deputy are blocked at the authorization layer (\texttt{mcp-tool-not-exposed}
403; \texttt{mcp-token-passthrough-forbidden} 403;
\texttt{InvalidAudienceError} 401); the output-exfiltration case is,
as pre-registered, detected but not prevented (accountability via the
signed RLA output digest). The rug-pull executed (HTTP~200): a
post-approval tool mutation is a \emph{temporal} \Itwo{} violation
(integrity holds at approval, breaks later) while the admission check
consults only approval-time state. Def.~\ref{def:evidence} already
attaches a validity window to every attested atom and requires
re-checking at use, so the fix is a tool-state attestation recomputed at
call time rather than new machinery; we report the measured behaviour
and leave the implementation to future work.

\subsection{Measured Overhead}
\label{sec:overhead}
We measure the cost of interposing the unmodified reference
monitor (source hashes published at \texttt{/env}) on a real
authorization path: real Keycloak RS256 tokens, real JWKS
validation, and real MCP JSON-RPC transport. Per the
pre-registered measurement protocol, arrivals are open-loop
Poisson at $\lambda \in \{5, 25, 100\}$\,req/s (latency taken from
the intended send time, so coordinated omission cannot understate
the tail), each rate held 60\,s after a discarded 20\,s warmup,
over three independent repetitions per arm with the process
restarted between them. Three arms share the identical request path:
A0 with the monitor removed, A1 enforcing \Ione--\Ithree, and A2
enforcing \Ione--\Ifour; the overhead is A2 minus A0, and the cost of
the \Ifour{} slot check is A2 minus A1. To isolate the monitor, the
upstream is a minimal tool stub with no model inference in the path, so
the relative overhead reported here is an \emph{upper bound}: the fixed
monitor cost amortizes against real tool and model latency in any
production deployment. The protocol requires zero errored requests in
every cell; all 27 cells passed with zero errored requests.

\begin{table}[t]
\centering
\caption{Authorization-path latency (ms) on the live stack (Keycloak
RS256 + JWKS + MCP \texttt{tools/call}; stub upstream, so relative
overhead is an upper bound). Open-loop Poisson arrivals; cells are mean
p50/p95 across three independently restarted repetitions, with zero
errored requests required; values flow from the harness output.}
\label{tab:overhead}
\scriptsize
\begin{threeparttable}
\setlength{\tabcolsep}{6pt}
\begin{tabular}{@{}l c c c@{}}
\toprule
& $\lambda{=}5$ & $\lambda{=}25$ & $\lambda{=}100$ req/s \\
\midrule
A0 (no monitor) p50/p95 & \RqFiveAZeroFiveCell & \RqFiveAZeroTwentyFiveCell & \RqFiveAZeroHundredCell \\
A1 (\Ione--\Ithree) p50/p95 & \RqFiveAOneFiveCell & \RqFiveAOneTwentyFiveCell & \RqFiveAOneHundredCell \\
A2 (\Ione--\Ifour) p50/p95 & \RqFiveATwoFiveCell & \RqFiveATwoTwentyFiveCell & \RqFiveATwoHundredCell \\
A2$-$A0 p50 (\%) & \RqFiveOverheadFive & \RqFiveOverheadTwentyFive & \RqFiveOverheadHundred \\
\bottomrule
\end{tabular}
\begin{tablenotes}[flushleft]\scriptsize
\item RLA sign / verify, mean p50/p95 across
\RqFiveRlaIterations{} iterations: \RqFiveRlaSignCell{} /
\RqFiveRlaVerifyCell{} ms; zero verification errors.
\end{tablenotes}
\end{threeparttable}
\end{table}

Table~\ref{tab:overhead} reports the result: A2$-$A0 p50 overhead is
\RqFiveOverheadFive\%, \RqFiveOverheadTwentyFive\%, and
\RqFiveOverheadHundred\% at 5, 25, and 100 req/s; A2$-$A1 is the
corresponding \Ifour{} increment. These are the constant factors behind
Proposition~\ref{prop:safety}; the reference implementation confirms
their asymptotic behaviour independently (Appendix~\ref{app:refsem}),
and every restarted repetition is retained in the deposited result file.


\section{The Counterfactual Audit Gate}
\label{sec:gate}
Corollary~\ref{cor:ratgap} prices the enforceable fragment: a laundered
choice is admitted, and no admission-time check can tell. The
impossibility is about single traces, not decisions, and a decision can
be re-run. This section deploys the auditable fragment as a
decision-local gate over irreversible classes: an action whose
commitment does not survive its own counterfactual is quarantined.

\subsection{Design}
The gate sits between admission and dispatch. For every admitted
action $a$ with cited support $S_c$ whose class is in the
pre-registered irreversible set $D = \{\mathrm{delete}, \mathrm{share},
\mathrm{send}\}$, the gateway issues re-commitments on perturbed
evidence planes and compares the result with $a$. If any re-commitment
returns an action $a' \neq a$, or leaves $a$ without qualifying
support, $a$ is quarantined and escalated to the human checkpoint that
already serves sub-$\tc$ admissions. Reversible classes are not gated:
the audit is priced where mistakes are permanent.
Algorithm~\ref{alg:gate} states the semantics; the gate adds no TCB,
since its inputs (admission record, labeler, ledgered atoms, commitment
protocol) are already inside the TCB of \S\ref{sec:tcb}.

\paragraph*{Which content is removed, and how the gate knows}
The gate never identifies malicious content. It removes evidence by
\emph{provenance class}, which is computable inside the TCB from the
ledger alone, and it asks whether the action stands on what remains.
Three perturbations are defined:
\begin{itemize}[leftmargin=1.15em,itemsep=1pt]
\item $\delta_1$, \emph{plant removal}: the evidence plane in which the
injected content never entered. This is an \emph{oracle} perturbation,
available only when the harness knows the injection vector; it is used
in the evaluation as the upper bound on what any deployed removal can
achieve.
\item $\delta_2$, \emph{cited-sub-required removal}: strike every
$x \in S_c$ with $\ell(x) \sqsubset \lreq(c(a))$ and re-commit against
what remains. (The version of this perturbation evaluated in
\S\ref{sec:rq2} struck atoms below $\lmin$ rather than below $\lreq$,
which in the deployed policy never fires; the $\lreq$ form is the
registered deployable version.)
\item $\delta_3$, \emph{uncited removal}: restrict the plane to the cited
support $S_c$ plus the trusted state and re-commit. If the action
depended on uncited content, it changes.
\end{itemize}
$\delta_2$ and $\delta_3$ are oracle-free. Together they cover both
ways a plant can influence a committed action: if the plant is uncited,
$\delta_3$ removes it; if it is cited, it sits below $\lreq$ for any
protected class whose required label exceeds the tool channel's, and
$\delta_2$ removes it. A plant cited \emph{at} $\lreq$ is channel
spoofing (\S\ref{sec:security}), not laundering. The registered
deployment \emph{enforces} $\delta_3$ and \emph{observes} $\delta_2$:
at atom granularity $\delta_2$ cannot tell a commitment that depends on
needed $\mathrm{MEDIUM}$ data from one that depends on a plant sharing
the same atom, so its flips are logged and reported as the gate's
resolution limit rather than enforced.

\paragraph*{Mixed content, stochasticity, and divergent state}
Removal is at atom granularity: a plant that shares an atom with the
data the decision needs (a planted claim inside the file listing the
user asked to be read) is removed together with that data under
$\delta_2$, and an honest commitment changes for the same reason a
laundered one does; a plant in an atom of its own is removed by
$\delta_3$ alone. The probe families of \S\ref{sec:rq2} are built in
both placements so that this limit is measured, not assumed; span-level
provenance~\cite{argus} is what would lift it. Re-commitments are issued at the
endpoint's deterministic setting (temperature~0, fixed seed where
supported); otherwise the registered rule is $k$ re-commitments with a
majority verdict, and the gate's inferences assume deterministic $J$.
Tool state cannot diverge between runs because the counterfactual
plane is assembled from ledgered atoms, never by re-querying tools. A
re-commitment that differs for a reason unrelated to the removed
content is quarantined anyway: the conservative default, priced by the
same false-quarantine rate.

\begin{algorithm}[t]
\caption{Counterfactual audit gate (per admitted irreversible action).}
\label{alg:gate}
\footnotesize
\begin{algorithmic}[1]
\Require admitted action $a$, cited support $S_c$, $c(a) \in D$; perturbation set $\Delta_{\mathrm{gate}}$ ($\{\delta_1\}$ in the oracle evaluation; $\{\delta_2, \delta_3\}$ deployed)
\For{$\delta \in \Delta_{\mathrm{gate}}$}
  \State $a' \gets \textsc{Recommit}(\mathrm{plane}(\mathrm{ctx}) \setminus \delta)$ \Comment{plane assembled from ledgered atoms}
  \If{$a' \neq a$ \textbf{ or } $\neg\textsc{Qualifies}(a)$ under $\delta$}
    \State \Return \textsc{Quarantine}$(a)$; escalate to human checkpoint
  \EndIf
\EndFor
\State \Return \textsc{Execute}$(a)$
\end{algorithmic}
\end{algorithm}

\subsection{Soundness and the Approximation Gap}
\begin{proposition}[Gate soundness, oracle case]
\label{prop:gatesound}
Under Assumptions~\ref{asm:mediation}--\ref{asm:recorded},
deterministic $J$, and oracle plant removal ($\delta_1 \in
\DeltaIrrel(q)$): if the gate quarantines an admitted action $a$ on
the plant-removal plane, the executions on $(q,d)$ and $(q, d \oplus
\delta_1)$ are an \Ifive{} witness pair at the level of
policy-admissible decisions---$J_{\Pi}(x)=J(x)$ if $J(x)$ qualifies
under admission and $\bot$ otherwise---with $\dist=1>\varepsilon$
under the discrete metric extended by $\dist(a,\bot)=1$.
\end{proposition}

\begin{proposition}[Approximation gap, deployable perturbations]
\label{prop:approxgap}
Let the deployed gate approximate the plant-removal plane by an
enforced set $E \subseteq \{\delta_2, \delta_3\}$, and over the
registered probe distribution let $\mathit{fp}$ be the rate at which an
honest commitment changes under some $\delta \in E$ and $\mathit{miss}$
the rate at which a laundered commitment survives every $\delta \in E$.
For deterministic $J$ the false-quarantine rate is at most $\mathit{fp}$
and the missed-quarantine rate at most $\mathit{miss}$; each
$\delta \in E$ is an exact recomputation inside the TCB and contributes
no gap of its own.
\end{proposition}
\noindent\emph{Proofs:} the artifact (faithful re-commitment; both
quarantine disjuncts discharged at the policy-admissible level;
qualification-preservation lemma and two event inclusions).
$\mathit{fp}$ and $\mathit{miss}$ are measured in \S\ref{sec:rq2}:
$\delta_3$ has $\mathit{fp}$ \RqTwoHonestDeltaThreeRate{} on actionable
honest evaluations and $\mathit{miss}$ \RqTwoDeltaThreeMissRate{} against
the oracle comparator.

\subsection{Relation to Counterfactual Defenses}
Two concurrent efforts bound the same phenomenon from opposite sides:
provenance-sensitivity auditing measures latent dependence offline and
is disclaimed as a guardrail~\cite{liao2026provenance}; causal tool
filtering minimizes which tools are visible~\cite{babu2026gistcmtf}.
Rationalization operates where filtering cannot help (the correct tool,
legitimately cited), and the gate acts where offline diagnosis cannot
(between admission and dispatch). Counterfactual re-execution defenses
perturb a different plane: MELON masks the user prompt~\cite{melon},
AgentSentry purifies the untrusted context~\cite{agentsentry},
CausalArmor ablates untrusted spans~\cite{causalarmor}; each asks
whether the action depends on content it should ignore, whereas the gate
asks whether the action stands on what it \emph{cites} and quarantines
irreversible classes before dispatch. Intervening at the (action, args)
decision point rather than in token or latent space~\cite{pona2025acf}
is what makes the audit cheap enough to run per action.


\begin{proposition}[Composition with sink authorization]
\label{prop:compose}
For $P_{\DEI}=\Ione\wedge\Itwo\wedge\Ithree\wedge\Ifour$ and any
EM-enforceable safety property $P_{\mathrm{sink}}$ over the extended
alphabet (a sink-authorization discipline~\cite{cxi}), the conjunction
$P_{\DEI}{\uparrow}\wedge P_{\mathrm{sink}}$ is EM-enforceable by a
product truncation automaton: every executed protected action is both
evidentially sound and manifest-authorized, and the conjunction
strictly refines each conjunct.
\end{proposition}
\noindent\emph{Proof:} the artifact (padding preserves safety; product
truncation automaton; commutativity; strictness witnesses in both
directions). The practical reading is that \DEI{} gates what may count
as a reason and CXI-class systems gate what a proposal may do; neither
subsumes the other.


\section{Evaluation}
\label{sec:eval}
Five questions. \textbf{RQ1}: does enforcing the safety fragment change
attack outcomes on a live model that follows the injection, and how
does it compare with simpler checks and an existing defense?
\textbf{RQ2}: does the gate add security where admission is provably
blind? \textbf{RQ3}: what does enforcement cost in benign utility, at
the power needed to resolve it? \textbf{RQ4}: do the answers hold
across models, suites, and attack families? \textbf{RQ5}: what does it
cost in latency? Every number flows from the episode logs through a
recompute pipeline that emits the macros used in the text.

\subsection{Setup and Registered Design}
\label{sec:evalsetup}
\paragraph*{Subjects, suites, attacks}
The primary subject is a frontier live endpoint, gemini-2.5-flash, on
the AgentDojo \emph{workspace} suite (email, calendar, cloud drive;
benchmark v1.2.x, package \texttt{agentdojo}~0.1.35) under the canonical
\texttt{important\_instructions} attack family, through a checkpointed
harness derived from the benchmark's own
\texttt{run\_task\_with\_pipeline} path. The pilot grid ran at the
endpoint's default sampling temperature (the pilot harness did not set
one); the registered harness sends temperature~0.7 and the derived seed
explicitly on every request. A dated proprietary snapshot
(\texttt{gpt-4o-2024-05-13}) is the positive calibration control (5/5
at $t{=}0.0$, 4/5 at $t{=}0.7$ on the known-susceptible
\texttt{workspace/user\_task\_0} injections), ruling out an inert
harness. A 0.5B open-weight driver (Qwen2.5-0.5B-Instruct, greedy)
developed the harness and the rationalization probes of
\S\ref{sec:rq2}; it completes no workspace task with or without
enforcement (0/16), so no utility or attack rate is claimed from it.
RQ4 adds a second frontier endpoint, a post-fix re-run of the
open-weight qwen3 family, and the benchmark's \texttt{tool\_knowledge}
attack family, the first stage of the cascade
of~\cite{bhagwatkar2025firewalls}; all three replications completed.
AgentDyn adds the completed dynamic-suite replication reported in
\S\ref{sec:rq4}; suite-generality remains bounded by its task and model
coverage~\cite{agentdyn}.

\paragraph*{Deployed policy objects}
$\lambda$ is as in Definition~\ref{def:policy}: the user turn and the
system prompt $\mathrm{TRUSTED}$, every tool result $\mathrm{MEDIUM}$,
model text no evidence at all. Support of an action is the set of atoms
naming one of its argument values plus the user atoms that express its
class (the registered directive lexicon); $S = \sum w / (1 + \sum w)$
over qualifying support with $\tc = 0.45$ at $\lreq = \mathrm{TRUSTED}$
for all protected classes, fixed before any attack run. This $S$
satisfies A1--A3 of Definition~\ref{def:i3} (sum over attested atoms
at or above $\lreq$; $x/(1+x)$ increasing with $w \ge 0$; weights
fixed in the TCB from the recorded prefix). $\pi$ authorizes no
post-ingress raise. $\bindpol$ binds targets that leave the trust
boundary (recipients, payees, users, channels, URLs, amounts) at
$\mathrm{TRUSTED}$ and internal handles that only tools can enumerate
(file, e-mail, and event ids; hotel and vendor names) at
$\mathrm{MEDIUM}$; free text and dates are unbound.
$D = \{\mathrm{delete}, \mathrm{share}, \mathrm{send}, \mathrm{book}\}$.
Two-root corroboration is not deployable on a single first-party
channel, so policy tightness is bounded by the attestation topology.
The offline sweep gives a flat region at $\tc\in\{0.30,0.45,0.50\}$:
\SpSweepTauFortyFiveSecAdmit/\SpSweepTauFortyFiveSecN{} security and
\SpSweepTauFortyFiveBenignAdmit/\SpSweepTauFortyFiveBenignN{} benign first protected
decisions admit. At $\tc=0.60$ and 0.67, both rates fall to zero, placing
the cliff between 0.50 and 0.60. Requiring every bound slot to be
TRUSTED reduces security admissions to \SpSweepEllAllTrustedSecAdmit/
\SpSweepEllAllTrustedSecN{} but leaves benign admissions unchanged at
\SpSweepEllAllTrustedBenignAdmit/\SpSweepEllAllTrustedBenignN.

\paragraph*{Development versus test}
Every case any development run touched is excluded from the held-out
pools and reported as development evidence: the four rationalization
probes, the five MCP authorization patterns of \S\ref{sec:redteam},
the 22-task small-model workload on which $\tau_c$ was fixed, the
post-fix open-weight run, and the 40-case/24-case pilot grid of
\S\ref{sec:rq1}, whose defended arms ran a \emph{protected-tool
filter}: the pilot harness never wired support extraction, so every
state-changing call was dropped regardless of evidence, and the one
target outside the filter's list (calendar-event creation) accounts for
all four defended-arm breaches. The held-out case sets (150 security
cases and the 56 benign tasks no development run touched, across the
four suites, drawn by a seeded procedure), the 20 choice-laundering
probes with their honest controls, the policy objects, the baseline and
gate configurations, and the utility minimal important difference are
deposited hash-sealed and publicly timestamped before the RQ1--RQ4 runs:
registration v2 (archive SHA-256 \texttt{ee88c72b\ldots c68}; manifest
\texttt{cbddbd33\ldots f0}; OpenTimestamps block 974214,
2026-11-02T05:06:44Z; the deleted original DOI is being re-registered at
final upload as DOI \url{https://doi.org/10.5281/zenodo.24075933}) and the replacement-2
testbed v4 (DOI \url{https://doi.org/10.5281/zenodo.23914405};
archive SHA-256 \texttt{d97f1369\ldots f3b}; manifest
\texttt{c67de998\ldots 071}; OpenTimestamps block 974389,
2026-11-03T13:03:48Z). The campaign hotfix is DOI
\url{https://doi.org/10.5281/zenodo.23917748}, manifest
\texttt{dfc691f8\ldots 13b}, OpenTimestamps block 974430,
2026-11-03T19:47:14Z.
An intermediate replacement-1 implementation and its results were
excluded because they lacked independent restarts and MCP transport. The $\tau_c$ workload touched only benign
traces; security cases on those workspace user tasks remain in the
pool and form a registered descriptive stratum beside the $\tau_c$
sweep, rather than being silently presented as untouched.

\paragraph*{Baselines}
Three baselines run under the identical harness and case sets:
\textbf{B1}, \emph{argument--instruction matching}, every target named
in a protected call must appear in the user's instruction (the simpler
check a decision-layer defense must beat); \textbf{B2}, a tool-filter
policy in the CaMeL lineage~\cite{camel,oob} (no tool-sourced value may
reach a protected sink unless the user turn names it); \textbf{B3},
\Ifour{} alone. Every arm's verdict is logged on every protected call
as a shadow; all seven arms completed the registered campaign.

\paragraph*{Statistics}
Wilson 95\,\% intervals on every rate; Fisher's exact test on the one
pre-registered primary comparison per grid, with risk difference (RD)
and risk ratio (RR) and their intervals beside every $p$-value; the
rule-of-three bound beside any zero-event claim; secondary comparisons
Holm-corrected. Five harness defects (among them the unset
temperature) are documented in the artifact; degenerate episodes (null
model output) are reported and excluded only in a sensitivity analysis.
The reference semantics' own coverage check (10/10 classes, 3/3 honest
baselines) is Appendix~\ref{app:refsem}.

\subsection{RQ1: Security on a Live Model}
\label{sec:rq1}
\paragraph*{The subject is susceptible}
gemini-2.5-flash follows the canonical injection in
\SecNdHit/\SecNdN{} no-defense pilot episodes (ASR \SecNdASR, Wilson
95\,\% CI $[\SecNdCIlo, \SecNdCIhi]$); a disjoint 30-case sweep
corroborates this ($23/60 = 0.383$ $[0.271, 0.510]$). The pilot
covers three workspace user tasks, a scope the held-out set removes.

\paragraph*{A protected-tool filter already collapses this family (pilot)}
On this subject (\SecNdN{} independent cases per arm), a filter that
drops every call to a state-changing tool reduces measured attack
success from \SecNdASR{} (\SecNdHit/\SecNdN) to \SecGateASR{}
(\SecGateHit/\SecGateN, Wilson $[\SecGateCIlo, \SecGateCIhi]$) in each of
two replicates: Fisher exact $p{=}\GateVsNdFisherP$, RD $= \HeadlineRD$
$[\HeadlineRDlo, \HeadlineRDhi]$, RR $= \HeadlineRR$
$[\HeadlineRRlo, \HeadlineRRhi]$ (Table~\ref{tab:livemodel}). The filter
verifiably engages (\SecMonFilt--\SecGateFilt{} protected calls dropped
per arm); every surviving breach is the one template whose target,
calendar-event creation, was not on its list, and excluding the four
degenerate episodes changes nothing (2/37, $p{=}6.4\times10^{-5}$).
This is B2 in its strictest form, and it sets the bar: admission must
match that security at lower benign cost. Security is measured across
seven arms on the held-out set; the powered RQ3 crossover decides the
comparative benign-cost bar.
The registered rows show that admission reduces ASR from
\SpSecNoDefenseASR{} to \SpSecDeiASR{} (Fisher
$p{=}\SpSecDeiFisherP$), but it does not beat B1: both are
\SpSecBOneASR, with Holm $p{=}\SpSecDeiVsBOneHolmP$. B2strict is lower
still at \SpSecBTwoStrictASR. RQ1 is therefore a result about
\Ifour{} rather than evidence of full-DEI superiority; the gate does not
change admission's ASR. The powered RQ3 contrasts, not the 56-task
descriptive arm, decide comparative benign cost.

\begin{table}[t]
\centering
\caption{RQ1 and RQ3. Pilot (development evidence): gemini-2.5-flash,
workspace, $n{=}\SecNdN$ cases per security arm, $n{=}\UtilNdN$ per
benign arm; the defended arms are two replicates of a protected-tool
filter. Registered campaign: held-out cases, seven arms.}
\label{tab:livemodel}
\fontsize{6.5}{7.5}\selectfont
\setlength{\tabcolsep}{0pt}
\begin{tabular}{@{}lccc@{}}
\toprule
Arm & ASR & 95\,\% CI & Filtered \\
\midrule
\multicolumn{4}{@{}l}{\emph{Pilot grid}}\\
No defense & $\SecNdASR$ ($\SecNdHit/\SecNdN$) & $[\SecNdCIlo, \SecNdCIhi]$ & $\SecNdFilt$ \\
Protected-tool filter, replicate 1 & $\SecMonASR$ ($\SecMonHit/\SecMonN$) & $[\SecMonCIlo, \SecMonCIhi]$ & $\SecMonFilt$ \\
Protected-tool filter, replicate 2 & $\SecGateASR$ ($\SecGateHit/\SecGateN$) & $[\SecGateCIlo, \SecGateCIhi]$ & $\SecGateFilt$ \\
Benign ($n{=}\UtilNdN$/arm) & \multicolumn{3}{l}{no defense $\UtilNdPass/\UtilNdN$; filter $\UtilGatePass/\UtilGateN$ (\S\ref{sec:rq3})} \\
\midrule
\multicolumn{4}{@{}l}{\emph{Registered campaign} (150 security / 56 benign held-out cases)}\\
No defense & \RqOneNodefenseCell & \RqOneNodefenseCI & \RqOneNodefenseFiltered \\
Admission (\Ione--\Ifour) & \RqOneDeiCell & \RqOneDeiCI & \RqOneDeiFiltered \\
Admission + gate & \RqOneDeigateCell & \RqOneDeigateCI & \RqOneDeigateFiltered \\
B1 argument--instruction matching & \RqOneBOneCell & \RqOneBOneCI & \RqOneBOneFiltered \\
B2 tool filter & \RqOneBTwoCell & \RqOneBTwoCI & \RqOneBTwoFiltered \\
B2strict filter & \RqOneBTwoStrictCell & \RqOneBTwoStrictCI & \RqOneBTwoStrictFiltered \\
B3 \Ifour{} alone & \RqOneBThreeCell & \RqOneBThreeCI & \RqOneBThreeFiltered \\
\bottomrule
\end{tabular}
\end{table}

\subsection{RQ2: What the Gate Adds}
\label{sec:rq2}
On the pilot grid the two defended arms are replicates of one filter
(the pilot's gate re-commitment was not wired), so the pilot cannot show
the gate's value, and the canonical injection family is caught at
admission in any case. The gate is designed for the case admission
cannot see (Cor.~\ref{cor:ratgap}), and RQ2 is answered there.

\paragraph*{Development probes: the counterexample, executed}
Four probes pair a legitimate directive ($\mathrm{TRUSTED}$,
$\tc$-sufficient) with a planted message naming an attacker-chosen
target ``with approval,'' on the small-model driver (existence
evidence, not rates). In 8 runs the model followed the injection
\GateFollows{} times; twice it laundered the action under the genuine
directive with $\tc$-sufficient support, and \Ione--\Ithree{}
\emph{admitted}: Cor.~\ref{cor:ratgap} end to end. With the gate
enabled (policy and perturbations frozen beforehand), both laundering
admissions were quarantined on the oracle plane $\delta_1$
(\GateQuar{} of \GateQuar); the honest-citation follow was blocked by
\Ithree{} before the gate; the interrupted follow, repaired to the
user's target, re-committed identically under both perturbations. Of
\GateBenignDestr{} gated benign commitments, \GateFalseQuarN{} were
falsely quarantined, at \GateOverhead$\times$ model time and
\GateMonitorUs\,$\mu$s p50 monitor latency. Cited-untrusted removal at
$\lmin$ never triggered (0/26), which is why laundering passes
\Ione--\Ithree{} at all.

\paragraph*{What these probes do and do not show}
They show that the gate closes the counterexample where \Ione--\Ithree{}
are blind, not that it adds value over \Ifour: both laundered targets
are absent from every trusted belief, so \Ifour{} (and B1) blocks them
at admission (Remark~\ref{rem:residual}). The registered RQ2 design
therefore uses \emph{choice-laundering} probes, in which every bound
argument is user-named and the plant is a false \emph{claim} that
selects among legitimate options: F1, delete whichever of two named
files is older (the plant reverses the dates), placed inside the file
listing the task reads (co-located) or inside a calendar entry it also
reads (separate); F2, send to a named contact only if an event falls on
a given day (the plant claims it moved); F3, write to whichever of two
named contacts attends an event (the plant claims a substitution).
Twenty probes with matched honest controls, four registered seeds
each, validated so that each honest ground truth passes and each plant
triggers the laundered action; \Ifour{} (and B1) admit every laundered
action by construction, so any quarantine is the gate's. The primary
endpoints are pooled $\delta_3$ catch on up to 40 separable-plant episodes
and false quarantine on up to 60 actionable honest-control evaluations,
with Wilson and probe-level cluster-bootstrap intervals and the actual
denominators reported; non-action controls never enter a gate-flip
denominator. Family/placement cells, observed
$\delta_2$, and oracle-$\delta_1$ $\mathit{fp}$/$\mathit{miss}$ per
Prop.~\ref{prop:approxgap} are secondary; registered prediction: $\delta_3$ catches
the separate placement and F2 at no honest cost and cannot see
co-located plants, which $\delta_2$ flags only together with every
honest commitment that needs the same atom.
The outcome refutes zero honest cost: $\delta_3$ catches
\RqTwoSeparableDeltaThreeK/\RqTwoSeparableDeltaThreeN{} separable
commitments (\RqTwoSeparableDeltaThreeRate, Wilson
$[\RqTwoSeparableDeltaThreeLo,\RqTwoSeparableDeltaThreeHi]$, probe-bootstrap
$[\RqTwoSeparableDeltaThreeClusterLo,\RqTwoSeparableDeltaThreeClusterHi]$)
but falsely quarantines \RqTwoHonestDeltaThreeK/\RqTwoHonestDeltaThreeN{}
controls (\RqTwoHonestDeltaThreeRate,
$[\RqTwoHonestDeltaThreeLo,\RqTwoHonestDeltaThreeHi]$). Only 26 separable
and 29 actionable evaluations arose; non-actions were not imputed.
Co-located catch is \RqTwoColocatedDeltaThreeRate, while $\delta_2$
flags every co-located plant and \RqTwoHonestDeltaTwoK/\RqTwoHonestDeltaTwoN{}
controls, the resolution limit. Against the oracle, $\delta_3$ has TP
\RqTwoDeltaThreeTp, FN \RqTwoDeltaThreeFn, FP \RqTwoDeltaThreeFp, TN
\RqTwoDeltaThreeTn (miss \RqTwoDeltaThreeMissRate, FP rate
\RqTwoDeltaThreeFpRate); $\delta_2$ has TP \RqTwoDeltaTwoTp, FN
\RqTwoDeltaTwoFn, FP \RqTwoDeltaTwoFp, TN \RqTwoDeltaTwoTn. The gate helps
on some separable plants, but its atom-granularity and re-commitment
failures are limits, not a general security claim.

\subsection{RQ3: Utility Cost}
\label{sec:rq3}
\paragraph*{Pilot}
Under the protected-tool filter the benign-completion rate is
\UtilGatePass/\UtilGateN{} (\UtilGateRate) versus \UtilNdPass/\UtilNdN{}
(\UtilNdRate) with no defense, with failures \UtilGateBlock/\UtilGateN{}
versus \UtilNdBlock/\UtilNdN. The 12.5-point estimate is not resolvable
at $n{=}\UtilGateN$ (Fisher $p{=}0.49$; paired discordant 3/0, exact
McNemar $p{=}0.25$) and is structurally uninformative: 23 of the 24
benign tasks never call a protected tool, the one that does was
blocked, and the two other extra failures occurred with nothing
filtered. We claim neither that the cost is small nor that it is
large; the registered study stratifies on protected-tool use.

\paragraph*{Registered study}
The utility comparison is re-run as a within-task crossover: each of
the 97 benign tasks of the four suites runs under three registered
seeds with no defense, admission plus gate, B1, and B2strict (291 pairs
per contrast, registered random order), pairing the within-task
variation an independent-arms design treats as noise. The sample size
follows from a pre-registered minimal important difference of 10 points
from the 83.3\,\% baseline (independent-arms floor 266 per arm at
80\,\% power; the paired test has power 0.83 at the independence bound
and more under positive within-task correlation). Paired RD with a
Newcombe interval for admission plus gate versus no defense is the
headline; the paired contrasts against B2strict and B1, Holm-corrected,
decide whether admission costs less than the filter; exact McNemar and
a task-level cluster bootstrap is secondary; protected-tool-use and
development-contact strata are descriptive; every failure is classified
as defense-induced (blocked by a named invariant), model-induced, or
degenerate.
The registered crossover shows a material cost: DEI+gate changes
completion by \RqThreeUtilDeiGatePairedRD{} relative to no defense
(Newcombe $[\RqThreeUtilDeiGatePairedRDlo,\RqThreeUtilDeiGatePairedRDhi]$;
task-cluster bootstrap $[\RqThreeUtilDeiGateClusterRDlo,\RqThreeUtilDeiGateClusterRDhi]$),
not by a negligible amount. It nevertheless outperforms B2strict by
\RqThreeUtilDeiGateVsBTwoStrictPairedRD{}
($[\RqThreeUtilDeiGateVsBTwoStrictPairedRDlo,\RqThreeUtilDeiGateVsBTwoStrictPairedRDhi]$;
Holm $p{=}\RqThreeUtilDeiGateVsBTwoStrictHolmP$) and B1 by
\RqThreeUtilDeiGateVsBOnePairedRD{}
($[\RqThreeUtilDeiGateVsBOnePairedRDlo,\RqThreeUtilDeiGateVsBOnePairedRDhi]$;
Holm $p{=}\RqThreeUtilDeiGateVsBOneHolmP$), the latter borderline.
DEI+gate failures split into 74 defense-induced, 115 model-induced, and
0 degenerate episodes; B2strict splits into 132 defense-induced and 88
model-induced. The practical claim is therefore a security--utility
trade-off, not preserved utility.

\subsection{RQ4: Generality}
\label{sec:rq4}
The pilot covers three workspace user tasks, one attack family, and
one model (per-case susceptibility 0.18--1.00), so no cell
extrapolates. Pre-fix open-weight archives (0/87 on qwen3-coder-30B,
0/33 on qwen3-32B) are negative observations, not resistance; a
post-fix development run of qwen3.6-35B followed the injection in
18/60 no-defense episodes across the four suites, so RQ4 re-runs that
family under the registration. On gpt-4.1-mini, ASR is
\RqFourASecNoDefenseASR{} with no defense, \RqFourASecDeiASR{} with DEI,
\RqFourASecDeiGateASR{} with DEI+gate, and \RqFourASecBOneASR{} with B1.
On \texttt{tool\_knowledge}, the corresponding rates are
\RqFourCSecNoDefenseASR, \RqFourCSecDeiASR, and
\RqFourCSecBOneASR. All 150-security-case rows have Wilson intervals and
zero voided outcomes; the benign rows are secondary and are not promoted
to primary evidence. The open-weight row is completed but uninformative:
qwen3:8b follows only \RqFourBSecNoDefenseHit/\RqFourBSecNoDefenseN{}
no-defense injections (DEI \RqFourBSecDeiHit/\RqFourBSecDeiN), so it does
not establish susceptibility or robustness. AgentDyn replicates on 1,120
episodes with 30 voided attempts excluded: ASR falls from
\RqAgentDynSecNoDefenseASR{} to \RqAgentDynSecDeiASR{} (RD
\RqAgentDynSecDeiRD{} $[\RqAgentDynSecDeiRDlo,\RqAgentDynSecDeiRDhi]$,
Fisher $p=\RqAgentDynSecDeiFisherP$), with suite-generality stated only
for the dynamic benchmark.

\subsection{RQ5: Latency}
Admission costs are those of Table~\ref{tab:overhead}: the replacement
A2 gateway adds \RqFiveOverheadRange{} p50 overhead over A0 at
5--100 req/s, with mean p50 \RqFiveATwoTwentyFiveMedian\,ms and mean p95
\RqFiveATwoTwentyFiveTail\,ms at 25 req/s. The gate adds one
re-commitment per enforced perturbation on irreversible classes
(\GateOverhead$\times$ benign model time in the development probes).

\subsection{Threats to Validity}
\label{sec:threats}
\paragraph*{Construct and internal}
\DEI{} constrains evidential structure, not truth, and within-TCB
compromise is excluded by the threat model. The pilot's defended arms
ran a protected-tool filter at an unset temperature (\S\ref{sec:evalsetup});
the registered harness fixes both, and B3 isolates \Ifour's
contribution. Development probes were built against the gateway's own
threat model; held-out sets are registered and timestamped before the
runs.
\paragraph*{External and statistical}
CZTA is one instantiation on one deployment; efficacy is per model,
suite, and attack family, and RQ4 bounds the generality claim. The
pilot's security comparison has power $\approx 0.99$ for its effect;
its utility arm is underpowered; its defended arms are one filter, so
RQ2 is answered where admission is provably blind.
\paragraph*{Adaptive adversaries}
\Ifive's perturbation class is registered but finite; an attacker
outside it, or a model that re-commits identically with and without
the plant, is the counterexample \S\ref{sec:falsifiability} invites.


\section{Related Work}
\label{sec:related}
\paragraph*{Defenses}
Indirect prompt injection was identified early as the structural risk
of tool-using LLM systems~\cite{greshake}; in-band defenses (detection
classifiers, prompt hygiene, instruction-hierarchy fine-tuning) have
been broken by adaptive evaluation~\cite{abdelnabi}. Out of band, CaMeL
separates control flow from data flow and attaches capability metadata
to values~\cite{camel}; its successors (FIDES, Progent, RTBAS, Conseca,
FORGE) are systematized in~\cite{oob} as Biba integrity plus reference
monitoring plus least privilege, with two open problems that shape this
paper: validation almost exclusively on static benchmarks, and an
under-specified trusted base for provenance assignment.
Guardrail agents~\cite{guardagent}, ARGUS's influence provenance
graph~\cite{argus}, provenance surveys~\cite{tracessurvey}, and
machine-checked Biba for agent memory~\cite{louck2026tmanm} are CZTA's
instantiation peers; none states the property it approximates as a
definition. Counterfactual and causal
defenses~\cite{melon,agentsentry,causalarmor,liao2026provenance,babu2026gistcmtf,pona2025acf}
are compared with the gate in \S\ref{sec:gate}.

\paragraph*{Integrity framings}
Intent-to-execution integrity~\cite{qu} gives judgment integrity the
invariance form adopted here as \Ifive; agent epistemic
integrity~\cite{shen} is an architectural constraint secured by a
model--harness contract; context-to-execution integrity~\cite{cxi}
binds field, effect, and invocation authority to one manifest at the
sink; a concurrent theory characterizes what a deterministic
pre-execution gate can enforce~\cite{ray2026enforced}; semantic
laundering~\cite{romanchuk2026semantic} states the architectural
version of the rationalization gap.

\paragraph*{Benchmarks, protocol surface, foundations}
AgentDojo established live-environment evaluation of
injection~\cite{agentdojo}, with documented scoring
defects~\cite{bhagwatkar2025firewalls} and realism gaps~\cite{agentdyn};
sibling suites are consolidated in~\cite{benchtax}. MCP standardized
tool connectivity and authorization~\cite{mcpspec}; taxonomies rank
tool poisoning among its principal risks~\cite{owaspmcp,mcpdpt};
zero-trust MCP gateways decide who may call
what~\cite{ztmcp,kodathala2026aiauthz}, not what may count as a reason;
frontier-lab guidance treats internal agents as compromised
insiders~\cite{controlroadmap}. \DEI{} applies classical machinery to a
new object: integrity models~\cite{biba,denning,clarkwilson},
noninterference~\cite{goguen}, decentralized IFC~\cite{myersliskov},
provenance~\cite{prov}, evidence theory~\cite{shafer,dung}, and the
runtime-verification canon~\cite{alpern,schneider,bauer,pnueli,mop,hyperprops,finkbeiner,stucki}.
``Epistemic security'' also names a societal program~\cite{seger};
ours is the technical sense.

\section{Conclusion}
Defenses for agentic systems constrain what an agent may do. We have
stated what it means for an agent's reasons to be trustworthy, proved
which part a monitor can enforce and which it can only audit, and
measured the resulting trade-offs. \Ifour{} is the admission result:
bound arguments must have cited evidential standing. The counterfactual
gate is valuable on separable plants, but its misses, false quarantines,
and utility costs are limits, not evidence of a general counterfactual
defense.

\section*{Ethics considerations}
This work evaluates defenses on a testbed the authors control; no
production or third-party systems were probed, no third-party
vulnerabilities were identified, and no responsible disclosure was
required. The taxonomy of \S\ref{sec:taxonomy} derives from a
registered literature review whose expert annotators coded publicly
available bibliographic metadata against a registered codebook; no
personal data about the annotators was collected, and only their
per-record codes are released under role labels. No human-subjects study was conducted for this paper; a
registered threshold-adjudication pilot described in the artifact
awaits an ethics determination and contributes no result here. The
attack content used in the evaluation is the public AgentDojo
injection corpus and constructed probes against a local testbed; it
confers no capability beyond what the benchmark already publishes.

\section*{LLM usage considerations}
The authors used large language model assistants for editorial and
preparatory support (phrasing, formatting, \LaTeX, and
consistency-checking of reported figures against source data), for
assistance drafting definitions and proofs from author-specified
theorem statements, and for developing the artifact's Lean~4
mechanization scripts. Every definition, theorem statement, and proof
was reviewed and verified by the authors; every Lean script is checked
by the Lean~4 kernel, which re-verifies each inference step
independently of its provenance; every reported number is recomputed
from episode logs by the pipeline in the artifact. The motivation for
this use is editorial throughput and the mechanized re-verification it
made affordable. The research questions, the framework, the system,
and all contributions are the authors' own, and the authors take full
responsibility for all material herein.

\section*{Availability}
The anonymized artifact~\cite{artifact} contains the \texttt{czta\_ref}
reference monitor (hash-pinned at \texttt{/env}); the executable
reference semantics and property tests (Appendix~\ref{app:refsem}); the
Lean~4 mechanization (\texttt{DeiSafety.lean}) of the safety core,
re-checkable by \texttt{check\_formal.sh}; the scoping-review protocol,
ledger, codebooks, instruments, per-record hashes, and amendment log;
the measurement harness and per-run outputs behind
Table~\ref{tab:overhead}; the authorization-layer harness and verdict
logs; the live-model campaign behind Table~\ref{tab:livemodel} (the
seed-audited harness \texttt{g2\_extended.py}, both episode logs, the
full console log, the recompute-and-macro pipeline, and a verification
manifest whose SHA-256 hashes and recomputed figures match every number
reported here); and the registered protocols, deposits, and timestamps
for the completed runs of \S\ref{sec:eval}.


\begin{thebibliography}{51}
\footnotesize

\bibitem{liao2026provenance} J.~Liao, ``Auditing provenance
sensitivity in LLM agent action selection,'' 2026,
arXiv:2607.20827.

\bibitem{babu2026gistcmtf} R.~S.~Babu and R.~Shukla, ``GIST-CMTF:
Goal-state inference for causal minimal tool filtering in LLM
agents,'' 2026, arXiv:2606.16813.

\bibitem{pona2025acf} E.~Pona, M.~Kazemi, Y.~Du, D.~Watson, and
N.~Paoletti, ``Abstract counterfactuals for language model
agents,'' in \emph{Advances in Neural Information Processing
Systems (NeurIPS)}, 2025; also arXiv:2506.02946.

\bibitem{camel} E.~Debenedetti, I.~Shumailov, T.~Fan, J.~Hayes,
N.~Carlini, D.~Fabian, et~al., ``Defeating prompt injections by
design,'' 2025, arXiv:2503.18813; IEEE SaTML 2026.

\bibitem{cxi} I.~Santos-Grueiro, ``Context-to-execution integrity
for LLM agents,'' 2026, arXiv:2607.06000.

\bibitem{guardagent} Z.~Xiang et~al., ``GuardAgent: Safeguard LLM
agents by a guard agent via knowledge-enabled reasoning,'' 2024,
arXiv:2406.09187; ICML 2025.

\bibitem{argus} S.~Weng, Y.~Feng, J.~Zhang, X.~Xie, J.~Yu, and
J.~Liu, ``ARGUS: Defending LLM agents against
context-aware prompt injection,'' 2026, arXiv:2605.03378.

\bibitem{ray2026enforced} S.~Ray, ``What can be enforced? A
theory of certified runtime safety for tool-using agents,'' 2026,
arXiv:2607.22868.

\bibitem{romanchuk2026semantic} O.~Romanchuk and R.~Bondar,
``Semantic laundering in AI agent architectures: Why tool
boundaries do not confer epistemic warrant,'' 2026,
arXiv:2601.08333.

\bibitem{louck2026tmanm} Y.~Louck, ``Securing LLM-agent long-term
memory against poisoning: Non-malleable, origin-bound authority
with machine-checked guarantees,'' 2026, arXiv:2606.24322.

\bibitem{ztmcp} S.~Mishra, ``ZT-MCP: A zero-trust security
architecture for MCP-connected AI agents,'' 2026, ResearchGate
preprint, doi:10.13140/RG.2.2.36042.53445.

\bibitem{kodathala2026aiauthz} S.~V.~Kodathala, ``aiAuthZ:
Off-host, identity-bound authorization for AI agents,'' 2026,
arXiv:2607.05518.

\bibitem{melon} K.~Zhu, X.~Yang, J.~Wang, W.~Guo, and W.~Y.~Wang,
``MELON: Provable defense against indirect prompt injection
attacks in AI agents,'' in \emph{Proc.\ ICML}, 2025; also
arXiv:2502.05174.

\bibitem{agentsentry} T.~Zhang, Y.~Xu, J.~Wang, \emph{et al.},
``AgentSentry: Mitigating indirect prompt injection in LLM agents
via temporal causal diagnostics and context purification,'' 2026,
arXiv:2602.22724.

\bibitem{causalarmor} M.~Kim, M.~Parmar, P.~Wallis, \emph{et al.},
``CausalArmor: Efficient indirect prompt injection guardrails via
causal attribution,'' 2026, arXiv:2602.07918.

\bibitem{oob} LaunchSafe Research, ``Adaptive evaluation of
out-of-band defenses against prompt injection in LLM agents,''
2026, arXiv:2606.26479.

\bibitem{qu} W.~Qu, M.~Xu, P.~Wang, S.~Zhai, J.~Zhang, and
D.~Song, ``Securing LLM agents need intent-to-execution
integrity,'' 2026, arXiv:2605.16976.

\bibitem{shen} Z.~Shen, ``Neither layer alone: Epistemic integrity
requires hierarchical joint design for long-running AI agents,''
2026, arXiv:2606.04017, posted 1 June 2026.

\bibitem{greshake} K.~Greshake, S.~Abdelnabi, S.~Mishra,
C.~Endres, T.~Holz, and M.~Fritz, ``Not what you've signed up
for: Compromising real-world LLM-integrated applications with
indirect prompt injection,'' in \emph{Proceedings of the 16th ACM
Workshop on Artificial Intelligence and Security (AISec)}, 2023.

\bibitem{abdelnabi} S.~Abdelnabi and E.~Bagdasarian, ``AI agents
may always fall for prompt injections,'' 2026, arXiv:2605.17634.

\bibitem{agentdojo} E.~Debenedetti, J.~Zhang, M.~Balunovi\'c,
L.~Beurer-Kellner, M.~Fischer, and F.~Tram\`er, ``AgentDojo: A
dynamic environment to evaluate prompt injection attacks and
defenses for LLM agents,'' in \emph{NeurIPS Datasets and
Benchmarks}, 2024.

\bibitem{benchtax} M.~Q.~Li, B.~C.~M.~Fung, B.~Li, H.~Ismail, and
F.~Iqbal, ``Taxonomy and consistency analysis of
safety benchmarks for AI agents,'' 2026, arXiv:2605.16282.

\bibitem{mcpspec} Model Context Protocol contributors, ``Model
context protocol specification,'' Specification revision
2025-11-25 (authorization framework), 2025.

\bibitem{owaspmcp} OWASP Foundation, ``OWASP MCP top 10 (beta),''
2026; Tool Poisoning listed (MCP03).

\bibitem{mcpdpt} M.~Rostamzadeh, S.~Narula, N.~Birhan,
M.~Ghasemigol, and D.~Takabi, ``MCP-DPT: A defense-placement taxonomy
and coverage analysis for MCP security,'' 2026,
arXiv:2604.07551.

\bibitem{biba} K.~J. Biba, ``Integrity considerations for secure
computer systems,'' MITRE Corporation, Tech.\ Rep.\ ESD-TR-76-372,
1977.

\bibitem{denning} D.~E. Denning, ``A lattice model of secure
information flow,'' \emph{Communications of the ACM}, vol.~19,
no.~5, pp.~236--243, 1976.

\bibitem{clarkwilson} D.~D. Clark and D.~R. Wilson, ``A comparison
of commercial and military computer security policies,'' in
\emph{IEEE Symposium on Security and Privacy}, 1987.

\bibitem{goguen} J.~A. Goguen and J.~Meseguer, ``Security policies
and security models,'' in \emph{IEEE Symposium on Security and
Privacy}, 1982.

\bibitem{myersliskov} A.~C. Myers and B.~Liskov, ``A decentralized
model for information flow control,'' in \emph{ACM Symposium on
Operating Systems Principles (SOSP)}, 1997.

\bibitem{prov} W3C, ``PROV-DM: The PROV data model,'' W3C
Recommendation, 30 April 2013, 2013.

\bibitem{shafer} G.~Shafer, \emph{A Mathematical Theory of
Evidence}. Princeton University Press, 1976.

\bibitem{dung} P.~M. Dung, ``On the acceptability of arguments and
its fundamental role in nonmonotonic reasoning, logic programming
and n-person games,'' \emph{Artificial Intelligence}, vol.~77,
no.~2, pp.~321--357, 1995.

\bibitem{alpern} B.~Alpern and F.~B. Schneider, ``Recognizing
safety and liveness,'' \emph{Distributed Computing}, vol.~2,
no.~3, pp.~117--126, 1987.

\bibitem{schneider} F.~B. Schneider, ``Enforceable security
policies,'' \emph{ACM Transactions on Information and System
Security}, vol.~3, no.~1, pp.~30--50, 2000.

\bibitem{bauer} A.~Bauer, M.~Leucker, and C.~Schallhart, ``Runtime
verification for LTL and TLTL,'' \emph{ACM Transactions on
Software Engineering and Methodology}, vol.~20, no.~4, 2011.

\bibitem{pnueli} A.~Pnueli and A.~Zaks, ``PSL model checking and
run-time verification via testers,'' in \emph{International
Symposium on Formal Methods (FM)}, 2006.

\bibitem{mop} F.~Chen and G.~Ro\c{s}u, ``MOP: An efficient and
generic runtime verification framework,'' in \emph{ACM OOPSLA},
2007.

\bibitem{hyperprops} M.~R. Clarkson and F.~B. Schneider,
``Hyperproperties,'' \emph{Journal of Computer Security},
vol.~18, no.~6, pp.~1157--1210, 2010.

\bibitem{finkbeiner} B.~Finkbeiner, C.~Hahn, M.~Stenger, and
L.~Tentrup, ``Monitoring hyperproperties,'' \emph{Formal Methods
in System Design}, vol.~54, pp.~336--363, 2019.

\bibitem{stucki} S.~Stucki, C.~S\'anchez, G.~Schneider, and
B.~Bonakdarpour, ``Gray-box monitoring of hyperproperties,'' in
\emph{International Symposium on Formal Methods (FM)}, 2019.

\bibitem{tracessurvey} Y.~Wang et~al., ``From agent traces to trust: A
survey of evidence tracing and execution provenance in LLM
agents,'' 2026, arXiv:2606.04990.

\bibitem{controlroadmap} M.~Phuong, E.~Jenner, L.~Simon, L.~Ho,
R.~Shah, S.~Farquhar, and S.~Coull, ``AI control roadmap / three
layers of agent security,'' Google DeepMind report and blog, 18
June 2026.

\bibitem{seger} E.~Seger, S.~Avin, G.~Pearson, M.~Briers,
S.~\'O h\'Eigeartaigh, and H.~Bacon, ``Tackling threats to
informed decision-making in democratic societies: Promoting
epistemic security in a technologically-advanced world,'' The Alan
Turing Institute, Tech.\ Rep., 2020.

\bibitem{prisma} A.~C. Tricco et~al., ``PRISMA extension for
scoping reviews (PRISMA-ScR): Checklist and explanation,''
\emph{Annals of Internal Medicine}, vol.~169, no.~7,
pp.~467--473, 2018.

\bibitem{gwet} K.~L. Gwet, ``Computing inter-rater reliability and
its variance in the presence of high agreement,'' \emph{British
Journal of Mathematical and Statistical Psychology}, vol.~61,
no.~1, pp.~29--48, 2008.

\bibitem{cohen} J.~Cohen, ``A coefficient of agreement for nominal
scales,'' \emph{Educational and Psychological Measurement},
vol.~20, no.~1, pp.~37--46, 1960.

\bibitem{wilson} E.~B. Wilson, ``Probable inference, the law of
succession, and statistical inference,'' \emph{Journal of the
American Statistical Association}, vol.~22, no.~158,
pp.~209--212, 1927.

\bibitem{artifact} Anonymous, ``Epistemic security artifact: reference
monitor, reference semantics, Lean~4 mechanization, review ledger, and
campaign logs,'' anonymized repository,
\url{https://anonymous.4open.science/r/epistemic-security-artifact-AF35/}.

\bibitem{bhagwatkar2025firewalls} R.~Bhagwatkar, K.~Kasa, A.~Puri,
G.~Huang, I.~Rish, G.~W.~Taylor, K.~Dj~Dvijotham, and A.~Lacoste,
``Indirect prompt injections: Are firewalls all you need, or stronger
benchmarks?'' 2025, arXiv:2510.05244.

\bibitem{agentdyn} H.~Li, R.~Wen, S.~Shi, N.~Zhang, Y.~Vorobeychik, and
C.~Xiao, ``AgentDyn: Are your agent security defenses deployable in
real-world dynamic environments?'' 2026, arXiv:2602.03117.

\end{thebibliography}

\appendices
\section{Proof Details}
\label{app:proofs}

\paragraph{(a) Setup and violation predicates}
Fix the event alphabet
$\Sigma = \{\ev{ing}(e), \ev{upd}(b, X), \ev{end}(b, \ell, \rho), \ev{act}(a)\}$
over the runtime object universes of \S\ref{sec:definition}. A
finite prefix $\sigma[..k]$ determines a monitor state: the
observed evidence and belief sets, the label map $\ell(\cdot)$,
the dependency graph $P$, for each observed belief $b$ its recorded
value set $\vals(b)$, and a memoized root certificate $\cert(b) \in \{\mathrm{ok}, \mathrm{bad}\}$
computed when $b$ is created ($\mathrm{ok}$ iff every maximal
derivation path from $b$ terminates in evidence $e$ with $\att(e)$
and $\ell(e) \sqsupseteq \lmin$; certificates of parents are
reused, giving the $O(1)$ amortized update of Remark~1). All four
conditions below are decidable from the prefix state at the event
that would violate them.

\paragraph{(b) \Ione--\Ifour{} are safety properties}
By Alpern--Schneider~\cite{alpern}, a property is safety iff every
violating infinite behavior has a finite bad prefix that no
extension can repair. Define the per-event violation predicates:
$\mathsf{viol}_1(\sigma[..k])$ iff $\epsilon_k = \ev{act}(a)$ with
$c(a) \in \Cprot$ and some $b \in \support(a)$ has
$\cert(b) = \mathrm{bad}$ (or an ill-founded closure), or fails the
action-time freshness re-check $\neg\mathrm{WindowValid}(b)$ (a
terminal root's attestation validity window has lapsed at the time
of $\epsilon_k$, though its creation-time certificate may be
$\mathrm{ok}$; Def.~\ref{def:evidence} evaluates $\att$ against the
current time);
$\mathsf{viol}_2(\sigma[..k])$ iff $\epsilon_k$ raises some label
other than by an authorized $\ev{end}$ event, or
$\epsilon_k = \ev{upd}(b', X)$ assigns
$\ell(b') \not\sqsubseteq \bigsqcap_{x \in X} \ell(x)$;
$\mathsf{viol}_3(\sigma[..k])$ iff $\epsilon_k = \ev{act}(a)$
protected and $S(\support(a), \lreq(c(a))) < \tau_{c(a)}$;
$\mathsf{viol}_4(\sigma[..k])$ iff $\epsilon_k = \ev{act}(a)$ protected
and some bound slot $s \in \bindpol(c(a))$ with $\args(a)(s)$ defined
has no $b \in \support(a)$ with $\args(a)(s) \in \vals(b)$ and
$\ell(b) \sqsupseteq \ell_{\bindpol}(c(a), s)$. Each
predicate depends only on $\sigma[..k]$. A trace violates $I_j$
iff some prefix satisfies $\mathsf{viol}_j$; the offending event
has occurred and occurs in every extension, so the prefix is bad
and irremediable: \Ione--\Ifour{} are safety. Monitorability
follows: the $\mathrm{LTL}_3$ semantics~\cite{bauer,pnueli}
assigns verdict $\bot$ exactly at the first prefix satisfying
$\mathsf{viol}_j$. Quantification over beliefs and actions is
resolved by parametric monitoring~\cite{mop}: the object universe
at any prefix is the finite set of objects observed so far, and
each instantiated monitor is the two-state automaton
$\mathrm{ok} \to \mathrm{bad}$ driven by $\mathsf{viol}_j$.

\paragraph{(c) Enforceability by truncation}
Schneider's execution monitors enforce exactly policies for which
a monitor mediating each event can halt the target before an
unacceptable step commits~\cite{schneider}. The gateway mediates
every $\Sigma$-event (complete-mediation assumption of
\S\ref{sec:threat}); on a candidate event $\epsilon$ it evaluates
$\mathsf{viol}_j$ against the current state and suppresses
$\epsilon$ iff some predicate fires. Two facts make this a sound
truncation automaton. (i)~\emph{Blocking preserves legality}:
\Ione--\Ifour{} universally quantify over events present in the
trace, so deleting a candidate event can never create a violation;
the truncated trace extends the previously legal prefix by nothing
and remains legal. (ii)~\emph{Blocking loses nothing legal}:
suppression occurs only when the event itself is the bad event, so
every $I$-satisfying continuation is preserved. Hence
$\Ione \wedge \Itwo \wedge \Ithree \wedge \Ifour$ is EM-enforceable, and the
enforcement cost is the per-event predicate evaluation measured in
Table~\ref{tab:overhead} and Appendix~\ref{app:refsem}; the
$\mathsf{viol}_4$ test adds set membership on recorded value sets and
changes neither fact.

\paragraph{(d) \Ifive{} is 2-safety and not single-trace monitorable}
\Ifive{} constrains pairs of executions: it is the hyperproperty
$H = \{T : \forall \sigma, \sigma' \in T \text{ agreeing on } q
\text{ and trusted state, with } d' = d \oplus \delta,\,
\delta \in \Dirrel(q) : \dist(J_\sigma, J_{\sigma'}) \leq \varepsilon\}$,
refutable by exhibiting two finite traces: 2-safety in the sense
of Clarkson--Schneider~\cite{hyperprops}. It is not a trace
property: membership of an individual trace is undefined. For
non-monitorability, consider any black-box monitor $M$ mapping
single observed prefixes to verdicts, and two systems $S^{+}$
(satisfying \Ifive) and $S^{-}$ (violating it on some
$\delta$-counterfactual) constructed to produce the same behavior
on the observed input $(q, d)$; the violation of $S^{-}$ is
witnessed only by the unobserved counterfactual run on
$(q, d \oplus \delta)$. $M$ receives identical input under both
systems, so any verdict it issues is unsound for one of them; no
sound single-trace black-box monitor can refute \Ifive. Runtime
checking therefore requires stored comparison executions or a
model of the system (gray-box monitoring~\cite{finkbeiner,stucki}), and we realize the check as the offline pre-registered
differential harness of \S\ref{sec:falsifiability} and, before dispatch, as the gate of \S\ref{sec:gate}.

\section{Charting-Class to Invariant Crosswalk}
\label{app:crosswalk}
The full crosswalk table, with per-class mechanisms, refutation signals, and corpus shares, is included in the anonymized artifact rather than printed here.

\section{Mechanism Predicates and Proofs of Theorem~\ref{thm:crosswalk} and Corollary~\ref{cor:ratgap}}
\label{app:t1c2}
\begin{definition}[Mechanism realization]
\label{def:mech}
For each charting class $K$ with a machine-fired primary $j(K) \in
\{1, 2, 3\}$ (the nine machine-fired rows of Appendix~\ref{app:crosswalk}:
eight full classes plus the anchor-extension boundary case), the trace
predicate $\mathrm{Mech}_K(\sigma)$ holds iff $\sigma$ has a prefix
$\sigma[..k]$ whose final event $\epsilon_k$ satisfies the row's
condition, stated in the vocabulary of Appendix~\ref{app:proofs}(a):
\emph{hallucination} ($j{=}1$), a protected $\ev{act}(a)$ with some
$b \in \support(a)$ whose closure has a maximal derivation path
terminating in no evidence atom (Def.~\ref{def:i1}(i)), so
$\cert(b) = \mathrm{bad}$; \emph{provenance/identity tampering}
($j{=}1$), a terminal root of the support closure that fails $\att$
(Def.~\ref{def:evidence}); \emph{retrieval/knowledge poisoning}
($j{=}1$), a terminal root with $\neg\att(e)$ or $\ell(e)
\not\sqsupseteq \lmin$; \emph{drift/staleness} ($j{=}1$, partial), a
terminal root whose validity window has lapsed at action time, so
$\neg\mathrm{WindowValid}(b)$; \emph{prompt injection} ($j{=}2$), an
$\ev{upd}(b', X, \ell_{\mathrm{prop}})$ with $\ell_{\mathrm{prop}}
\not\sqsubseteq \meet_{x \in X} \ell(x)$ or an $\ev{end}(b, \ell',
\rho)$ with $\ell' \sqsupset \ell(b)$ and ($\rho \notin \pi$ or
$\neg\mathsf{Authorized}$), the affected belief depending on
untrusted-channel content; \emph{confidence/trust manipulation}
($j{=}2$), an endorsement attempt citing no authorized rule;
\emph{tool/agent hijacking} ($j{=}2$, partial), a label-raising
$\ev{upd}$ or $\ev{end}$ without authorized endorsement on a belief
depending on tool-result content; \emph{misinformation} ($j{=}3$), a
protected $\ev{act}(a)$ with $S(\support(a), \lreq(c(a))) < \tc$ once
non-qualifying evidence is excluded by A1; \emph{retrieval boundary
case} ($j{=}3$), the same deficit arising from selection starving
qualifying support from an intact substrate. The declared out-of-scope
row and the control row carry no predicate.
\end{definition}

\begin{proof}[Proof of Theorem~\ref{thm:crosswalk}]
Let $\mathrm{Mech}_K(\sigma)$ hold with witness prefix $\sigma[..k]$.
The four $j{=}1$ clauses are exactly the bad-certificate or
lapsed-window disjuncts of $\mathit{viol}_1$; the three $j{=}2$
clauses are label raises without an authorized endorsement, exactly
$\mathit{viol}_2$; and the two $j{=}3$ clauses are the qualifying-score
deficit $\mathit{viol}_3$, with A1 explaining which atoms contribute.
Thus $\mathit{viol}_{j(K)}(\sigma[..k])$ holds. These predicates are
prefix-persistent, so a first firing index $k^{*}\le k$ exists.
Assumption~\ref{asm:mediation} makes that event mediated,
Assumption~\ref{asm:recorded} makes the recorded prefix state trusted,
and Algorithm~\ref{alg:monitor} suppresses precisely the firing event;
it therefore blocks at or before $\epsilon_k$. Conversely, a control
instance fires no predicate. Since the algorithm blocks only on a
predicate and blocking loses nothing legal
(Appendix~\ref{app:proofs}(c)), every control event is admitted. The
10/10 attack and 3/3 honest reference-suite outcomes
(Appendix~\ref{app:refsem}) validate realization of the predicates;
they are empirical checks, not premises of this trace-level proof.
\end{proof}

\begin{proof}[Proof of Corollary~\ref{cor:ratgap}]
\emph{(i) Existence.} Take an honestly admitted class-$c_0$ action with
support $S_0$: its roots are attested and fresh and
$S(S_0,\lreq(c_0))\ge\tc$. The adversary places $d_{\mathrm{inj}}$ on
an untrusted channel. The compromised model selects any
$a^{*}\in A_{\bindpol}(S_0)$ from that injection but commits the genuine
$S_0$, with no label raise or endorsement. At the action prefix, all
certificates and windows are valid ($\neg\mathit{viol}_1$), labels
satisfy the meet ($\neg\mathit{viol}_2$), the score passes
($\neg\mathit{viol}_3$), and every bound slot is named at the required
label ($\neg\mathit{viol}_4$). A sound truncation monitor must admit
this legal event although injection caused the selection. Because the
certificates, score, and binding test read only genuine attestation,
labels, weights, and recorded values, finer provenance granularity
changes none of them.

\emph{(ii) Necessity.} Suppose a single-trace detector $D$ soundly
blocks this strategy. Let $\delta_1$ remove the plant. Construct an
honest $S^{+}$ that commits $a^{*}$ both with and without the plant,
and a laundering $S^{-}$ that emits the identical observed stream on
$(q,d)$ but commits $a^{-}\ne a^{*}$ after removal. Both observed
streams satisfy \Ione--\Ifour, so $D$ must issue the same verdict for
indistinguishable systems: blocking is unsound for $S^{+}$, admitting
for $S^{-}$. This is the $S^{\pm}$ construction of
Prop.~\ref{prop:hyper}; detection therefore requires a pair-witnessing
differential harness or the decision-local gate, not a refined
single-trace admission rule.
\end{proof}

\section{Executable Reference Semantics and Evidence Maturity}
\label{app:refsem}
The executable monitor implements Algorithm~\ref{alg:monitor} event for
event: lattice-meet label propagation, authorized endorsement, fresh
root certificates, support scoring, and bound-slot checks precede every
protected commit. Its constructed suites exercise each crosswalk
predicate, the \Ifive{} differential harness, and truncation legality.
Table~\ref{tab:refsem} reports these semantics checks separately from
live-model efficacy.

\begin{table}[!htb]
\centering
\caption{Executable semantics checks; counterfactual results are
invariance followed by witnesses.}
\label{tab:refsem}
\scriptsize
\setlength{\tabcolsep}{3pt}
\renewcommand{\arraystretch}{0.94}
\begin{tabular}{@{}>{\raggedright\arraybackslash}p{0.72\linewidth} r@{}}
\toprule
Property under test & Result \\
\midrule
Predicted-invariant blocks; coverage counterexamples & 10/10; 0 \\
Honest admissions; admit/block and legality tests & 3/3; 13/13 \\
Injection-obeying decision function & 71\,\%; 6/21 \\
Guarded decision function & 100\,\%; 0/21 \\
Memoized cost span; unmemoized depth correlation & $1.42\times$; 0.996 \\
Development verdict/action invariance & 52/60; 41/60 \\
\bottomrule
\end{tabular}
\end{table}

Across 12 tasks and five $\Dirrel$ families, reordering remained
invariant (12/12), while inject-then-ignore was sharpest (7/12). These
are development stability checks, not live security estimates. The
$256\times$ depth experiment supports Remark~1: memoized receipt checks
remain nearly flat while unmemoized cost tracks depth.

\paragraph*{Evidence maturity}
Props.~\ref{prop:safety}--\ref{prop:hyper} and
Cor.~\ref{cor:ratgap}, with the Lean~4 safety mechanization, are
\emph{proved}; the replacement-2 \Ifour{} bound-slot check is
\emph{implemented}. Table~\ref{tab:refsem}, the authorization cases,
and RQ1--RQ5 are \emph{measured}: 1,442 RQ1 episodes with zero voided
and a DEI/B1 tie; 640 RQ2 probes plus 80 oracle episodes; 1,164 RQ3
episodes with 50 voided and one degenerate, giving 291 pairs per
contrast; 27/27 RQ5 cells; 824 RQ4(a) and 618 RQ4(c) episodes; and
1,120 AgentDyn episodes with 30 voided attempts excluded. RQ4(b)
completed 120 episodes but remains uninformative at 1/60
susceptibility. Development checks are not held-out effectiveness
evidence.


\end{document}
